% Lesson 50 — Grammar: General revision — Anglais Tle A — Cours signé OnBuch+
% Programme officiel MINESEC. Compiler avec : tectonic EN50-grammar-general-revision.tex
\newcommand{\DOCMATIERE}{Anglais}
\newcommand{\DOCNIVEAU}{Tle A}
\newcommand{\DOCMODULE}{Chapter 5 : Media and Communication}
\newcommand{\DOCLECON}{EN50}
\newcommand{\DOCTITRE}{Lesson 50 — Grammar: General revision}
% OnBuch+ — préambule « cours signés » ENRICHI (compilé avec tectonic / XeLaTeX)
% Système visuel « Pop » d'OnBuch+ : fond crème (jamais blanc), bordures encre
% épaisses, ombres « dures » décalées, titres Archivo Black en capitales.
% Version MATHÉMATIQUES : graphes (pgfplots/tikz), boîtes pédagogiques
% (définition, propriété, méthode, attention, le savais-tu, exercice résolu,
% exercice, TP, bilan), page de garde et signature OnBuch+ ancrée en bas.
%
% Macros attendues dans main.tex AVANT \input{preamble} :
%   \DOCMATIERE  \DOCNIVEAU  \DOCMODULE  \DOCLECON (numéro)  \DOCTITRE
\documentclass[11pt]{article}

\usepackage{fontspec}
\usepackage[english,french]{babel} % français = langue principale ; l'anglais via \eng{..} / textebox
\usepackage[a4paper,margin=2cm,top=3.4cm,bottom=2.8cm,headheight=1.4cm,footskip=1.3cm]{geometry}
\usepackage{amsmath,amssymb,amsfonts}
\usepackage{mathtools} % \xmapsto, \coloneqq... souvent utilisés par les modèles
\usepackage{cancel} % simplifications barrées (bilans de forces)
% \abs{z} (module, valeur absolue) et \norm{u} (norme), utilisés par les modèles
\providecommand{\abs}{}\let\abs\relax\DeclarePairedDelimiter{\abs}{\lvert}{\rvert}
\providecommand{\norm}{}\let\norm\relax\DeclarePairedDelimiter{\norm}{\lVert}{\rVert}
\usepackage[table]{xcolor} % [table] : fournit aussi \rowcolors (alternance de lignes)
\usepackage{pagecolor}
\usepackage{tikz}
\usetikzlibrary{arrows.meta,positioning,calc,shapes.geometric,shapes.misc,shapes.arrows,decorations.pathmorphing,decorations.pathreplacing,decorations.markings,patterns,shadows,mindmap,trees,fit,backgrounds,matrix,intersections,angles,quotes}
\usepackage{pgfplots}
\usepackage[version=4]{mhchem} % équations nucléaires \ce{^{235}_{92}U}
\usepackage[european,siunitx]{circuitikz} % schémas électriques
\pgfplotsset{compat=1.18}
\usepgfplotslibrary{fillbetween,groupplots} % aires entre courbes (calcul intégral)
\usepackage{enumitem}
\usepackage{fancyhdr}
% [most] charge toutes les bibliothèques tcolorbox (listings, minted, raster,
% documentation...) et gonfle inutilement la mémoire TeX sur un document long
% (~300 boîtes) : on ne charge que ce qui est réellement utilisé.
\usepackage[skins,breakable]{tcolorbox}
\usepackage{graphicx}
\usepackage{colortbl}
\usepackage{booktabs}
\usepackage{tabularx}
\usepackage{diagbox} % case d'en-tête diagonale des tableaux à double entrée
% Colonnes extensibles centrée / gauche / droite (C, L, R), souvent utilisées par les modèles
\newcolumntype{C}{>{\centering\arraybackslash}X}
\newcolumntype{L}{>{\raggedright\arraybackslash}X}
\newcolumntype{R}{>{\raggedleft\arraybackslash}X}
\usepackage{array}
\usepackage{multicol}
\usepackage{multirow}
\usepackage{siunitx}
% parse-numbers=false : accepte aussi \SI{2,5 \times 10^{-3}}{...}
\sisetup{locale=FR,output-decimal-marker={,},group-minimum-digits=4,parse-numbers=false}
\DeclareSIUnit\ohm{\ensuremath{\symup{\Omega}}} % U+2126 absent de la police
\DeclareSIUnit\division{div} % réglages d'oscilloscope (ms/div, V/div)
\DeclareSIUnit\tr{tr} % tours (vitesse de rotation, ex. \SI{1500}{\tr\per\minute})
\DeclareSIUnit\molar{M} % concentration molaire (ex. \SI{3}{\molar}, \SI{1}{\micro\molar})
\DeclareSIUnit\an{an} % année (le modèle confond parfois avec \year, un compteur TeX) — ex. \SI{5}{\kilo\meter\per\an}
\DeclareSIUnit\FCFA{FCFA} % franc CFA (ex. \SI{500}{\FCFA\per\kilo\gram})
\DeclareSIUnit\degreeBrix{\textdegree Brix} % teneur en sucre d'un sirop/jus (réfractométrie)
\DeclareSIUnit\month{mois} % le modèle utilise parfois \month, un compteur TeX, comme unité de durée
\usepackage{qrcode}
% \degree : commande fréquemment utilisée par les modèles pour un angle ou une
% température en dehors de \SI{}{} (non fournie par les paquets chargés).
\providecommand{\degree}{^{\circ}}
% \qed (fin de démonstration), utilisé par les modèles sans amsthm
\providecommand{\qed}{\hfill$\square$}
% \barre : double barre des valeurs interdites dans un tableau de variations (tkz-tab)
\providecommand{\barre}{\ensuremath{\|}}
% environnement proof (amsthm non chargé) utilisé par les modèles
\ifdefined\proof\else\newenvironment{proof}[1][Démonstration]{\par\noindent\textit{#1.}\ }{\qed\par}\fi
\usepackage{lastpage}
\usepackage{hyperref}
\hypersetup{hidelinks}

% ============================================================
% Palette « Pop » OnBuch+ — mêmes hex que la classe P (plus_ui.dart)
% ============================================================
\definecolor{poporange}{HTML}{F59321}
\definecolor{popdark}{HTML}{E07A0C}
\definecolor{popcream}{HTML}{FFF6E8}
\definecolor{popink}{HTML}{1C1714}
\definecolor{poppurple}{HTML}{7A5AE0}
\definecolor{popgreen}{HTML}{2E9E6A}
\definecolor{poppink}{HTML}{E0457B}
\definecolor{popblue}{HTML}{2F7DE1}
\definecolor{popgold}{HTML}{C98A12}
\definecolor{popmuted}{HTML}{8A7F76}
\definecolor{popline}{HTML}{EADFD2}
\definecolor{popgray}{HTML}{8A7F76}
\colorlet{popgrey}{popgray}
% Alias tolérants (le modèle nomme parfois les couleurs différemment)
\colorlet{popwhite}{white}
\colorlet{popblack}{popink}
\colorlet{popred}{poppink}
\colorlet{popyellow}{popgold}
% Teintes claires (fonds de tableaux/encadrés) — le modèle nomme parfois
% les couleurs avec un suffixe L (ex. popblueL) sans les avoir définies.
\colorlet{poporangeL}{poporange!20}
\colorlet{popdarkL}{popdark!20}
\colorlet{poppurpleL}{poppurple!20}
\colorlet{popgreenL}{popgreen!20}
\colorlet{poppinkL}{poppink!20}
\colorlet{popblueL}{popblue!20}
\colorlet{popgoldL}{popgold!20}
\colorlet{popredL}{popred!20}
\colorlet{popgrayL}{popgray!20}
% Teintes claires pour fonds de tableaux / zones
\colorlet{popblueL}{popblue!10!white}
\colorlet{poporangeL}{poporange!14!white}
\colorlet{popgreenL}{popgreen!12!white}
\colorlet{poppurpleL}{poppurple!12!white}
\colorlet{poppinkL}{poppink!12!white}
\colorlet{popgoldL}{popgold!14!white}

\pagecolor{popcream}

% ============================================================
% Typographie
% ============================================================
\defaultfontfeatures{Ligatures=TeX}
\setmainfont{PlusJakartaSans-Medium.ttf}[Path=../../../fonts/,BoldFont=PlusJakartaSans-Bold.ttf]
% Maths en sans-serif (Fira Math) avec chiffres et lettres droites de Plus
% Jakarta, pour que les formules se fondent dans le texte.
\usepackage[math-style=ISO,bold-style=ISO]{unicode-math}
\setmathfont{FiraMath-Regular.otf}[Path=../../../fonts/]
\setmathfont{PlusJakartaSans-Medium.ttf}[Path=../../../fonts/,range={up/num,up/{Latin,latin},\mathup}]
\usepackage{icomma} % virgule décimale sans espace en mode math
% Caractères mathématiques Unicode absents des polices de texte (Plus Jakarta,
% Archivo Black) : rendus par leur équivalent mathématique, en texte comme en
% formule — sinon ils s'affichent en carré vide (ex. « Arithmétique dans ℤ »).
\usepackage{newunicodechar}
\newunicodechar{ℝ}{\ensuremath{\mathbb{R}}}
\newunicodechar{ℤ}{\ensuremath{\mathbb{Z}}}
\newunicodechar{ℂ}{\ensuremath{\mathbb{C}}}
\newunicodechar{ℕ}{\ensuremath{\mathbb{N}}}
\newunicodechar{ℚ}{\ensuremath{\mathbb{Q}}}
\newunicodechar{𝔻}{\ensuremath{\mathbb{D}}}
\newunicodechar{α}{\ensuremath{\alpha}}
\newunicodechar{β}{\ensuremath{\beta}}
\newunicodechar{γ}{\ensuremath{\gamma}}
\newunicodechar{δ}{\ensuremath{\delta}}
\newunicodechar{ε}{\ensuremath{\varepsilon}}
\newunicodechar{θ}{\ensuremath{\theta}}
\newunicodechar{λ}{\ensuremath{\lambda}}
\newunicodechar{μ}{\ensuremath{\mu}}
\newunicodechar{σ}{\ensuremath{\sigma}}
\newunicodechar{τ}{\ensuremath{\tau}}
\newunicodechar{φ}{\ensuremath{\varphi}}
\newunicodechar{ψ}{\ensuremath{\psi}}
\newunicodechar{ω}{\ensuremath{\omega}}
\newunicodechar{Σ}{\ensuremath{\Sigma}}
% majuscules produites par \MakeUppercase dans les titres (α-aminés -> Α)
\newunicodechar{Α}{\ensuremath{\alpha}}
\newunicodechar{Β}{\ensuremath{\beta}}
\newunicodechar{Γ}{\ensuremath{\Gamma}}
\newunicodechar{Δ}{\ensuremath{\Delta}}
\newunicodechar{Ε}{\ensuremath{\varepsilon}}
\newunicodechar{Θ}{\ensuremath{\Theta}}
\newunicodechar{Λ}{\ensuremath{\Lambda}}
\newunicodechar{Μ}{\ensuremath{\mu}}
\newunicodechar{Τ}{\ensuremath{\tau}}
\newunicodechar{Φ}{\ensuremath{\Phi}}
\newunicodechar{Ψ}{\ensuremath{\Psi}}
\newunicodechar{Ω}{\ensuremath{\Omega}}
\newunicodechar{↦}{\ensuremath{\mapsto}}
\newunicodechar{∈}{\ensuremath{\in}}
\newunicodechar{∉}{\ensuremath{\notin}}
\newunicodechar{∧}{\ensuremath{\wedge}}
\newunicodechar{⇔}{\ensuremath{\Leftrightarrow}}
\newunicodechar{⟺}{\ensuremath{\Longleftrightarrow}}
\newunicodechar{⇒}{\ensuremath{\Rightarrow}}
\newunicodechar{⟹}{\ensuremath{\Longrightarrow}}
\newunicodechar{∘}{\ensuremath{\circ}}
\newunicodechar{∪}{\ensuremath{\cup}}
\newunicodechar{∩}{\ensuremath{\cap}}
\newunicodechar{≡}{\ensuremath{\equiv}}
\newunicodechar{⊂}{\ensuremath{\subset}}
\newunicodechar{⊥}{\ensuremath{\perp}}
\newunicodechar{∃}{\ensuremath{\exists}}
\newunicodechar{∀}{\ensuremath{\forall}}
\newunicodechar{∛}{\ensuremath{\sqrt[3]{\,}}}
\newunicodechar{∜}{\ensuremath{\sqrt[4]{\,}}}
\newunicodechar{⃗}{\ensuremath{{}^{\to}}}
\newunicodechar{ⁿ}{\ifmmode^{n}\else\textsuperscript{\ensuremath{n}}\fi}
\newunicodechar{ˣ}{\ifmmode^{x}\else\textsuperscript{\ensuremath{x}}\fi}
\newunicodechar{ᵇ}{\ifmmode^{b}\else\textsuperscript{\ensuremath{b}}\fi}
\newunicodechar{ʳ}{\ifmmode^{r}\else\textsuperscript{\ensuremath{r}}\fi}
\newunicodechar{ᵘ}{\ifmmode^{u}\else\textsuperscript{\ensuremath{u}}\fi}
\newunicodechar{ʸ}{\ifmmode^{y}\else\textsuperscript{\ensuremath{y}}\fi}
\newunicodechar{ᵃ}{\ifmmode^{a}\else\textsuperscript{\ensuremath{a}}\fi}
\newunicodechar{ᵉ}{\ifmmode^{e}\else\textsuperscript{\ensuremath{e}}\fi}
\newunicodechar{ᵏ}{\ifmmode^{k}\else\textsuperscript{\ensuremath{k}}\fi}
\newunicodechar{ᵖ}{\ifmmode^{p}\else\textsuperscript{\ensuremath{p}}\fi}
\newunicodechar{ᴺ}{\ifmmode^{N}\else\textsuperscript{\ensuremath{N}}\fi}
\newunicodechar{ᶜ}{\ifmmode^{c}\else\textsuperscript{\ensuremath{c}}\fi}
\newunicodechar{ʰ}{\ifmmode^{h}\else\textsuperscript{\ensuremath{h}}\fi}
\newunicodechar{ᵠ}{\ifmmode^{q}\else\textsuperscript{\ensuremath{q}}\fi}
\newunicodechar{ₙ}{\ifmmode_{n}\else\textsubscript{\ensuremath{n}}\fi}
\newunicodechar{ₐ}{\ifmmode_{a}\else\textsubscript{\ensuremath{a}}\fi}
\newunicodechar{ᵢ}{\ifmmode_{i}\else\textsubscript{\ensuremath{i}}\fi}
\newunicodechar{ᵣ}{\ifmmode_{r}\else\textsubscript{\ensuremath{r}}\fi}
\newunicodechar{ₖ}{\ifmmode_{k}\else\textsubscript{\ensuremath{k}}\fi}
\newunicodechar{ᵦ}{\ifmmode_{\beta}\else\textsubscript{\ensuremath{\beta}}\fi}
\newunicodechar{ₓ}{\ifmmode_{x}\else\textsubscript{\ensuremath{x}}\fi}
\newfontfamily\popdisplay{ArchivoBlack-Regular.ttf}[Path=../../../fonts/]
\newfontfamily\poplogo{SpaceGrotesk-Bold.ttf}[Path=../../../fonts/]
\newfontfamily\popsemi{PlusJakartaSans-SemiBold.ttf}[Path=../../../fonts/]
\newfontfamily\popxbold{PlusJakartaSans-ExtraBold.ttf}[Path=../../../fonts/]
\newfontfamily\popxboldls{PlusJakartaSans-ExtraBold.ttf}[Path=../../../fonts/,LetterSpace=3.0]

\setlist[enumerate]{leftmargin=1.7em,itemsep=5pt,topsep=4pt}
\setlist[itemize]{leftmargin=1.7em,itemsep=4pt,topsep=4pt,label={\color{poporange}\textbullet}}
\setlength{\parindent}{0pt}
\setlength{\parskip}{5pt}
\renewcommand{\arraystretch}{1.25}

% pgfplots : style Pop par défaut
\pgfplotsset{
  every axis/.append style={axis line style={popink,line width=1pt},tick style={popink},
    grid style={popline,line width=0.5pt},label style={font=\small},tick label style={font=\footnotesize}},
  popcurve/.style={poporange,line width=1.8pt,no markers,smooth},
}
% Même style disponible dans un \draw TikZ simple (hors pgfplots)
\tikzset{popcurve/.style={poporange,line width=1.8pt}}
\tikzset{popgrid/.style={popline,very thin}}
\colorlet{popcurve}{poporange} % le nom du style est parfois utilisé comme couleur

% ============================================================
% Pill / squiggle
% ============================================================
\newcommand{\poppill}[3]{%
  \tikz[baseline=-0.55ex]\node[draw=popink,line width=1.2pt,rounded corners=2.6pt,%
    fill=#2,inner xsep=6.5pt,inner ysep=3pt]{%
    \popxboldls\fontsize{7}{7}\selectfont\color{#3}\MakeUppercase{#1}};%
}
\newcommand{\popsquiggle}[1]{%
  \tikz[baseline=-0.4ex]{
    \draw[#1,line width=2.6pt,line cap=round]
      plot [smooth] coordinates {(0,0.09) (0.34,0.01) (0.68,0.21) (1.02,0.01) (1.36,0.10)};
  }%
}
% Mot-clé surligné (marqueur orange sous le texte)
\newcommand{\cle}[1]{{\popxbold\color{popdark}#1}}

% ============================================================
% En-tête / pied
% ============================================================
\pagestyle{fancy}
\fancyhf{}
\renewcommand{\headrulewidth}{0pt}
\renewcommand{\footrulewidth}{0pt}
\lhead{\raisebox{-0.15ex}{\poplogo\fontsize{13}{13}\selectfont\color{popink}OnBuch\color{poporange}+}}
\rhead{\poppill{\DOCMATIERE\ \textbullet\ \DOCNIVEAU\ \textbullet\ Leçon \DOCLECON}{white}{popink}}
\lfoot{\popxbold\fontsize{7.6}{7.6}\selectfont\color{popmuted}\MakeUppercase{Cours signé OnBuch+}\ \textbullet\ {\popsemi\DOCTITRE}}
\rfoot{\tikz[baseline=-0.6ex]\node[rounded corners=8.5pt,fill=popink,minimum height=17pt,minimum width=17pt,inner xsep=5pt,inner sep=0]{\popxbold\fontsize{8}{8}\selectfont\color{popcream}\thepage\,/\,\pageref*{LastPage}};}

% ============================================================
% Titres
% ============================================================
\newcommand{\chaptitre}[2]{%
  \begin{tcolorbox}[enhanced,colback=poporange,colframe=popink,boxrule=2.6pt,arc=11pt,
      drop shadow=popink,left=15pt,right=15pt,top=12pt,bottom=11pt,before skip=3pt,after skip=18pt]
    \poppill{#2}{popink}{popcream}\par\vspace{9pt}
    {\raggedright\hyphenpenalty=10000\exhyphenpenalty=10000\popdisplay\fontsize{22}{24}\selectfont\color{popcream}\MakeUppercase{#1}\par}
  \end{tcolorbox}
}
% Section numérotée : pastille orange avec le numéro + titre Archivo
\newcounter{popsec}
\newcommand{\coursec}[1]{%
  \stepcounter{popsec}\setcounter{popsub}{0}%
  \par\vspace{16pt}\noindent
  \begin{minipage}{\linewidth}
  \tikz[baseline=-0.7ex]\node[draw=popink,line width=1.6pt,fill=poporange,rounded corners=4pt,minimum size=22pt,inner sep=0]{\popdisplay\fontsize{12}{12}\selectfont\color{popcream}\thepopsec};%
  \hspace{8pt}{\popdisplay\fontsize{15}{17}\selectfont\color{popink}\MakeUppercase{#1}}\par
  \vspace{4pt}\noindent{\color{poporange}\rule{36pt}{3.6pt}}
  \end{minipage}\par\nopagebreak\vspace{8pt}
}
\newcounter{popsub}
\newcommand{\courssub}[1]{%
  \stepcounter{popsub}%
  \par\vspace{9pt}\noindent{\popxbold\fontsize{11.5}{13}\selectfont\color{popink}{\color{poporange}\thepopsec.\thepopsub}\hspace{6pt}#1}\par\nopagebreak\vspace{4pt}
}

% ============================================================
% Boîtes pédagogiques — même recette : fond blanc, bordure épaisse colorée,
% ombre dure encre, badge plein. Titre optionnel [#1].
% ============================================================
\tcbset{popbase/.style={breakable,enhanced,colback=white,boxrule=2.3pt,arc=9pt,
  shadow={3pt}{-3pt}{0pt}{popink},left=14pt,right=14pt,top=11pt,bottom=11pt,before skip=11pt,after skip=15pt,
  fontupper=\popsemi\fontsize{10.6}{14.5}\selectfont\color{popink},
  fontlower=\popsemi\fontsize{10.6}{14.5}\selectfont\color{popink}}}
% Coche dessinée (le glyphe ✓ n'existe pas dans les polices du document)
\newcommand{\popcheck}[1]{\tikz[baseline=-0.5ex]\draw[#1,line width=1.6pt,line cap=round,line join=round](-0.13,0.0)--(-0.04,-0.09)--(0.13,0.1);}
\providecommand{\coche}{\popcheck{popgreen}}
\providecommand{\resultat}[1]{\par\smallskip\noindent\fcolorbox{popgreen}{popgreenL}{\parbox{\dimexpr\linewidth-2\fboxsep-2\fboxrule\relax}{\textbf{Résultat.} #1}}\par\smallskip}
\newcommand{\popboxhead}[3]{\poppill{#1}{#2}{white}\ifx\relax#3\relax\else\hspace{6pt}{\popxbold\fontsize{10.6}{12}\selectfont\color{popink}#3}\fi\par\vspace{7pt}}

\newtcolorbox{aretenir}[1][]{popbase,colframe=popgreen,before upper={\popboxhead{À retenir}{popgreen}{#1}}}
% Texte en anglais (document de lecture, dialogue, extrait) : cadre
% d'étude. Un titre optionnel (source, auteur) s'affiche dans l'en-tête.
\newtcolorbox{textebox}[1][]{popbase,colframe=popblue,before upper={\popboxhead{Text}{popblue}{#1}\begin{otherlanguage}{english}},after upper={\end{otherlanguage}}}
% Marques ✓ / ✗ (forme correcte / fautive) souvent utilisées par les modèles
\providecommand{\cross}{\textcolor{poppink}{\ensuremath{\times}}}
\providecommand{\xmark}{\cross}\providecommand{\crossmark}{\cross}\providecommand{\cmark}{\ensuremath{\checkmark}}
% Ligne de réponse / justification à compléter (inventée par les modèles)
\providecommand{\justification}[1]{\par\noindent\textit{Justification~:} \dotfill\par}
% \textipa (package tipa non chargé) : la police ne couvre pas l'API -> simple passage
\providecommand{\textipa}[1]{#1}
% Soulignement (ulem non chargé) : \uline, \ul
\providecommand{\uline}[1]{\underline{#1}}\providecommand{\ul}[1]{\underline{#1}}
% \glossaire{\item ...} inventée par les modèles : liste de glossaire
\providecommand{\glossaire}[1]{\begin{itemize}#1\end{itemize}}
% \alert (beamer) inventée par les modèles : texte mis en évidence
\providecommand{\alert}[1]{\textbf{\textcolor{poppink}{#1}}}
% variantes ASCII d'ordinaux écrites en macro
\providecommand{\iere}{\textsuperscript{re}}\providecommand{\ieme}{\textsuperscript{e}}\providecommand{\ere}{\textsuperscript{re}}\providecommand{\eme}{\textsuperscript{e}}\providecommand{\er}{\textsuperscript{er}}\providecommand{\ie}{\textsuperscript{e}}
% Ordinaux français écrits en macro par les modèles : 1\ière, 3\ième
\newcommand{\ière}{\textsuperscript{re}}\newcommand{\ième}{\textsuperscript{e}}
% \exemple (inventée par les modèles) : étiquette « Exemple : »
\providecommand{\exemple}{\textbf{Exemple~:}\ }
% \square utilisable aussi hors mode math (cases à cocher)
\AtBeginDocument{\let\squareMath\square\renewcommand{\square}{\ensuremath{\squareMath}}}
% Mot ou phrase en anglais à l'intérieur d'un paragraphe français
\newcommand{\eng}[1]{\ifmmode\text{\textit{#1}}\else\begingroup\selectlanguage{english}\itshape #1\endgroup\fi}
\let\esp\eng % alias toléré
% flèches d'intonation inventées par les modèles
\providecommand{\textsearrow}{}\renewcommand{\textsearrow}{\ensuremath{\searrow}}\providecommand{\textnearrow}{}\renewcommand{\textnearrow}{\ensuremath{\nearrow}}
% \fr{..} : mot français (inventé par les modèles)
\providecommand{\fr}[1]{\textit{#1}}
\newtcolorbox{exemplebox}[1][]{popbase,colframe=poporange,before upper={\popboxhead{Exemple}{poporange}{#1}}}
\newtcolorbox{definition}[1][]{popbase,colframe=popblue,before upper={\popboxhead{Définition}{popblue}{#1}}}
\newtcolorbox{propriete}[1][]{popbase,colframe=poppurple,before upper={\popboxhead{Propriété}{poppurple}{#1}}}
\newtcolorbox{methode}[1][]{popbase,colframe=popgold,before upper={\popboxhead{Méthode}{popgold}{#1}}}
\newtcolorbox{attention}[1][]{popbase,colframe=poppink,before upper={\popboxhead{Attention}{poppink}{#1}}}
\newtcolorbox{experience}[1][]{popbase,colframe=popblue,colback=popblueL,before upper={\popboxhead{Expérience}{popblue}{#1}}}
% « Le savais-tu ? » : carte violette pleine (contexte, culture, Cameroun)
\newtcolorbox{savaistu}[1][]{popbase,colback=poppurple,colframe=popink,
  fontupper=\popsemi\fontsize{10.4}{14.2}\selectfont\color{white},
  before upper={\poppill{Le savais-tu\,?}{white}{poppurple}\ifx\relax#1\relax\else\hspace{6pt}{\popxbold\color{white}#1}\fi\par\vspace{7pt}}}
% Exercice résolu : énoncé en haut, \tcblower puis la solution sur fond crème
\newcounter{popexo}\newcounter{popexr}
\newtcolorbox{exoresolu}[1][]{popbase,colframe=poporange,colbacklower=popcream,
  segmentation style={popink,line width=1.2pt,dashed},
  before upper={\stepcounter{popexr}\popboxhead{Exercice résolu \thepopexr}{poporange}{#1}},
  before lower={\poppill{Solution}{popink}{popcream}\par\vspace{6pt}}}
% Exercice (non résolu en place) — la correction est dans la partie Corrigés
\newtcolorbox{exercice}[1][]{popbase,colframe=popink,
  before upper={\stepcounter{popexo}\popboxhead{Exercice \thepopexo}{popink}{#1}}}
% Corrigé
\newtcolorbox{corrige}[1][]{popbase,colframe=popgreen,colback=popgreenL,
  before upper={\popboxhead{Corrigé}{popgreen}{#1}}}
% Objectifs / prérequis / bilan
\newtcolorbox{objectifs}{popbase,colframe=popink,colback=poporangeL,
  before upper={\popboxhead{Objectifs de la leçon}{popink}{}}}
\newtcolorbox{prerequis}{popbase,colframe=popink,colback=popblueL,
  before upper={\popboxhead{Prérequis}{popblue}{}}}
\newtcolorbox{bilan}{popbase,colframe=popink,colback=white,boxrule=2.8pt,
  before upper={\popboxhead{Fiche bilan}{poporange}{L'essentiel à connaître pour l'examen}}}
% Figure encadrée avec légende
\newtcolorbox{popfigure}[1][]{enhanced,breakable=false,colback=white,colframe=popline,boxrule=1.4pt,arc=8pt,
  left=10pt,right=10pt,top=10pt,bottom=8pt,before skip=10pt,after skip=12pt,halign=center,
  lower separated=false,halign lower=center,
  fontlower=\popsemi\fontsize{9}{11.5}\selectfont\color{popmuted},
  #1}
\newcommand{\legende}[1]{\par\vspace{6pt}{\popsemi\fontsize{9}{11.5}\selectfont\color{popmuted}#1}}

% ============================================================
% Signature OnBuch+
% ============================================================
\newcommand{\poplogoinline}[1]{{\poplogo\fontsize{#1}{#1}\selectfont\color{popink}OnBuch\color{poporange}+}}
\newcommand{\signatureonbuch}{%
  \begin{tcolorbox}[enhanced,colback=popink,colframe=popink,boxrule=2.2pt,arc=10pt,
      drop shadow=poporange,left=14pt,right=14pt,top=10pt,bottom=10pt,before skip=0pt,after skip=0pt]
    \tikz[baseline=-0.7ex]\node[circle,fill=popgreen,draw=popcream,line width=1.3pt,minimum size=22pt,inner sep=0]{\popcheck{white}};%
    \hspace{9pt}\begin{minipage}[c]{0.62\linewidth}
      {\popxbold\fontsize{12}{13}\selectfont\color{popcream}Signé\ \textbullet\ {\poplogo OnBuch\color{poporange}+}}\par\vspace{2pt}
      {\raggedright\popsemi\fontsize{8}{10.5}\selectfont\color{popline}\MakeUppercase{Équipe pédagogique OnBuch+}\\\MakeUppercase{Conforme au programme officiel MINESEC}\par}
    \end{minipage}\hfill
    {\poplogo\fontsize{22}{22}\selectfont\color{popcream}OnBuch\color{poporange}+}
  \end{tcolorbox}%
}

% ============================================================
% Page de garde
% #1 titre · #2 matière · #3 niveau · #4 module · #5 n° leçon · #6 durée ·
% #7 description / objectifs (texte ou itemize)
% ============================================================
\newcommand{\pagegarde}[7]{%
  \thispagestyle{empty}
  \newgeometry{a4paper,margin=1.7cm,top=1.6cm,bottom=1.4cm}%
  \begin{tikzpicture}[remember picture,overlay]
    \fill[poporange!18] ([xshift=-3.2cm,yshift=-3.6cm]current page.north east) circle (4.6cm);
    \fill[poppurple!14] ([xshift=2.4cm,yshift=7.6cm]current page.south west) circle (3.8cm);
  \end{tikzpicture}%
  {\poplogo\fontsize{30}{30}\selectfont\color{popink}OnBuch\color{poporange}+}\hfill
  \raisebox{0.6ex}{\poppill{Cours signé \textbullet\ Premium}{popink}{popcream}}\par
  \vspace{1pt}\popsquiggle{poppurple}\par
  \vspace{8pt}
  \begin{tcolorbox}[enhanced,colback=poporange,colframe=popink,boxrule=2.8pt,arc=13pt,
      drop shadow=popink,left=18pt,right=18pt,top=12pt,bottom=13pt,after skip=10pt]
    \poppill{#2 \textbullet\ #3}{popink}{popcream}\hspace{5pt}\poppill{#4}{white}{popink}\par\vspace{8pt}
    {\popdisplay\fontsize{14}{14}\selectfont\color{popink}LEÇON #5}\par\vspace{4pt}
    {\raggedright\hyphenpenalty=10000\exhyphenpenalty=10000\popdisplay\fontsize{25}{27}\selectfont\color{popcream}\MakeUppercase{#1}\par}
    \vspace{8pt}
    {\popxbold\fontsize{9.5}{9.5}\selectfont\color{popink}\MakeUppercase{Durée conseillée : #6}}
  \end{tcolorbox}
  \begin{tcolorbox}[enhanced,colback=white,colframe=popink,boxrule=2.2pt,arc=10pt,
      drop shadow=popink,left=16pt,right=16pt,top=10pt,bottom=10pt,after skip=8pt]
    \poppill{Dans cette leçon}{poppurple}{white}\par\vspace{5pt}
    \popsemi\fontsize{9.3}{12.6}\selectfont\color{popink}#7
  \end{tcolorbox}
  \noindent\begin{minipage}[t]{0.487\linewidth}
    \begin{tcolorbox}[enhanced,colback=popgreen,colframe=popink,boxrule=2.2pt,arc=10pt,
        drop shadow=popink,left=11pt,right=11pt,top=10pt,bottom=10pt,height=3.1cm,valign=center]
      \tcbox[colback=white,colframe=popink,boxrule=1.3pt,arc=5pt,boxsep=0pt,nobeforeafter,
        left=3pt,right=3pt,top=3pt,bottom=3pt,valign=center]{\qrcode[height=1.9cm]{https://play.google.com/store/apps/details?id=com.luvvix.onbuch&hl=fr}}%
      \hspace{8pt}\begin{minipage}[c]{0.5\linewidth}
        {\popdisplay\fontsize{11}{12.5}\selectfont\color{white}\MakeUppercase{Télécharge OnBuch}}\par\vspace{4pt}
        {\raggedright\popsemi\fontsize{8.4}{11}\selectfont\color{white}Gratuit \textbullet\ Google Play\\Tuteur IA, annales, concours, résultats\par}
      \end{minipage}
    \end{tcolorbox}
  \end{minipage}\hfill
  \begin{minipage}[t]{0.487\linewidth}
    \begin{tcolorbox}[enhanced,colback=poppurple,colframe=popink,boxrule=2.2pt,arc=10pt,
        drop shadow=popink,left=11pt,right=11pt,top=10pt,bottom=10pt,height=3.1cm,valign=center]
      \tikz[baseline=-0.7ex]\node[circle,fill=white,draw=popink,line width=1.3pt,minimum size=1.6cm,inner sep=0]{\popdisplay\fontsize{24}{24}\selectfont\color{poppurple}+};%
      \hspace{8pt}\begin{minipage}[c]{0.6\linewidth}
        {\popdisplay\fontsize{11}{12.5}\selectfont\color{white}\MakeUppercase{Passe à OnBuch+}}\par\vspace{4pt}
        {\raggedright\popsemi\fontsize{8.4}{11}\selectfont\color{white}Tous les cours signés, Tuteur IA illimité, fiches PDF — onglet OnBuch+ de l'app\par}
      \end{minipage}
    \end{tcolorbox}
  \end{minipage}
  \vspace{8pt}
  \popcontenu
  \vfill
  \signatureonbuch
  \restoregeometry
  \newpage
}

% Rangée « Ce cours contient » (page de garde)
\newcommand{\poptile}[3]{% #1 couleur #2 gros texte #3 légende
  \begin{tcolorbox}[enhanced,colback=#1,colframe=popink,boxrule=2pt,arc=9pt,shadow={2.5pt}{-2.5pt}{0pt}{popink},
      left=6pt,right=6pt,top=9pt,bottom=9pt,halign=center,valign=center,height=1.75cm,width=\linewidth,nobeforeafter]
    {\popdisplay\fontsize{12.5}{13}\selectfont\color{white}\MakeUppercase{#2}}\par\vspace{3pt}
    {\centering\popsemi\fontsize{7.8}{9.5}\selectfont\color{white}#3\par}
  \end{tcolorbox}}
\newcommand{\popcontenu}{%
  \par\noindent{\popxboldls\fontsize{7.5}{7.5}\selectfont\color{popmuted}CE COURS SIGNÉ CONTIENT}\par\vspace{7pt}
  \noindent\begin{minipage}[t]{0.235\linewidth}\poptile{popblue}{Cours}{complet, illustré, pas à pas}\end{minipage}\hfill
  \begin{minipage}[t]{0.235\linewidth}\poptile{poporange}{Méthodes}{et exercices résolus}\end{minipage}\hfill
  \begin{minipage}[t]{0.235\linewidth}\poptile{poppink}{Exercices}{type examen + corrigés}\end{minipage}\hfill
  \begin{minipage}[t]{0.235\linewidth}\poptile{popgreen}{Fiche bilan}{l'essentiel à retenir}\end{minipage}\par}

% Sommaire
\newtcolorbox{sommaire}{enhanced,colback=white,colframe=popink,boxrule=2.2pt,arc=10pt,
  drop shadow=popink,left=18pt,right=18pt,top=13pt,bottom=13pt,before skip=6pt,after skip=20pt,
  before upper={\poppill{Sommaire}{poppurple}{white}\par\vspace{9pt}\popsemi\fontsize{10.6}{14.5}\selectfont\color{popink}}}

% Signature de fin de cours — poussée tout en bas de la dernière page
\newcommand{\findecours}{%
  \par\vfill\nopagebreak
  \begin{center}\popsquiggle{poporange}\end{center}\vspace{4pt}
  \signatureonbuch
}
\AtBeginDocument{\def\llbracket{\mathopen{[\mkern-3.5mu[}}\def\rrbracket{\mathclose{]\mkern-3.5mu]}}} % intervalles d'entiers (glyphe ⟦ absent de la police)
% Glyphes absents de Fira Math : redéfinis à partir de glyphes disponibles
\AtBeginDocument{%
  \def\setminus{\mathbin{\backslash}}\let\smallsetminus\setminus
  \def\vdots{\mathord{\vbox{\baselineskip3.2pt\lineskiplimit0pt\kern4pt\hbox{.}\hbox{.}\hbox{.}}}}%
  \def\ddots{\mathinner{\mkern1mu\raise6.5pt\hbox{.}\mkern2mu\raise3.6pt\hbox{.}\mkern2mu\raise0.7pt\hbox{.}\mkern1mu}}%
  \def\triangle{\mathord{\scalebox{0.9}{$\symup{\Delta}$}}}%
  \def\nabla{\mathord{\raisebox{\depth}{\scalebox{1}[-1]{$\symup{\Delta}$}}}}%
  \def\checkmark{\popcheck{popdark}}%
  \def\star{\mathbin{\ast}}\def\bigstar{\mathord{\ast}}%
  \def\diamond{\mathbin{\scalebox{0.6}{$\Diamond$}}}%
  \def\bigcup{\mathop{\mathchoice{\raisebox{-0.3em}{\scalebox{1.7}{$\cup$}}}{\scalebox{1.25}{$\cup$}}{\cup}{\cup}}\displaylimits}%
  \def\bigcap{\mathop{\mathchoice{\raisebox{-0.3em}{\scalebox{1.7}{$\cap$}}}{\scalebox{1.25}{$\cap$}}{\cap}{\cap}}\displaylimits}%
  \def\lesseqqgtr{\mathrel{\lessgtr}}\def\bigoplus{\mathop{\oplus}\displaylimits}%
}
\providecommand{\ssi}{si et seulement si} % abréviation fréquente
\AtBeginDocument{\providecommand{\wideparen}{\overparen}} % arc de cercle

\begin{document}

\pagegarde{Lesson 50 — Grammar: General revision}{Anglais}{Tle A}{Chapter 5 : Media and Communication}{EN50}{2 h}{Leçon de révision générale de grammaire pour la Terminale A : synthèse des temps, de la voix passive, du discours rapporté, des modaux et des propositions relatives, appliqués au thème des médias et de la communication. Observation, règles, comparaison avec le français et entraînement type Bac.
\vspace{4pt}
\begin{multicols}{2}\small
\begin{itemize}
\item À la fin de cette leçon, je sais identifier dans un texte les principaux temps anglais (present simple, present continuous, past simple, present perfect, future) et justifier leur emploi
\item À la fin de cette leçon, je sais conjuguer correctement les verbes réguliers et les verbes irréguliers les plus fréquents
\item À la fin de cette leçon, je sais transformer une phrase active en phrase passive et inversement
\item À la fin de cette leçon, je sais rapporter des paroles au discours indirect avec les changements de temps, de pronoms et de marqueurs temporels
\item À la fin de cette leçon, je sais employer les modaux (can, must, should, may, might) pour exprimer capacité, obligation, conseil et probabilité
\item À la fin de cette leçon, je sais construire des propositions relatives avec who, which, whose, where
\end{itemize}\end{multicols}}

\begin{sommaire}
\begin{multicols}{2}
\begin{enumerate}
\item Observation et diagnostic : repérer les structures dans un texte média
\item Révision des temps verbaux
\item Voix passive et discours rapporté
\item Modaux, relatifs et comparaison
\item Comparaison français/anglais et pièges des francophones
\item Entraînement intégré type Bac
\item Activité d'intégration
\item Méthodes et exercices résolus
\item Exercices
\item Corrigés des exercices
\item Fiche bilan
\end{enumerate}
\end{multicols}
\end{sommaire}

\begin{objectifs}
\begin{itemize}
\item À la fin de cette leçon, je sais identifier dans un texte les principaux temps anglais (present simple, present continuous, past simple, present perfect, future) et justifier leur emploi
\item À la fin de cette leçon, je sais conjuguer correctement les verbes réguliers et les verbes irréguliers les plus fréquents
\item À la fin de cette leçon, je sais transformer une phrase active en phrase passive et inversement
\item À la fin de cette leçon, je sais rapporter des paroles au discours indirect avec les changements de temps, de pronoms et de marqueurs temporels
\item À la fin de cette leçon, je sais employer les modaux (can, must, should, may, might) pour exprimer capacité, obligation, conseil et probabilité
\item À la fin de cette leçon, je sais construire des propositions relatives avec who, which, whose, where
\item À la fin de cette leçon, je sais former comparatifs et superlatifs et éviter les doubles marques (more better)
\item À la fin de cette leçon, je sais repérer et corriger les erreurs typiques des francophones (s de la 3e personne, be + V-ing, ordre des mots, faux amis)
\item À la fin de cette leçon, je sais produire un court texte correct sur les médias en mobilisant plusieurs structures révisées
\end{itemize}
\end{objectifs}

\begin{prerequis}
\begin{itemize}
\item Les temps de base étudiés en Première : present simple, present continuous, past simple, present perfect
\item La voix passive (be + participe passé) vue en Première
\item Le discours rapporté (reported speech) et la concordance des temps
\item Les auxiliaires modaux et les pronoms relatifs
\item Le vocabulaire des médias et de la communication (chapitre 5)
\end{itemize}
\end{prerequis}

\newpage

\coursec{Observation et diagnostic : repérer les structures dans un texte média}

\courssub{Situation d'accroche : le post Facebook de Manka}

Pendant les vacances, Manka, élève de Tle A à Bamenda, écrit un post Facebook pour féliciter son cousin admis à l'université de Buea :

\begin{textebox}[Manka's Facebook post]
“I am knowing him since ten years, he has work hardly and yesterday he has win a scholarship. My cousin is more intelligent than any student I know. Congratulations bro!”
\end{textebox}

Ses amis anglophones sourient : presque chaque verbe contient une erreur de temps ou de forme ! En classe, le professeur projette le message et demande : \emph{combien de fautes pouvez-vous corriger ?}

\textbf{Problématique.} Cette leçon de révision générale va vous permettre de mobiliser, en un seul regard, tous les points de grammaire étudiés depuis la Seconde : \cle{temps verbaux}, \cle{voix passive}, \cle{discours rapporté}, \cle{modaux}, \cle{propositions relatives} et \cle{comparaison}. Objectif : écrire et parler un anglais correct et fluide, comme l'exige l'épreuve du Baccalauréat. Mais avant de réviser chaque point en détail, apprenons d'abord à \emph{diagnostiquer} un texte : savoir repérer une structure, c'est déjà savoir l'employer.

\courssub{Lecture et observation : un article sur les médias au Cameroun}

Lisons ce court article original. Il contient \textbf{volontairement} toutes les structures que nous allons réviser dans cette leçon. Lisez-le une première fois pour le sens, puis une deuxième fois en soulignant mentalement tous les verbes.

\begin{textebox}[“The changing face of Cameroonian media” — school press review]
The CRTV has broadcast news in English and French since 1987. For decades, families in Yaoundé and Douala gathered around the radio every evening to listen to the news. Nowadays, many young Cameroonians are getting their information from social media. Last week, a journalist from Buea said that the press was becoming more digital every year. Articles which are shared online can be read by millions of people in a few hours. A report on fake news was published by the Ministry of Communication last month, and it warned that false stories spread faster than true ones. Everyone must check sources before sharing news. If citizens verify information carefully, social media will remain a useful tool for democracy. Media literacy is now taught in many schools, and teachers believe that young people will become more responsible readers than their parents were.
\end{textebox}

\textbf{Glossaire.}

\begin{center}
\begin{tabularx}{\linewidth}{|l|l|X|}
\hline
\rowcolor{popblueL} Mot anglais & Nature & Traduction française \\
\hline
to broadcast & verbe (irrégulier : broadcast, broadcast, broadcast) & diffuser (radio, télévision) \\
\hline
news & nom (toujours singulier !) & les informations, les nouvelles \\
\hline
to gather & verbe & se rassembler \\
\hline
a source & nom & une source (d'information) \\
\hline
fake news & nom & fausses informations, intox \\
\hline
to spread & verbe (irrégulier : spread, spread, spread) & se répandre, diffuser \\
\hline
to warn & verbe & avertir, mettre en garde \\
\hline
to check / to verify & verbes & vérifier \\
\hline
media literacy & nom & éducation aux médias \\
\hline
\end{tabularx}
\end{center}

\begin{attention}[Attention au mot \eng{news}]
\eng{News} se termine par un \emph{s} mais est \textbf{singulier} : \eng{The news is good.} « Les nouvelles sont bonnes. » On ne dit jamais \eng{a news} ; pour compter, on dit \eng{a piece of news} « une nouvelle ». De même, \eng{information} est indénombrable : \eng{a piece of information}, jamais \eng{an information}.
\end{attention}

\textbf{Questions de compréhension (reading).}

\begin{enumerate}
\item Since when has the CRTV broadcast news in both official languages?
\item Where do many young Cameroonians get their information nowadays?
\item What did the journalist from Buea say about the press?
\item Who published a report on fake news, and when?
\item What must everyone do before sharing news?
\end{enumerate}

\courssub{Méthode : comment diagnostiquer un texte}

\begin{methode}[Diagnostiquer les structures grammaticales d'un texte]
\begin{enumerate}
\item \textbf{Souligner tous les verbes} du texte (y compris les auxiliaires : \eng{has}, \eng{are}, \eng{was}, \eng{can}, \eng{must}...).
\item \textbf{Nommer chaque temps} : present simple, present continuous, present perfect, past simple, futur avec \eng{will}...
\item \textbf{Justifier l'emploi} de chaque temps : habitude ? action en cours ? bilan avec lien au présent ? fait passé daté ?
\item \textbf{Repérer les structures spéciales} : voix passive (\eng{be + participe passé}), discours rapporté (\eng{said that...}), modaux (\eng{must}, \eng{can}, \eng{will}), propositions relatives (\eng{which}, \eng{who}, \eng{that}), comparatifs (\eng{more... than}, \eng{-er than}).
\item \textbf{Classer} chaque structure dans un tableau de révision pour savoir ce qu'il reste à travailler.
\end{enumerate}
\end{methode}

Appliquons cette méthode au texte, phrase par phrase.

\textbf{Étapes 1 et 2 : les temps verbaux du texte.}

\begin{center}
\begin{tabularx}{\linewidth}{|X|l|X|}
\hline
\rowcolor{popblueL} Verbe souligné & Temps & Justification de l'emploi \\
\hline
\eng{has broadcast} & present perfect & bilan : action commencée en 1987 et qui continue ; \eng{since} exige le present perfect \\
\hline
\eng{gathered} & past simple & habitude passée révolue (\eng{for decades}, époque terminée) \\
\hline
\eng{are getting} & present continuous & tendance en cours de nos jours (\eng{nowadays}) \\
\hline
\eng{said} & past simple & parole située à un moment précis (\eng{last week}) \\
\hline
\eng{was becoming} & past continuous (discours rapporté) & recul d'un temps après \eng{said that} \\
\hline
\eng{are shared / can be read} & present passif + modal passif & l'action est subie par les articles \\
\hline
\eng{was published} & past passif & fait passé daté (\eng{last month}), sujet non agent \\
\hline
\eng{spread} & present simple & vérité générale, constat \\
\hline
\eng{must check} & modal + base verbale & obligation générale \\
\hline
\eng{verify / will remain} & present simple + futur & conditionnel de type 1 : \eng{if + present, will + BV} \\
\hline
\eng{is taught} & present passif & habitude, sujet non agent \\
\hline
\eng{believe / will become} & present simple + futur & opinion et prédiction \\
\hline
\end{tabularx}
\end{center}

\begin{aretenir}[Justifier un temps : les quatre questions à se poser]
\begin{itemize}
\item \textbf{Habitude ou vérité générale} $\rightarrow$ present simple : \eng{False stories spread fast.} « Les fausses informations se répandent vite. »
\item \textbf{Action en cours au moment où l'on parle} $\rightarrow$ present continuous : \eng{Many young people are getting their news online.} « Beaucoup de jeunes s'informent en ligne (en ce moment, tendance actuelle). »
\item \textbf{Bilan, expérience, action liée au présent} $\rightarrow$ present perfect : \eng{The CRTV has broadcast news since 1987.} « La CRTV diffuse les informations depuis 1987. »
\item \textbf{Fait passé daté et terminé} $\rightarrow$ past simple : \eng{A report was published last month.} « Un rapport a été publié le mois dernier. »
\end{itemize}
\end{aretenir}

\textbf{Étape 4 : les structures spéciales du texte.}

\begin{itemize}
\item \textbf{Voix passive} : \eng{Articles which are shared online can be read by millions of people.} — le sujet \eng{articles} subit l'action ; l'agent est introduit par \eng{by}. Autre exemple : \eng{A report was published by the Ministry.} « Un rapport a été publié par le ministère. »
\item \textbf{Discours rapporté} : \eng{A journalist said that the press was becoming more digital every year.} — après un verbe au passé (\eng{said}), le temps rapporté recule : \eng{is becoming} $\rightarrow$ \eng{was becoming}. « Un journaliste a déclaré que la presse devenait plus numérique chaque année. »
\item \textbf{Modal} : \eng{Everyone must check sources before sharing news.} — \eng{must} exprime l'obligation ; il est suivi de la base verbale sans \eng{to}. « Chacun doit vérifier les sources avant de partager une information. »
\item \textbf{Proposition relative} : \eng{Articles which are shared online...} — \eng{which} reprend \eng{articles} (chose) ; pour une personne on utiliserait \eng{who} : \eng{The journalist who wrote this article lives in Buea.} « Le journaliste qui a écrit cet article vit à Buea. »
\item \textbf{Comparatif} : \eng{false stories spread faster than true ones} et \eng{more responsible readers than their parents}. — adjectif court : \eng{fast} $\rightarrow$ \eng{faster than} ; adjectif long : \eng{responsible} $\rightarrow$ \eng{more responsible than}. « Les fausses histoires se répandent plus vite que les vraies. »
\end{itemize}

\courssub{Carte mentale : les six points de la révision générale}

\begin{popfigure}
\begin{tikzpicture}[mindmap, grow cyclic, every node/.style={concept}, root concept/.append style={concept color=popblue, font=\bfseries, minimum size=3.2cm, align=center}, level 1/.append style={level distance=3.6cm, sibling angle=60}]
\node[root concept, text=white]{General Grammar Revision}
  child[concept color=poporange] { node[concept, text=white, align=center]{Tenses} }
  child[concept color=popgreen] { node[concept, text=white, align=center]{Passive voice} }
  child[concept color=poppurple] { node[concept, text=white, align=center]{Reported speech} }
  child[concept color=poppink] { node[concept, text=white, align=center]{Modals} }
  child[concept color=popgold] { node[concept, text=white, align=center]{Relatives} }
  child[concept color=popblue!60] { node[concept, text=white, align=center]{Comparison} };
\end{tikzpicture}
\legende{Carte mentale des six points de grammaire à réviser : temps, voix passive, discours rapporté, modaux, relatifs, comparaison.}
\end{popfigure}

\courssub{Correction collective : les cinq erreurs du post de Manka}

Revenons à l'accroche. Le post de Manka contient \textbf{cinq erreurs} : trois erreurs de temps ou de forme verbale, une confusion de vocabulaire et une erreur de comparaison. Cherchons-les avant de lire la correction.

\begin{exoresolu}[Corriger le post Facebook de Manka]
Énoncé. Corrigez les cinq erreurs du message : \eng{“I am knowing him since ten years, he has work hardly and yesterday he has win a scholarship. My cousin is more intelligent than any student I know.”}
\tcblower
Solution détaillée.
\begin{enumerate}
\item \eng{I am knowing him since ten years} $\rightarrow$ \eng{I have known him for ten years.} Deux erreurs en une : le verbe \eng{to know} est un \textbf{verbe d'état} qui ne se met pas à la forme continue ; et « depuis dix ans » se traduit par \eng{for ten years} (durée) avec un \textbf{present perfect}. « Je le connais depuis dix ans. »
\item \eng{he has work} $\rightarrow$ \eng{he has worked.} Après \eng{has}, il faut le \textbf{participe passé} : \eng{worked}. « Il a travaillé. »
\item \eng{hardly} $\rightarrow$ \eng{hard.} Faux ami célèbre : \eng{hardly} signifie « à peine » ; l'adverbe « dur » est \eng{hard} (identique à l'adjectif). \eng{He has worked hard.} « Il a travaillé dur. »
\item \eng{yesterday he has win} $\rightarrow$ \eng{yesterday he won.} Avec un repère de temps passé précis (\eng{yesterday}), le present perfect est impossible : on utilise le \textbf{past simple}. De plus, le participe passé de \eng{win} est \eng{won} (win, won, won). « Hier, il a gagné une bourse. »
\item \eng{more intelligent than any student I know} $\rightarrow$ \eng{the most intelligent student I know.} Le sens voulu (« le plus intelligent de tous ») appelle le \textbf{superlatif} : \eng{the most + adjectif long}. « L'étudiant le plus intelligent que je connaisse. » (Remarque : avec le comparatif, il faudrait dire \eng{more intelligent than any other student}, car on ne se compare pas à soi-même.)
\end{enumerate}
Version corrigée complète : \eng{“I have known him for ten years, he has worked hard and yesterday he won a scholarship. My cousin is the most intelligent student I know. Congratulations bro!”}
\end{exoresolu}

\begin{exemplebox}[Exemples traduits : formes fausses et formes correctes]
\begin{itemize}
\item \eng{I am knowing him.} (faux) $\rightarrow$ \eng{I have known him for ten years.} « Je le connais depuis dix ans. »
\item \eng{He has work hardly.} (faux) $\rightarrow$ \eng{He has worked hard.} « Il a travaillé dur. »
\item \eng{Yesterday he has won.} (faux) $\rightarrow$ \eng{Yesterday he won.} « Hier, il a gagné. »
\item \eng{She is more tall than her brother.} (faux) $\rightarrow$ \eng{She is taller than her brother.} « Elle est plus grande que son frère. »
\item \eng{The news was announced by CRTV.} (correct) « La nouvelle a été annoncée par la CRTV. » — voix passive au passé.
\end{itemize}
\end{exemplebox}

\begin{attention}[Piège des francophones : « depuis » et le present perfect]
En français, « depuis » s'emploie avec le présent : \emph{j'habite ici depuis 2020}. En anglais, c'est impossible : on n'utilise jamais le present simple avec \eng{since} ou \eng{for} pour une action qui dure encore. L'anglais exige le \textbf{present perfect} :
\begin{itemize}
\item \eng{I have lived here since 2020.} « J'habite ici depuis 2020. » (\eng{since} + point de départ)
\item \eng{I have lived here for four years.} « J'habite ici depuis quatre ans. » (\eng{for} + durée)
\item Faux : \eng{I live here since 2020.} / \eng{I am living here since 2020.}
\end{itemize}
Autre confusion fréquente : \eng{since} (depuis, point de départ) contre \eng{for} (durée) contre \eng{ago} (il y a, avec le past simple) : \eng{He won the scholarship two days ago.} « Il a gagné la bourse il y a deux jours. »
\end{attention}

\courssub{Bilan : la liste des points à réviser}

Le diagnostic du texte et la correction du post de Manka nous donnent le programme des sections suivantes de cette leçon :

\begin{center}
\begin{tabularx}{\linewidth}{|l|X|X|}
\hline
\rowcolor{popblueL} Point de grammaire & Exemple du texte & Question clé à retenir \\
\hline
Tenses (temps) & \eng{has broadcast / are getting / gathered} & habitude, action en cours, bilan ou fait daté ? \\
\hline
Passive voice & \eng{was published by the Ministry} & le sujet fait-il ou subit-il l'action ? \\
\hline
Reported speech & \eng{said that the press was becoming} & comment les temps reculent-ils ? \\
\hline
Modals & \eng{must check / can be read} & obligation, possibilité, conseil ? \\
\hline
Relatives & \eng{articles which are shared} & \eng{who}, \eng{which} ou \eng{that} ? \\
\hline
Comparison & \eng{faster than / more digital} & adjectif court (\eng{-er}) ou long (\eng{more}) ? \\
\hline
\end{tabularx}
\end{center}

\begin{exercice}[À vous de jouer : diagnostic rapide]
Dans les phrases suivantes, soulignez le verbe, nommez la structure et justifiez l'emploi :
\begin{enumerate}
\item \eng{Students in Bamenda have used this library since 2015.}
\item \eng{The match was watched by thousands of fans last Saturday.}
\item \eng{The teacher said that the exam would be easier than the test.}
\item \eng{You should listen to English radio every day.}
\item \eng{The girl who won the prize lives in Douala.}
\end{enumerate}
\end{exercice}

\begin{corrige}[Exercice : diagnostic rapide]
\begin{enumerate}
\item \eng{have used} : present perfect — bilan d'une action commencée en 2015 et qui continue (\eng{since}).
\item \eng{was watched} : voix passive au past simple — le sujet \eng{the match} subit l'action ; repère passé \eng{last Saturday} ; agent introduit par \eng{by}.
\item \eng{said / would be} : discours rapporté — après \eng{said}, \eng{will} recule à \eng{would} ; comparatif \eng{easier than} (adjectif court).
\item \eng{should listen} : modal \eng{should} + base verbale — expression du conseil.
\item \eng{who won} : proposition relative — \eng{who} reprend \eng{the girl} (personne) ; \eng{won} au past simple (fait passé).
\end{enumerate}
\end{corrige}

\begin{aretenir}[Ce qu'il faut retenir de cette section]
Diagnostiquer un texte, c'est souligner les verbes, nommer les temps, justifier chaque emploi (habitude, action en cours, bilan, fait passé daté) puis repérer les structures spéciales : passive, discours rapporté, modaux, relatives, comparaison. Le post de Manka nous a rappelé trois réflexes essentiels : les verbes d'état comme \eng{know} refusent la forme continue ; « depuis » exige le present perfect avec \eng{since} ou \eng{for} ; un repère passé précis comme \eng{yesterday} impose le past simple. Nous pouvons maintenant passer à la révision détaillée des temps verbaux.
\end{aretenir}

\coursec{Révision des temps verbaux}

Dans la section précédente, nous avons repéré dans un texte de presse plusieurs structures verbales : \eng{the minister announced}, \eng{journalists are investigating}, \eng{the channel has broadcast}. Pourquoi le rédacteur a-t-il choisi tel temps plutôt qu'un autre ? C'est toute la question de cette section. En anglais, le choix du temps verbal n'est pas une affaire de style : c'est un choix de \cle{sens}. Chaque temps répond à une question précise : l'action est-elle habituelle ? en cours ? terminée et datée ? encore liée au présent ? Révisons systématiquement les sept temps essentiels du programme de Terminale, avec leurs formes, leurs emplois, leurs marqueurs temporels et les pièges qu'ils tendent aux francophones.

\courssub{Vue d'ensemble : la frise des temps}

Avant d'entrer dans le détail, visualisons l'ensemble. L'anglais organise les temps autour de trois zones : le passé, le moment présent (\eng{now}) et le futur. Chaque temps « habite » une zone précise.

\begin{popfigure}
\begin{tikzpicture}[scale=1.0]
\draw[-{Stealth}, thick, popdark] (-0.5,0) -- (12.5,0);
\fill[popblue] (2,0) circle (0.09);
\fill[poporange] (6,0) circle (0.14);
\fill[popgreen] (10,0) circle (0.09);
\node[below, popblue, font=\bfseries] at (2,-0.15) {PAST};
\node[below, poporange, font=\bfseries] at (6,-0.15) {NOW};
\node[below, popgreen, font=\bfseries] at (10,-0.15) {FUTURE};
% past simple : point passé
\fill[popblue] (1.2,0) circle (0.12);
\node[above, popblue, align=center, font=\small] at (1.2,0.35) {past simple\\ \eng{I saw it yesterday}};
% present perfect : flèche du passé vers NOW
\draw[-{Stealth}, very thick, poppurple] (3.2,0.9) .. controls (4.4,1.5) and (5.2,1.5) .. (5.9,0.25);
\node[above, poppurple, align=center, font=\small] at (4.6,1.75) {present perfect\\ \eng{I have seen it}};
% present continuous : trait autour de NOW
\draw[very thick, poporange] (5.2,0) -- (6.8,0);
\draw[very thick, poporange] (5.2,-0.18) -- (5.2,0.18);
\draw[very thick, poporange] (6.8,-0.18) -- (6.8,0.18);
\node[below, poporange, align=center, font=\small] at (6,-0.9) {present continuous\\ \eng{I am watching now}};
% will / going to : zone future
\draw[very thick, popgreen, dashed] (8.6,0.55) -- (11.8,0.55);
\fill[poppink] (9.4,0) circle (0.12);
\fill[popgold] (11,0) circle (0.12);
\node[above, popgreen, align=center, font=\small] at (10.2,0.75) {will / be going to\\ \eng{I will watch / I am going to watch}};
% past continuous
\draw[very thick, popblue!60] (0.6,-0.75) -- (2.6,-0.75);
\draw[very thick, popblue!60] (0.6,-0.93) -- (0.6,-0.57);
\draw[very thick, popblue!60] (2.6,-0.93) -- (2.6,-0.57);
\fill[poppink] (1.8,-0.75) circle (0.1);
\node[below, popblue!60!popdark, align=center, font=\small] at (1.6,-1.1) {past continuous + interruption\\ \eng{I was watching when the phone rang}};
\end{tikzpicture}
\legende{La frise des temps anglais : chaque temps occupe une position précise par rapport à \eng{now}.}
\end{popfigure}

\begin{aretenir}[L'idée directrice]
En français, on dit souvent « je l'ai vu » sans préciser ; l'anglais, lui, \textbf{oblige} à choisir : \eng{I saw it} (action datée, coupée du présent) ou \eng{I have seen it} (bilan, expérience, lien avec le présent). Le temps anglais porte une information que le français laisse dans le vague.
\end{aretenir}

\courssub{Le present simple : habitudes et vérités générales}

\textbf{Observation.} Lisons ces phrases tirées d'un article sur les médias :

\eng{She reads the news every morning.} « Elle lit les nouvelles chaque matin. »

\eng{The BBC broadcasts in more than forty languages.} « La BBC diffuse dans plus de quarante langues. »

\eng{Most teenagers in Douala check their phones before breakfast.} « La plupart des adolescents de Douala consultent leur téléphone avant le petit-déjeuner. »

Qu'ont ces phrases en commun ? Elles décrivent des \cle{habitudes}, des routines ou des \cle{vérités générales} — pas une action en train de se faire à l'instant.

\begin{propriete}[Present simple : forme et emplois]
\textbf{Forme} : base verbale ; à la 3e personne du singulier (\eng{he, she, it}), on ajoute \textbf{-s} (ou \textbf{-es} après -s, -sh, -ch, -o, -x : \eng{watch → watches}, \eng{go → goes}). Négation : \eng{do/does + not + base verbale}. Question : \eng{Do/Does + sujet + base verbale ?}

\textbf{Emplois} :
\begin{itemize}
\item habitudes et routines : \eng{He watches TV every evening.} « Il regarde la télé chaque soir. »
\item vérités générales et faits permanents : \eng{The sun rises in the east.} « Le soleil se lève à l'est. »
\item horaires et programmes officiels : \eng{The news starts at 8 p.m.} « Le journal commence à 20 h. »
\end{itemize}
\textbf{Marqueurs typiques} : \eng{always, usually, often, sometimes, never, every day, every morning, on Mondays, twice a week}.
\end{propriete}

\begin{center}
\begin{tabular}{|l|l|l|}
\hline
\rowcolor{popblueL} Affirmation & Négation & Question \\
\hline
I watch & I do not (don't) watch & Do I watch? \\
\hline
You watch & You don't watch & Do you watch? \\
\hline
He/She watch\textbf{es} & He \textbf{doesn't} watch & \textbf{Does} he watch? \\
\hline
We watch & We don't watch & Do we watch? \\
\hline
They watch & They don't watch & Do they watch? \\
\hline
\end{tabular}
\end{center}

Remarque essentielle : à la forme négative et interrogative, le \textbf{-s} migre sur l'auxiliaire \eng{does} ; le verbe revient à sa base : \eng{Does she watch?} (jamais \eng{Does she watches?}).

\courssub{Le present continuous : l'action en cours}

\textbf{Observation.}

\eng{They are recording a programme now.} « Ils enregistrent une émission en ce moment. »

\eng{Listen! The presenter is interviewing a guest.} « Écoute ! Le présentateur interviewe un invité. »

Ici, l'action est \cle{en cours} au moment où l'on parle : elle a commencé, elle n'est pas finie.

\begin{propriete}[Present continuous : forme et emplois]
\textbf{Forme} : \eng{be} (am / is / are) + verbe en \textbf{-ing}. \eng{I am working, she is working, they are working.}

\textbf{Emplois} :
\begin{itemize}
\item action en cours au moment de la parole : \eng{I am writing an article right now.} « Je suis en train d'écrire un article. »
\item situation temporaire : \eng{She is staying with her aunt this week.} « Elle loge chez sa tante cette semaine (temporairement). »
\item rendez-vous futur fixé (voir plus bas) : \eng{We are meeting the editor tomorrow.} « Nous rencontrons le rédacteur en chef demain. »
\end{itemize}
\textbf{Marqueurs} : \eng{now, right now, at the moment, today, this week, Look!, Listen!}
\end{propriete}

\begin{attention}[Les verbes d'état refusent le continuous]
Certains verbes expriment un \cle{état} et non une action : \eng{know, like, love, hate, want, need, believe, understand, mean, own, belong, seem, prefer}. Ils ne se mettent \textbf{jamais} au present continuous.

\eng{I know the answer.} (jamais \eng{I am knowing}) « Je connais la réponse. »

\eng{She wants a new phone.} (jamais \eng{She is wanting}) « Elle veut un nouveau téléphone. »

\eng{I believe you.} « Je te crois. » — et non \eng{I am believing you}.
\end{attention}

\courssub{Le past simple : l'action passée, terminée, datée}

\textbf{Observation.}

\eng{The reporter interviewed the minister yesterday.} « Le journaliste a interviewé le ministre hier. »

\eng{In 2019, the station launched its first podcast.} « En 2019, la station a lancé son premier podcast. »

L'action est \cle{terminée} et située dans un passé \cle{daté}, sans lien avec le présent.

\begin{propriete}[Past simple : forme et emplois]
\textbf{Forme} : verbes réguliers : base + \textbf{-ed} (\eng{watch → watched, publish → published}). Verbes irréguliers : 2e colonne de la liste (\eng{go → went, see → saw, write → wrote, speak → spoke}). Négation et question avec \eng{did} + base verbale : \eng{Did you see the report? — No, I didn't see it.}

\textbf{Emplois} : action passée achevée, avec un repère de temps passé explicite ou sous-entendu.

\textbf{Marqueurs} : \eng{yesterday, last week, last year, in 2019, two days ago, when I was young}.
\end{propriete}

\courssub{Le present perfect : le bilan présent}

\textbf{Observation.}

\eng{We have just published the article.} « Nous venons de publier l'article. »

\eng{She has worked for this newspaper for ten years.} « Elle travaille pour ce journal depuis dix ans. »

\eng{Have you ever interviewed a celebrity?} « As-tu déjà interviewé une célébrité ? »

Dans les trois cas, le passé \cle{touche le présent} : l'article vient de paraître (résultat visible maintenant), elle travaille \emph{encore} pour ce journal, et la question porte sur toute la vie jusqu'à aujourd'hui.

\begin{propriete}[Present perfect : forme et emplois]
\textbf{Forme} : \eng{have/has} + participe passé (\eng{I have seen, she has written}).

\textbf{Emplois} :
\begin{itemize}
\item bilan ou expérience sans date : \eng{I have read that book.} « J'ai lu ce livre (à un moment ou à un autre). »
\item action récente au résultat présent : \eng{He has just sent the email.} « Il vient d'envoyer le courriel. »
\item action commencée dans le passé et qui continue : \eng{They have lived in Buea since 2015.} « Ils habitent à Buea depuis 2015. »
\end{itemize}
\textbf{Marqueurs} : \eng{for, since, already, just, never, ever, yet, recently, so far}.
\end{propriete}

\begin{propriete}[La règle d'or : present perfect ou past simple ?]
On emploie le present perfect pour un bilan \textbf{sans date précise}. Dès qu'une \textbf{date passée} apparaît (\eng{yesterday, last year, in 2020, when...}), on revient obligatoirement au past simple.

\eng{I have seen the report.} « Je l'ai vu » (bilan, pas de date) — mais \eng{I saw it yesterday.} « Je l'ai vu hier. »

\eng{She has visited London twice.} « Elle a visité Londres deux fois » (dans sa vie) — mais \eng{She visited London in 2022.} « Elle a visité Londres en 2022. »
\end{propriete}

\begin{center}
\begin{tabularx}{\linewidth}{|X|X|}
\hline
\rowcolor{poppurpleL} \textbf{since + point de départ} & \textbf{for + durée} \\
\hline
\eng{since 2015} (depuis 2015) & \eng{for ten years} (depuis dix ans / pendant dix ans) \\
\hline
\eng{since Monday} (depuis lundi) & \eng{for a long time} (depuis longtemps) \\
\hline
\eng{since last January} (depuis janvier dernier) & \eng{for two weeks} (depuis deux semaines) \\
\hline
\eng{since I was a child} (depuis mon enfance) & \eng{for ages} (depuis une éternité) \\
\hline
\end{tabularx}
\end{center}

\begin{attention}[Le piège du « depuis » français]
Le français dit « j'habite ici \emph{depuis} dix ans » avec un présent. L'anglais exige le present perfect : \eng{I have lived here for ten years.} — jamais \eng{I live here for ten years}. Et attention : \eng{since} ne s'emploie jamais avec une durée (\eng{since ten years} est faux ; dites \eng{for ten years}).
\end{attention}

\courssub{Le futur : will, be going to, present continuous}

L'anglais n'a pas un futur mais \cle{trois façons} de parler de l'avenir, selon le sens :

\begin{propriete}[Choisir la bonne forme de futur]
\begin{itemize}
\item \textbf{will + base verbale} : décision spontanée, promesse, prédiction d'opinion. \eng{The phone is ringing — I'll answer it!} « Le téléphone sonne — je vais répondre ! » \eng{I think online media will replace newspapers.} « Je pense que les médias en ligne remplaceront les journaux. »
\item \textbf{be going to + base verbale} : projet, intention déjà formée, ou conséquence visible. \eng{We are going to launch a school radio next term.} « Nous allons lancer une radio scolaire le trimestre prochain. » \eng{Look at those clouds — it is going to rain.} « Regarde ces nuages — il va pleuvoir. »
\item \textbf{present continuous} : rendez-vous fixé, organisation concrète (date, lieu, personnes déjà prévus). \eng{I am meeting the journalist at 5 p.m. tomorrow.} « Je rencontre le journaliste demain à 17 h. »
\end{itemize}
\end{propriete}

\courssub{Le past continuous : l'action interrompue}

\textbf{Observation.}

\eng{While I was watching TV, the phone rang.} « Pendant que je regardais la télé, le téléphone a sonné. »

\eng{They were listening to the radio when the lights went out.} « Ils écoutaient la radio quand la lumière s'est éteinte. »

Le schéma est toujours le même : une \cle{action longue en cours} (past continuous : \eng{was/were + V-ing}) est \cle{interrompue} par une action courte (past simple). L'action longue s'introduit souvent par \eng{while} ou \eng{when} ; l'interruption par \eng{when}.

\begin{aretenir}[La combinaison à mémoriser]
\eng{While + past continuous} (le décor) \;+\; \eng{past simple} (l'événement qui coupe).

\eng{She was recording her podcast when the microphone stopped working.} « Elle enregistrait son podcast quand le micro a cessé de fonctionner. »
\end{aretenir}

\courssub{Tableau récapitulatif : temps, marqueurs et exemples}

\begin{center}
\begin{tabularx}{\linewidth}{|l|X|X|}
\hline
\rowcolor{popblueL} \textbf{Temps} & \textbf{Marqueurs typiques} & \textbf{Exemple} \\
\hline
Present simple & \eng{always, every day, usually, never} & \eng{She reads the news every morning.} \\
\hline
Present continuous & \eng{now, at the moment, Look!} & \eng{They are recording a programme now.} \\
\hline
Past simple & \eng{yesterday, last week, in 2019, ago} & \eng{He interviewed the minister yesterday.} \\
\hline
Present perfect & \eng{for, since, just, already, ever, never, yet} & \eng{We have just published the article.} \\
\hline
Past continuous & \eng{while, when (interruption)} & \eng{I was watching TV when the phone rang.} \\
\hline
Futur (will) & \eng{tomorrow, soon, I think...} & \eng{Online media will grow.} \\
\hline
Futur (going to) & \eng{next week, this evening} & \eng{We are going to launch a podcast.} \\
\hline
\end{tabularx}
\end{center}

\begin{center}
\begin{tabularx}{\linewidth}{|X|X|X|}
\hline
\rowcolor{popgreenL} \textbf{Français} & \textbf{English} & \textbf{Remarque} \\
\hline
Elle lit les nouvelles chaque matin. & \eng{She reads the news every morning.} & -s à la 3e personne \\
\hline
Ils enregistrent une émission en ce moment. & \eng{They are recording a programme now.} & action en cours \\
\hline
Nous venons de publier l'article. & \eng{We have just published the article.} & « venir de » = \eng{have just} \\
\hline
Je l'ai vu hier. & \eng{I saw it yesterday.} & date passée → past simple \\
\hline
J'habite ici depuis dix ans. & \eng{I have lived here for ten years.} & « depuis » + durée = \eng{for} \\
\hline
Je regardais la télé quand le téléphone a sonné. & \eng{I was watching TV when the phone rang.} & longue + courte \\
\hline
\end{tabularx}
\end{center}

\courssub{Entraînement}

\begin{exoresolu}[Mettre les verbes au bon temps]
Énoncé — Complétez avec le temps correct du verbe entre parenthèses :
\begin{enumerate}
\item Every evening, my father \dots\ (watch) the news on CRTV.
\item Look! The reporter \dots\ (interview) the mayor.
\item I \dots\ (never / visit) a television studio.
\item They \dots\ (launch) their website in 2021.
\item She \dots\ (work) for this radio station since 2018.
\item While we \dots\ (listen) to the debate, the power \dots\ (go) off.
\item I think print newspapers \dots\ (disappear) one day.
\end{enumerate}
\tcblower
Solution détaillée :
\begin{enumerate}
\item \eng{watches} — habitude (\eng{every evening}) → present simple, 3e personne : -es après -ch.
\item \eng{is interviewing} — \eng{Look!} signale une action en cours → present continuous.
\item \eng{have never visited} — expérience de vie sans date → present perfect.
\item \eng{launched} — date passée (\eng{in 2021}) → past simple obligatoire.
\item \eng{has worked} — \eng{since 2018}, action qui continue → present perfect.
\item \eng{were listening / went} — action longue interrompue : past continuous + past simple.
\item \eng{will disappear} — prédiction d'opinion (\eng{I think}) → \eng{will}.
\end{enumerate}
\end{exoresolu}

\begin{attention}[Bilan des erreurs fréquentes des francophones]
\begin{itemize}
\item Oublier le \textbf{-s} de la 3e personne : \eng{he watch} au lieu de \eng{he watches}.
\item Employer le present continuous avec un verbe d'état : \eng{I am knowing} au lieu de \eng{I know}.
\item Mettre une date passée avec le present perfect : \eng{I have seen him yesterday} au lieu de \eng{I saw him yesterday}.
\item Traduire « depuis » par \eng{since} devant une durée : \eng{since ten years} au lieu de \eng{for ten years}.
\item Oublier que \eng{did} absorbe le passé : \eng{Did you went?} au lieu de \eng{Did you go?}
\end{itemize}
\end{attention}

\begin{savaistu}[Le present perfect dans les titres de presse]
Au Nigeria et au Ghana, la presse anglophone utilise massivement le present perfect dans les titres de dépêches : \eng{Government has announced new media regulations} « Le gouvernement a annoncé de nouvelles règles pour les médias ». Pourquoi ce choix ? Parce que la nouvelle vient d'arriver et que ses conséquences sont encore d'actualité : c'est exactement la logique du present perfect, un bilan frais relié au présent. Quand le journal revient sur l'événement le lendemain, il passe au past simple : \eng{The government announced the regulations on Tuesday.}
\end{savaistu}

\begin{exercice}[À vous de jouer]
Mettez les verbes entre parenthèses au temps qui convient, puis traduisez les phrases 2 et 5 en français :
\begin{enumerate}
\item My sister usually \dots\ (read) the newspaper online.
\item We \dots\ (just / receive) an email from the editor.
\item The journalist \dots\ (arrive) in Bamenda two days ago.
\item Listen! Someone \dots\ (speak) on the radio.
\item They \dots\ (know) each other for many years.
\item While I \dots\ (write) my article, my computer \dots\ (crash).
\item We \dots\ (meet) the cameraman at 9 a.m. tomorrow — it is already arranged.
\end{enumerate}
\end{exercice}

\begin{corrige}[Exercice « À vous de jouer »]
\begin{enumerate}
\item \eng{reads} — habitude (\eng{usually}) → present simple, 3e personne.
\item \eng{have just received} — action récente au résultat présent → present perfect.
\item \eng{arrived} — date passée (\eng{two days ago}) → past simple.
\item \eng{is speaking} — \eng{Listen!} → action en cours, present continuous.
\item \eng{have known} — \eng{for many years} + verbe d'état → present perfect simple (jamais \eng{have been knowing}).
\item \eng{was writing / crashed} — action longue interrompue : past continuous + past simple.
\item \eng{are meeting} — rendez-vous fixé et organisé → present continuous à valeur de futur.
\end{enumerate}
Traductions : (2) « Nous venons de recevoir un courriel du rédacteur en chef. » (5) « Ils se connaissent depuis de nombreuses années. »
\end{corrige}

\coursec{Voix passive et discours rapporté}

Dans la section précédente, nous avons révisé les temps verbaux anglais et vu comment chaque temps exprime un point de vue sur l'action. Mais dans les médias, le journaliste ne dit pas toujours \emph{qui} fait l'action : il met souvent en avant le \emph{résultat}. Il ne dit pas non plus toujours les paroles mot pour mot : il les \emph{rapporte}. Ces deux outils — la \cle{voix passive} et le \cle{discours rapporté} — sont les deux piliers grammaticaux de l'anglais de la presse, et deux questions quasi systématiques au Baccalauréat.

\courssub{Observation : repérer les structures dans un texte de presse}

\begin{textebox}[Reading text — “Local Radio Station Reopens in Bamenda”]
Good news for media lovers in the North-West Region. The community radio station “Highlands FM”, which was closed two years ago, was officially reopened on Monday by the regional authorities. The building had been damaged during a storm, and most of the equipment had been stolen. New transmitters have now been installed, and programmes will be broadcast in English, Pidgin and several local languages.

Speaking at the ceremony, the station manager said that the staff were delighted to be back on air. He explained that the repairs had been financed by local businesses and by listeners abroad. “We will train young journalists here,” he announced, “and the first courses will start next month.” He added that the station had always been considered the voice of the community.

A media expert interviewed after the ceremony declared that community radios were being given a second chance across the country. She said that many stations had been closed in recent years, but that new licences would soon be granted. “Information is power,” she concluded, “and local voices must be heard.”
\end{textebox}

\textbf{Glossaire.} \emph{to reopen} (v.) : rouvrir ; \emph{equipment} (n. sing.) : équipement, matériel ; \emph{transmitter} (n.) : émetteur ; \emph{to broadcast} (v.) : diffuser ; \emph{staff} (n. sing.) : personnel ; \emph{licence} (n.) : licence, autorisation ; \emph{to grant} (v.) : accorder.

\textbf{Questions d'observation.}
\begin{enumerate}
\item Soulignez tous les verbes construits avec \eng{be} + participe passé. Qui fait l'action dans chaque cas ? Est-ce précisé ?
\item Relevez les phrases où le journaliste rapporte les paroles du directeur. Que devient \eng{will} ? Que devient \eng{next month} ?
\end{enumerate}

On constate que le texte est plein de formes comme \eng{was reopened}, \eng{had been stolen}, \eng{will be broadcast} : c'est la \cle{voix passive}. Et les paroles du directeur ne sont presque jamais citées directement : elles sont \emph{rapportées} avec \eng{said that}, \eng{explained that}, \eng{announced} : c'est le \cle{discours rapporté}.

\courssub{La voix passive : formation}

\begin{propriete}[Formation de la voix passive]
La voix passive se construit ainsi :
\begin{center}
\textbf{Sujet + auxiliaire \eng{be} (au temps voulu) + participe passé (+ \eng{by} + agent)}
\end{center}
C'est \eng{be} qui porte la marque du temps, de la négation et de l'interrogation ; le participe passé ne change jamais.
\end{propriete}

\begin{center}
\begin{tabular}{|l|l|l|}
\hline
\rowcolor{popblueL} \textbf{Temps} & \textbf{Actif} & \textbf{Passif} \\
\hline
Present simple & Journalists write articles. & Articles \textbf{are written}. \\
\hline
Present continuous & They are repairing the studio. & The studio \textbf{is being repaired}. \\
\hline
Past simple & The storm damaged the building. & The building \textbf{was damaged}. \\
\hline
Past perfect & Thieves had stolen the equipment. & The equipment \textbf{had been stolen}. \\
\hline
Future (\eng{will}) & They will broadcast the news. & The news \textbf{will be broadcast}. \\
\hline
Modal (\eng{must}) & We must hear local voices. & Local voices \textbf{must be heard}. \\
\hline
\end{tabular}
\end{center}

\begin{propriete}[Emplois de la voix passive]
On emploie le passif quand :
\begin{enumerate}
\item l'\textbf{agent est inconnu} : \eng{My bicycle has been stolen.} (On a volé mon vélo — on ne sait pas qui.)
\item l'agent est \textbf{sans importance ou évident} : \eng{The suspect was arrested.} (Le suspect a été arrêté — évidemment par la police.)
\item on veut \textbf{mettre en avant le résultat} ou le fait, très fréquent dans la presse et les textes officiels : \eng{The law was voted yesterday.} (La loi a été votée hier.)
\end{enumerate}
L'agent, quand il est mentionné, est introduit par \eng{by} : \eng{The station was reopened by the authorities.}
\end{propriete}

\begin{attention}[Passif et français]
Le français utilise souvent le pronom \emph{on} ou une tournure pronominale là où l'anglais exige le passif : \emph{On a volé mon vélo} = \eng{My bicycle has been stolen} (jamais \eng{they have stolen my bicycle} si le voleur est inconnu). De même, \emph{cela se dit} = \eng{it is said}. Attention aussi : seuls les verbes \textbf{transitifs} (suivis d'un complément d'objet direct) peuvent se mettre au passif. On ne peut pas dire \eng{was arrived} : \eng{arrive} est intransitif.
\end{attention}

\courssub{Transformation actif $\rightarrow$ passif}

\begin{methode}[Passer de la voix active à la voix passive]
\begin{enumerate}
\item Repérer le \textbf{complément d'objet} du verbe actif : il devient le \textbf{sujet} de la phrase passive.
\item Mettre \eng{be} au \textbf{même temps} que le verbe actif, en l'accordant avec le nouveau sujet.
\item Ajouter le \textbf{participe passé} du verbe.
\item Ajouter éventuellement \eng{by} + l'ancien sujet (l'agent), souvent omis.
\end{enumerate}
\end{methode}

\begin{popfigure}
\begin{tikzpicture}[node distance=0.6cm, >=Stealth]
\node[draw=popblue, fill=popblueL, rounded corners, align=center, font=\small] (active) {\textbf{Active}\\ Journalists \underline{write} \underline{articles}.\\ \tiny sujet \quad verbe \quad objet};
\node[draw=poporange, fill=poporangeL, rounded corners, align=center, font=\small, right=4.6cm of active] (passive) {\textbf{Passive}\\ \underline{Articles} \underline{are written} (by journalists).\\ \tiny sujet \quad be + participe passé \quad agent};
\draw[->, very thick, popdark] (active) -- node[above, align=center, font=\footnotesize]{objet $\rightarrow$ sujet\\ verbe $\rightarrow$ be + P.P.} (passive);
\end{tikzpicture}
\legende{Schéma de transformation de la voix active à la voix passive.}
\end{popfigure}

\begin{exemplebox}[Exemples traduits]
\eng{The interview was broadcast yesterday.} « L'interview a été diffusée hier. »

\eng{The minister announced the decision.} $\rightarrow$ \eng{The decision was announced by the minister.} « La décision a été annoncée par le ministre. »

\eng{People speak English in Cameroon.} $\rightarrow$ \eng{English is spoken in Cameroon.} « On parle anglais au Cameroun. »
\end{exemplebox}

\courssub{Le discours rapporté : la concordance des temps}

Quand un journaliste rapporte des paroles, il passe du \emph{discours direct} (paroles exactes, entre guillemets) au \emph{discours indirect} (paroles intégrées à sa phrase).

\begin{propriete}[Règle de la concordance des temps]
Si le verbe introducteur est au \textbf{passé} (\eng{said}, \eng{told}, \eng{asked}), le temps de la parole rapportée \textbf{recule d'un cran} vers le passé :
\eng{“I am tired,” he said.} $\rightarrow$ \eng{He said (that) he was tired.} « Il a dit qu'il était fatigué. »
Si le verbe introducteur est au présent (\eng{he says}), il n'y a \textbf{aucun changement} de temps.
\end{propriete}

\begin{center}
\begin{tabular}{|l|l|l|}
\hline
\rowcolor{popblueL} \textbf{Discours direct} & \textbf{Discours rapporté} & \textbf{Exemple} \\
\hline
Present simple & Past simple & \eng{“I work here.”} $\rightarrow$ He said he \textbf{worked} there. \\
\hline
Present continuous & Past continuous & \eng{“I am listening.”} $\rightarrow$ She said she \textbf{was listening}. \\
\hline
Past simple & Past perfect & \eng{“I saw the film.”} $\rightarrow$ He said he \textbf{had seen} the film. \\
\hline
Present perfect & Past perfect & \eng{“I have finished.”} $\rightarrow$ She said she \textbf{had finished}. \\
\hline
\eng{will} & \eng{would} & \eng{“I will come.”} $\rightarrow$ He said he \textbf{would come}. \\
\hline
\eng{can} & \eng{could} & \eng{“I can help.”} $\rightarrow$ She said she \textbf{could help}. \\
\hline
\eng{must} & \eng{had to} & \eng{“I must go.”} $\rightarrow$ He said he \textbf{had to go}. \\
\hline
\end{tabular}
\end{center}

\textbf{Exception.} Si la parole rapportée exprime une vérité générale ou un fait toujours vrai, le recul est facultatif : \eng{The teacher said that water boils at 100 degrees.} (Le professeur a dit que l'eau bout à 100 degrés.)

\courssub{Pronoms, possessifs et marqueurs de temps et de lieu}

Le passage au discours rapporté entraîne aussi l'adaptation des \cle{pronoms} (je $\rightarrow$ il/elle selon le contexte) et des \cle{marqueurs spatio-temporels}, car on ne parle plus au même moment ni au même endroit.

\begin{center}
\begin{tabularx}{\linewidth}{|X|X|X|}
\hline
\rowcolor{popblueL} \textbf{Discours direct} & \textbf{Discours rapporté} & \textbf{Français} \\
\hline
now & then & maintenant $\rightarrow$ alors \\
\hline
today & that day & aujourd'hui $\rightarrow$ ce jour-là \\
\hline
yesterday & the day before & hier $\rightarrow$ la veille \\
\hline
tomorrow & the next day & demain $\rightarrow$ le lendemain \\
\hline
here & there & ici $\rightarrow$ là \\
\hline
this / these & that / those & ceci $\rightarrow$ cela \\
\hline
ago & before & il y a (durée) $\rightarrow$ avant \\
\hline
\end{tabularx}
\end{center}

\begin{exemplebox}[Exemples traduits]
\eng{“We will publish the story tomorrow,” the editor said.} $\rightarrow$ \eng{The editor said (that) they would publish the story the next day.} « Le rédacteur a dit qu'ils publieraient l'article le lendemain. »

\eng{The minister announced that the law would change.} « Le ministre a annoncé que la loi changerait. »

\eng{“I interviewed the mayor here yesterday,” she said.} $\rightarrow$ \eng{She said she had interviewed the mayor there the day before.} « Elle a dit qu'elle avait interviewé le maire là-bas la veille. »
\end{exemplebox}

\courssub{Les verbes introducteurs : say, tell, ask, explain, announce}

\begin{propriete}[Say ou tell ?]
\begin{itemize}
\item \eng{say (that) + phrase} : \emph{sans} complément de personne. \eng{He said that he was tired.}
\item \eng{tell + someone (that) + phrase} : \emph{avec} complément de personne obligatoire. \eng{He told \textbf{me} that he was tired.}
\item \eng{ask + someone + if/whether...} pour les questions fermées ; \eng{ask + someone + wh-...} pour les questions ouvertes : \eng{She asked me if I liked radio. / She asked me where I lived.}
\item \eng{explain (that)}, \eng{announce (that)}, \eng{add (that)}, \eng{declare (that)} : fréquents dans la presse.
\end{itemize}
\end{propriete}

\begin{methode}[Comment rapporter une parole en quatre étapes]
\begin{enumerate}
\item \textbf{Choisir} le verbe introducteur : \eng{say}, \eng{tell} + personne, ou \eng{ask} pour une question.
\item \textbf{Reculer} le temps du verbe rapporté d'un cran (present $\rightarrow$ past, past $\rightarrow$ past perfect, \eng{will} $\rightarrow$ \eng{would}, \eng{can} $\rightarrow$ \eng{could}).
\item \textbf{Adapter} les pronoms, les possessifs et les marqueurs temporels (\eng{now} $\rightarrow$ \eng{then}, \eng{today} $\rightarrow$ \eng{that day}, \eng{tomorrow} $\rightarrow$ \eng{the next day}).
\item \textbf{Supprimer} les guillemets et les deux-points ; relier avec \eng{that} (facultatif) ou \eng{if/whether} pour une question fermée.
\end{enumerate}
\end{methode}

\begin{exoresolu}[Transformation au discours rapporté et au passif]
Énoncé. Mettez à la forme demandée :
\begin{enumerate}
\item \eng{“I will broadcast the news tonight,” the presenter said.} (discours rapporté)
\item \eng{The police arrested the journalist.} (voix passive)
\end{enumerate}
\tcblower
Solution détaillée.
\begin{enumerate}
\item Verbe introducteur au passé (\eng{said}) : on recule \eng{will} $\rightarrow$ \eng{would}, on adapte le pronom \eng{I} $\rightarrow$ \eng{she/he} et le marqueur \eng{tonight} $\rightarrow$ \eng{that night}. On obtient : \eng{The presenter said (that) she would broadcast the news that night.} « La présentatrice a dit qu'elle diffuserait les informations ce soir-là. »
\item L'objet \eng{the journalist} devient sujet ; \eng{arrested} (past simple) devient \eng{was} + participe passé : \eng{The journalist was arrested by the police.} « Le journaliste a été arrêté par la police. »
\end{enumerate}
\end{exoresolu}

\begin{attention}[Erreurs fréquentes des francophones]
\begin{itemize}
\item Ne dites jamais \eng{it was sayed} : \eng{say} est irrégulier — \textbf{say – said – said} (« sé – sèd – sèd »). De même \eng{tell – told – told}.
\item Ne confondez pas \eng{say} et \eng{tell} : on dit \eng{He said that...} mais \eng{He told \textbf{me} that...} — jamais \eng{he said me} ni \eng{he told that}.
\item N'oubliez pas le recul des temps : \eng{He said he \textbf{will} come} est fautif ; il faut \eng{He said he \textbf{would} come}.
\item Au passif, n'oubliez jamais l'auxiliaire \eng{be} : \eng{the news broadcast yesterday} est incorrect ; il faut \eng{the news \textbf{was} broadcast yesterday}.
\item Attention aux participes passés irréguliers fréquents dans les médias : \eng{broadcast – broadcast – broadcast}, \eng{write – wrote – written}, \eng{steal – stole – stolen}.
\end{itemize}
\end{attention}

\begin{exercice}[À vous de jouer]
\begin{enumerate}
\item Mettez à la voix passive : \eng{The students read the newspaper every morning.} / \eng{Someone has stolen my phone.}
\item Rapportez : \eng{“We are preparing a new programme,” the producer told us.} / \eng{“Did you watch the debate yesterday?” she asked me.}
\end{enumerate}
\end{exercice}

\begin{corrige}[Exercice « À vous de jouer »]
\begin{enumerate}
\item \eng{The newspaper is read by the students every morning.} (Le journal est lu par les élèves chaque matin.) / \eng{My phone has been stolen.} (On a volé mon téléphone — agent inconnu, donc pas de \eng{by}.)
\item \eng{The producer told us (that) they were preparing a new programme.} (Le producteur nous a dit qu'ils préparaient une nouvelle émission.) / \eng{She asked me if I had watched the debate the day before.} (Elle m'a demandé si j'avais regardé le débat la veille.)
\end{enumerate}
\end{corrige}

\begin{aretenir}[L'essentiel de la section]
Passif = \eng{be} au temps voulu + participe passé ; l'objet devient sujet, l'agent est introduit par \eng{by} ou omis. Discours rapporté : verbe introducteur au passé $\Rightarrow$ recul des temps (present $\rightarrow$ past, past $\rightarrow$ past perfect, \eng{will} $\rightarrow$ \eng{would}, \eng{can} $\rightarrow$ \eng{could}), adaptation des pronoms et des marqueurs (\eng{now} $\rightarrow$ \eng{then}, \eng{tomorrow} $\rightarrow$ \eng{the next day}), suppression des guillemets. \eng{Say} sans personne, \eng{tell} + personne.
\end{aretenir}

\coursec{Modaux, relatifs et comparaison}

Dans la section précédente, nous avons révisé la voix passive (\eng{The article was published}) et le discours rapporté (\eng{She said that the news had been verified}). Ces deux structures sont omniprésentes dans les médias. Mais un article de presse bien écrit mobilise aussi trois autres outils grammaticaux fondamentaux : les \cle{verbes modaux} pour exprimer l'obligation, le conseil ou la probabilité (\eng{Journalists must respect ethics}), les \cle{pronoms relatifs} pour relier les idées (\eng{the reporter who interviewed me}), et les \cle{comparatifs et superlatifs} pour évaluer et hiérarchiser l'information (\eng{the most reliable source}). C'est l'objet de cette quatrième section de révision générale.

\courssub{Les verbes modaux : observation}

\begin{textebox}[Editorial — “Responsible journalism”]
Every journalist must check the facts before publishing. Reporters should protect their sources, and they mustn't invent quotations. Readers, for their part, should not believe everything they see on social media: a viral post may be true, but it might also be completely false. In Cameroon, anyone can create a blog today, so we could soon face an ocean of unverified information. Fortunately, you don't have to be a professional to spot fake news — you only need to ask simple questions: who wrote this? where was it published? can the facts be confirmed elsewhere?
\end{textebox}

\textbf{Glossaire :} \eng{to check} (vérifier), \eng{a source} (une source), \eng{a quotation} (une citation), \eng{viral} (viral), \eng{to spot} (repérer), \eng{elsewhere} (ailleurs).

Observez les verbes en gras implicite du texte : \eng{must check}, \eng{should protect}, \eng{mustn't invent}, \eng{may be}, \eng{might be}, \eng{can create}, \eng{could face}, \eng{don't have to be}. Tous ces verbes sont des \cle{modaux} (ou semi-auxiliaires) : ils n'expriment pas une action mais l'\emph{attitude} du locuteur envers l'action (obligation, conseil, capacité, probabilité, interdiction).

\begin{propriete}[Forme des modaux : trois règles d'or]
\begin{enumerate}
\item Le modal est \textbf{invariable} : il ne prend jamais de \textbf{-s} à la troisième personne. \eng{He must check the facts.} (et non \emph{he musts}).
\item Il est suivi de la \textbf{base verbale sans to} : \eng{You should check the source.} (et non \emph{you should to check}).
\item La négation et l'interrogation se font \textbf{sans do} : \eng{You mustn't share this.} / \eng{Can you confirm the story?}
\end{enumerate}
Exception de forme : \eng{have to} n'est pas un vrai modal — il se conjugue (\eng{has to}, \eng{had to}) et prend l'auxiliaire \eng{do} : \eng{Do you have to pay? / She doesn't have to pay.}
\end{propriete}

\begin{attention}[L'erreur classique du francophone]
En français on dit « il doit vérifier » avec l'infinitif ; la tentation est grande de traduire par \emph{he must to check}. Faux ! Le modal absorbe déjà la fonction de l'infinitif : \eng{he must check}. De même, \emph{he cans speak} est impossible : \eng{he can speak}.
\end{attention}

\begin{popfigure}
\begin{center}
\begin{tikzpicture}[
  head/.style={fill=popblue, text=white, font=\bfseries, minimum height=0.8cm, align=center},
  cell/.style={fill=popblueL, minimum height=0.8cm, align=center, font=\small},
  ex/.style={fill=popcream, minimum height=0.8cm, align=center, font=\small\itshape}]
\node[head] (h0) at (0,0) {Modal};
\node[head, right=0.05cm of h0, minimum width=3.2cm] (h1) {Valeur};
\node[head, right=0.05cm of h1, minimum width=6.2cm] (h2) {Exemple};
\node[cell, below=0.05cm of h0, minimum width=2.2cm] (c1) {can / could};
\node[cell, right=0.05cm of c1, minimum width=3.2cm] (c1b) {capacité, permission};
\node[ex, right=0.05cm of c1b, minimum width=6.2cm] (c1c) {Anyone can create a blog.};
\node[cell, below=0.05cm of c1, minimum width=2.2cm] (c2) {must / have to};
\node[cell, right=0.05cm of c2, minimum width=3.2cm] (c2b) {obligation};
\node[ex, right=0.05cm of c2b, minimum width=6.2cm] (c2c) {Journalists must respect ethics.};
\node[cell, below=0.05cm of c2, minimum width=2.2cm] (c3) {should};
\node[cell, right=0.05cm of c3, minimum width=3.2cm] (c3b) {conseil};
\node[ex, right=0.05cm of c3b, minimum width=6.2cm] (c3c) {You should check the source.};
\node[cell, below=0.05cm of c3, minimum width=2.2cm] (c4) {may / might};
\node[cell, right=0.05cm of c4, minimum width=3.2cm] (c4b) {probabilité};
\node[ex, right=0.05cm of c4b, minimum width=6.2cm] (c4c) {This news might be false.};
\end{tikzpicture}
\legende{Tableau à double entrée : modaux, valeurs et exemples dans le contexte des médias.}
\end{center}
\end{popfigure}

\begin{exemplebox}[Exemples traduits]
\begin{itemize}
\item \eng{You should check the source before sharing.} « Tu devrais vérifier la source avant de partager. » (conseil)
\item \eng{Journalists must respect ethics.} « Les journalistes doivent respecter la déontologie. » (obligation morale, venant du locuteur)
\item \eng{This news might be false.} « Cette nouvelle pourrait être fausse. » (probabilité incertaine)
\item \eng{She can speak three languages, so she could work for the BBC.} « Elle sait parler trois langues, donc elle pourrait travailler pour la BBC. » (capacité au présent, puis hypothèse)
\item \eng{Citizens have to pay for a TV licence in some countries.} « Dans certains pays, les citoyens doivent payer une redevance télé. » (obligation extérieure, règlement)
\end{itemize}
\end{exemplebox}

\begin{propriete}[Nuances d'emploi]
\begin{itemize}
\item \textbf{Capacité} : \eng{can} (présent), \eng{could} (passé ou hypothèse). \eng{I could read English when I was ten.}
\item \textbf{Permission} : \eng{can} (familier), \eng{may} (formel). \eng{May I ask a question, Sir?}
\item \textbf{Obligation} : \eng{must} = obligation ressentie ou imposée par le locuteur ; \eng{have to} = obligation extérieure (loi, règlement). \eng{I must finish this article} (je me l'impose) ; \eng{Students have to wear a uniform} (règlement de l'école).
\item \textbf{Conseil} : \eng{should} / \eng{shouldn't}. \eng{You shouldn't believe rumours.}
\item \textbf{Probabilité} : \eng{may} (possible), \eng{might} (encore plus incertain), \eng{must} (certitude logique : « doit sûrement »). \eng{He must be tired — he worked all night.} « Il doit être fatigué. »
\end{itemize}
\end{propriete}

\begin{attention}[mustn't $\neq$ don't have to — le piège capital]
C'est l'une des erreurs les plus graves au Bac, car le sens change totalement :
\begin{itemize}
\item \eng{You mustn't share fake news.} = « Il est \textbf{interdit} de partager les fausses informations. » (ne le fais pas !)
\item \eng{You don't have to buy the paper.} = « Tu n'es \textbf{pas obligé} d'acheter le journal. » (tu peux, mais ce n'est pas nécessaire)
\end{itemize}
En français, « il ne faut pas » correspond à \eng{mustn't}, tandis que « ce n'est pas la peine de / on n'est pas obligé de » correspond à \eng{don't have to}. Ne traduisez jamais « tu n'as pas à le faire » par \emph{you mustn't do it}.
\end{attention}

\begin{exoresolu}[Complétez avec le modal approprié]
Énoncé — Complétez avec \eng{must}, \eng{mustn't}, \eng{don't have to}, \eng{should} ou \eng{might} : (1) You \dots\ share this video; it contains hate speech. (2) You \dots\ subscribe to the channel, but it helps the creator. (3) A good reporter \dots\ always verify his facts. (4) Be careful: that website \dots\ not be reliable. (5) You look exhausted; you \dots\ take a break.
\tcblower
Solution : (1) \eng{mustn't} — interdiction absolue (discours haineux). (2) \eng{don't have to} — absence d'obligation : c'est facultatif. (3) \eng{must} — obligation professionnelle forte. (4) \eng{might} — probabilité incertaine (« pourrait ne pas être fiable »). (5) \eng{should} — conseil amical. Phrases complètes : \eng{You mustn't share this video; it contains hate speech. / You don't have to subscribe to the channel, but it helps the creator. / A good reporter must always verify his facts. / Be careful: that website might not be reliable. / You look exhausted; you should take a break.}
\end{exoresolu}

\courssub{Les pronoms relatifs : who, which, whose, where, that}

Observez ces deux phrases tirées de notre éditorial : \eng{The reporter \textbf{who} interviewed me works for CRTV.} et \eng{The article \textbf{which} I read was well documented.} Le mot en gras relie deux propositions : c'est un \cle{pronom relatif}. Il évite la répétition et rend la phrase plus fluide — exactement comme « qui », « que », « dont », « où » en français.

\begin{propriete}[Règle des relatifs]
\begin{itemize}
\item \eng{who} = pour les \textbf{personnes} (sujet) : \eng{The reporter who interviewed me.} « Le journaliste qui m'a interviewé. »
\item \eng{which} = pour les \textbf{choses et animaux} : \eng{The article which I read.} « L'article que j'ai lu. »
\item \eng{whose} = \textbf{possession} (« dont », « duquel ») : \eng{The journalist whose camera was stolen.} « Le journaliste dont la caméra a été volée. »
\item \eng{where} = \textbf{lieux} (« où ») : \eng{The studio where the news is recorded.} « Le studio où le journal est enregistré. »
\item \eng{that} = peut remplacer \eng{who} ou \eng{which} dans les propositions \textbf{définissantes} : \eng{The news that shocked the country.} « La nouvelle qui a choqué le pays. »
\end{itemize}
Attention : après une \textbf{virgule} (proposition explicative), on utilise obligatoirement \eng{who} ou \eng{which}, jamais \eng{that} : \eng{My uncle, who lives in Buea, is a cameraman.}
\end{propriete}

\begin{center}
\begin{tabularx}{\linewidth}{|X|X|X|}
\hline
\rowcolor{popblueL} Français & English & Remarque \\
\hline
le journaliste qui m'a interviewé & the reporter who interviewed me & personne : who \\
\hline
l'article que j'ai lu & the article (which) I read & chose : which, omissible si complément \\
\hline
la ville où je suis né & the town where I was born & lieu : where \\
\hline
le journaliste dont la caméra & the journalist whose camera & possession : whose \\
\hline
la nouvelle qui a choqué le pays & the news that shocked the country & that possible (définissant) \\
\hline
\end{tabularx}
\end{center}

\begin{attention}[Pièges des francophones]
\begin{enumerate}
\item Ne doublez pas le complément : \emph{The article which I read it} est faux — \eng{which} remplace déjà \eng{it}. On écrit \eng{The article which I read.}
\item \eng{whose} ne s'emploie pas seulement pour les personnes : \eng{a country whose economy is growing} est correct.
\item Quand le relatif est \textbf{complément}, on peut l'omettre : \eng{The article I read was interesting.} Mais jamais quand il est \textbf{sujet} : \emph{The reporter interviewed me works here} est impossible.
\end{enumerate}
\end{attention}

\courssub{Comparatifs et superlatifs}

Les médias comparent sans cesse : audiences, prix, fiabilité. Observez : \eng{Online news is \textbf{faster than} printed newspapers, but radio is still \textbf{the most popular} medium in many villages.}

\begin{propriete}[Formation]
\begin{itemize}
\item \textbf{Adjectifs courts} (1 syllabe, ou 2 syllabes en -y) : comparatif en \textbf{-er + than}, superlatif en \textbf{the -est}. \eng{fast $\rightarrow$ faster than $\rightarrow$ the fastest} ; \eng{easy $\rightarrow$ easier $\rightarrow$ the easiest}.
\item \textbf{Adjectifs longs} (2 syllabes et plus) : \textbf{more ... than}, \textbf{the most ...}. \eng{reliable $\rightarrow$ more reliable than $\rightarrow$ the most reliable}.
\item \textbf{Irréguliers} : \eng{good $\rightarrow$ better $\rightarrow$ the best} ; \eng{bad $\rightarrow$ worse $\rightarrow$ the worst} ; \eng{far $\rightarrow$ further $\rightarrow$ the furthest}.
\item Orthographe : consonne finale doublée après voyelle brève (\eng{big $\rightarrow$ bigger $\rightarrow$ the biggest}) ; \eng{-y} devient \eng{-i} (\eng{happy $\rightarrow$ happier}).
\end{itemize}
\end{propriete}

\begin{propriete}[Structures de comparaison]
\begin{itemize}
\item \textbf{Égalité} : \eng{as ... as}. \eng{Radio is as important as television in rural areas.} « La radio est aussi importante que la télévision en zone rurale. » Négation : \eng{not as ... as} / \eng{not so ... as}.
\item \textbf{Infériorité} : \eng{less ... than}. \eng{Blogs are less reliable than newspapers.} « Les blogs sont moins fiables que les journaux. »
\item \textbf{Proportion} : \eng{the more ..., the more ...}. \eng{The more you read, the more vocabulary you learn.} « Plus tu lis, plus tu apprends de vocabulaire. »
\end{itemize}
\end{propriete}

\begin{attention}[Jamais de double marque de comparaison]
L'anglais ne cumule jamais les deux procédés. \emph{more better} est \textbf{faux} : on dit simplement \eng{better}. De même, \emph{the most easiest} est \textbf{faux} : on dit \eng{the easiest}. Retenez : soit la marque en \eng{-er / -est}, soit \eng{more / most}, jamais les deux ensemble. Autre piège : après un superlatif, on utilise \eng{in} (et non \eng{of}) pour le groupe : \eng{the best student in the class}.
\end{attention}

\begin{savaistu}[La structure qui impressionne le correcteur]
La construction \eng{the more ..., the more ...} est très appréciée à l'écrit du Bac car elle montre une maîtrise avancée de la syntaxe. Exemple à réutiliser dans vos compositions sur les médias : \eng{The more you read English newspapers, the more vocabulary you learn.} « Plus tu lis les journaux anglais, plus tu apprends de vocabulaire. » Variante : \eng{The faster news travels, the harder it is to verify.} « Plus l'information circule vite, plus elle est difficile à vérifier. »
\end{savaistu}

\begin{exoresolu}[Comparatifs, superlatifs et relatifs — entraînement type Bac]
Énoncé — (a) Mettez l'adjectif entre parenthèses à la forme correcte : (1) The internet is \dots\ (fast) than the post office. (2) This is \dots\ (bad) article I have ever read. (3) Television is \dots\ (expensive) medium than radio. (b) Reliez les phrases avec le relatif approprié : (4) The journalist won a prize. He reported on the cocoa crisis. (5) The newspaper was closed. It published fake news.
\tcblower
Solution : (1) \eng{faster} — adjectif court, comparatif en -er : \eng{The internet is faster than the post office.} (2) \eng{the worst} — irrégulier : \eng{bad $\rightarrow$ worse $\rightarrow$ the worst} ; \eng{This is the worst article I have ever read.} (3) \eng{more expensive} — adjectif long : \eng{Television is a more expensive medium than radio.} (4) Personne, sujet : \eng{The journalist who reported on the cocoa crisis won a prize.} (5) Chose : \eng{The newspaper which (that) published fake news was closed.}
\end{exoresolu}

\begin{exercice}[À vous de jouer]
\begin{enumerate}
\item Complétez avec un modal justifié : (a) You \dots\ verify information before reposting it (conseil). (b) Passengers \dots\ smoke in a bendskin (interdiction). (c) You \dots\ buy a TV to watch the news; a phone is enough (absence d'obligation).
\item Reliez avec \eng{who}, \eng{which}, \eng{whose} ou \eng{where} : (a) The reporter \dots\ microphone broke continued the interview. (b) The market \dots\ I buy my newspapers is near the school.
\item Corrigez les erreurs : (a) \emph{This newspaper is more better than that one.} (b) \emph{She must to check her sources.} (c) \emph{The article which I read it was false.}
\item Traduisez : « Plus tu regardes les informations en anglais, plus ta compréhension orale s'améliore. »
\end{enumerate}
\end{exercice}

\begin{corrige}[Exercice « À vous de jouer »]
1. (a) \eng{should} (conseil) ; (b) \eng{mustn't} (interdiction) ; (c) \eng{don't have to} (facultatif). 2. (a) \eng{whose} (possession : son micro) ; (b) \eng{where} (lieu). 3. (a) \eng{This newspaper is better than that one.} (pas de double marque) ; (b) \eng{She must check her sources.} (base verbale sans to) ; (c) \eng{The article which I read was false.} (pas de double complément). 4. \eng{The more you watch the news in English, the more your listening comprehension improves.}
\end{corrige}

\begin{aretenir}[L'essentiel de la section]
\begin{itemize}
\item Modal invariable + base verbale sans \eng{to} : \eng{he must check}.
\item \eng{mustn't} = interdiction ; \eng{don't have to} = absence d'obligation.
\item Relatifs : \eng{who} (personnes), \eng{which} (choses), \eng{whose} (possession), \eng{where} (lieux), \eng{that} (définissant).
\item Comparatif : \eng{-er than} / \eng{more ... than} ; superlatif : \eng{the -est} / \eng{the most} ; jamais les deux à la fois.
\item Irréguliers : \eng{good $\rightarrow$ better $\rightarrow$ the best} ; \eng{bad $\rightarrow$ worse $\rightarrow$ the worst} ; \eng{far $\rightarrow$ further $\rightarrow$ the furthest}.
\item Structures : \eng{as ... as}, \eng{less ... than}, \eng{the more ..., the more ...}.
\end{itemize}
\end{aretenir}

\coursec{Comparaison français/anglais et pièges des francophones}

Dans les sections précédentes, nous avons révisé les temps verbaux, la voix passive, le discours rapporté, les modaux, les propositions relatives et la comparaison. Nous abordons maintenant le problème le plus redoutable pour un élève francophone : les \cle{calques}, c'est-à-dire les phrases anglaises construites mot à mot sur le modèle du français. Ce sont elles qui coûtent le plus de points au baccalauréat. Cette section vous donne une méthode de vérification systématique pour les éliminer.

\courssub{Observer avant de corriger : un paragraphe « truffé » de calques}

Lisez ce court paragraphe écrit par un candidat francophone. Il contient de nombreuses erreurs typiques. Saurez-vous toutes les repérer avant de lire la correction ?

\begin{textebox}[A candidate's paragraph — spot the mistakes]
« I am agree with the journalist: actually, the medias give too much informations. I listen often the radio and I am interested by the politics. The news are sometimes surprising, but the advices of experts are useful. Young people depend of their phones and they don't know nothing about traditional newspapers. »
\end{textebox}

Correction : \eng{I agree} (pas de \eng{am}) ; \eng{in fact} ou \eng{nowadays} selon le sens (\eng{actually} ne veut pas dire « actuellement ») ; \eng{the media give too much information} (\eng{media} pas de \eng{-s}, \eng{information} indénombrable) ; \eng{I often listen to the radio} (adverbe avant le verbe + préposition \eng{to}) ; \eng{interested in politics} (sans \eng{the}, préposition \eng{in}) ; \eng{the news is} (\eng{news} singulier) ; \eng{the advice of experts} (\eng{advice} indénombrable) ; \eng{depend on} ; \eng{they know nothing} (négation unique). Une douzaine d'erreurs en quatre phrases : c'est dire l'importance de cette révision !

\courssub{Pas d'auxiliaire \eng{do} en français : questions et négations}

\begin{propriete}[L'auxiliaire do / does / did]
Le français forme ses questions par inversion (\emph{Regardes-tu la télévision ?}) ou par intonation. L'anglais, lui, utilise obligatoirement l'auxiliaire \cle{do} (\eng{does} à la 3\textsuperscript{e} personne du singulier, \eng{did} au prétérit) pour construire questions et négations au présent simple et au prétérit :
\begin{itemize}
\item Question : \textbf{Do/Does/Did + sujet + base verbale} ? \eng{Do you watch TV every evening?} « Regardes-tu la télévision chaque soir ? »
\item Négation : \textbf{sujet + don't/doesn't/didn't + base verbale}. \eng{He doesn't agree with the editor.} « Il n'est pas d'accord avec le rédacteur en chef. »
\end{itemize}
Attention : après \eng{does} et \eng{did}, le verbe revient à la base verbale : \eng{Does she work?} (jamais \emph{does she works}) ; \eng{Did you see it?} (jamais \emph{did you saw}).
Exceptions : on n'utilise pas \eng{do} avec \eng{be} (\eng{Are you tired?}), avec les modaux (\eng{Can you help me?}), ni quand le mot interrogatif est le sujet (\eng{Who wrote this article?} « Qui a écrit cet article ? »).
\end{propriete}

\begin{attention}[« I am agree » : l'erreur reine des francophones]
En français on dit \emph{je suis d'accord}, avec le verbe \emph{être}. En anglais, \eng{agree} est un \textbf{verbe} : on dit \eng{I agree}, \eng{she agrees}, et la négation est \eng{I don't agree}. Jamais \emph{I am agree}. De même, \emph{cela dépend de} se traduit \eng{it depends on} (verbe \eng{depend} + préposition \eng{on}) : \emph{it depend of} est doublement faux (manque du \eng{-s} et mauvaise préposition).
\end{attention}

\courssub{L'ordre des mots : adjectif et adverbe}

\begin{propriete}[Deux ordres inversés par rapport au français]
\begin{enumerate}
\item \textbf{L'adjectif précède le nom} en anglais, alors qu'il suit souvent en français : \eng{a famous journalist} « un journaliste célèbre » ; \eng{a Cameroonian newspaper} « un journal camerounais ». L'adjectif anglais est invariable : \eng{famous journalists} (jamais \emph{famouses}).
\item \textbf{L'adverbe de fréquence} (\eng{always, usually, often, sometimes, never}) se place \textbf{avant le verbe} principal, mais \textbf{après le verbe be} : \eng{I often read the newspaper.} « Je lis souvent le journal. » / \eng{She is always punctual.} « Elle est toujours ponctuelle. »
\end{enumerate}
\end{propriete}

\begin{exemplebox}[Exemples traduits]
\eng{I am interested in social media.} « Je m'intéresse aux réseaux sociaux. »\par
\eng{She is good at interviewing people.} « Elle est douée pour interviewer les gens. »\par
\eng{They never watch the news on television.} « Ils ne regardent jamais les informations à la télévision. »\par
\eng{He bought an interesting magazine.} « Il a acheté un magazine intéressant. »
\end{exemplebox}

\courssub{Les faux amis du domaine des médias}

\begin{definition}[Faux amis (false friends)]
Mots anglais qui ressemblent à des mots français mais ont un sens différent. Ils provoquent des contresens graves à l'écrit comme à l'oral.
\end{definition}

\begin{center}
\begin{tabularx}{\linewidth}{|X|X|X|}
\hline
\rowcolor{popblueL} Mot anglais & Sens réel & Piège français \\
\hline
\eng{actually} & en fait, en réalité & pas « actuellement » (= \eng{currently, at present}) \\
\hline
\eng{to assist (someone)} & aider (quelqu'un) & « assister à » = \eng{to attend} \\
\hline
\eng{a library} & une bibliothèque & « une librairie » = \eng{a bookshop} \\
\hline
\eng{sensible} & raisonnable, sensé & « sensible » = \eng{sensitive} \\
\hline
\eng{to demand} & exiger & « demander » = \eng{to ask (for)} \\
\hline
\eng{a magazine} & un magazine & « un magasin » = \eng{a shop, a store} \\
\hline
\eng{the press} & la presse & « le prestige » = \eng{prestige} \\
\hline
\eng{a reporter} & un reporter, un journaliste & « un rapporteur » = \eng{a rapporteur} \\
\hline
\end{tabularx}
\end{center}

\begin{exemplebox}[Exemples traduits]
\eng{Actually, I prefer the radio to television.} « En fait, je préfère la radio à la télévision. »\par
\eng{A journalist assisted me during the interview.} « Un journaliste m'a aidé pendant l'interview. »\par
\eng{I borrowed this book from the school library.} « J'ai emprunté ce livre à la bibliothèque de l'école. »\par
\eng{It is not sensible to believe everything you read online.} « Il n'est pas raisonnable de croire tout ce qu'on lit en ligne. »\par
\eng{The press reported the event accurately.} « La presse a rapporté l'événement avec exactitude. »
\end{exemplebox}

\courssub{Prépositions piégeuses et indénombrables}

\begin{propriete}[Prépositions à mémoriser avec le verbe]
En anglais, la préposition ne se devine pas depuis le français : elle s'apprend avec le verbe.
\begin{itemize}
\item \eng{to listen \textbf{to}} : \eng{We listen to the radio every morning.} « Nous écoutons la radio chaque matin. »
\item \eng{to wait \textbf{for}} : \eng{They are waiting for the results.} « Ils attendent les résultats. »
\item \eng{to depend \textbf{on}} : \eng{It depends on the price.} « Cela dépend du prix. »
\item \eng{to be interested \textbf{in}} : \eng{I am interested in politics.} « Je m'intéresse à la politique. »
\item \eng{to be good \textbf{at}} : \eng{She is good at writing articles.} « Elle est douée pour écrire des articles. »
\item \eng{to look \textbf{at}} / \eng{to watch} : \eng{Look at the headline!} « Regarde le titre ! » / \eng{Watch the news.} « Regarde le journal télévisé. »
\end{itemize}
\end{propriete}

\begin{propriete}[Indénombrables : singuliers en anglais, pluriels en français]
\eng{information}, \eng{news}, \eng{advice}, \eng{homework}, \eng{furniture}, \eng{money}, \eng{equipment}, \eng{coverage} sont \textbf{indénombrables} en anglais : pas de \eng{-s}, verbe au singulier, pas d'article \eng{a/an}. Pour les compter : \eng{a piece of news / advice / information}, \eng{an item of news}, \eng{a piece of equipment}.
\end{propriete}

\begin{exemplebox}[Exemples traduits]
\eng{The news is surprising.} « Les nouvelles sont surprenantes. » (\eng{news} est singulier)\par
\eng{He gave me useful advice.} « Il m'a donné des conseils utiles. » (\eng{advice} est indénombrable)\par
\eng{This information is not reliable.} « Ces informations ne sont pas fiables. »\par
\eng{The equipment is expensive.} « L'équipement est cher. »
\end{exemplebox}

\begin{propriete}[Négation unique]
Le français utilise deux négations (\emph{je ne sais rien}, \emph{il ne voit personne}). L'anglais n'en admet qu'\textbf{une seule} : soit \eng{not} avec le verbe, soit un mot négatif (\eng{nobody, never, nothing, no one, neither}) avec un verbe \textbf{affirmatif} : \eng{I know nothing.} « Je ne sais rien. » (jamais \emph{I don't know nothing}) ; \eng{Nobody came.} « Personne n'est venu. » ; \eng{She never lies.} « Elle ne ment jamais. » ; \eng{I have neither time nor money.} « Je n'ai ni temps ni argent. »
\end{propriete}

\begin{methode}[Le contrôle anti-calque en quatre vérifications]
Avant de rendre toute production écrite, relisez chaque phrase en posant ces quatre questions :
\begin{enumerate}
\item \textbf{Auxiliaire} : ma question ou ma négation a-t-elle bien \eng{do/does/did} (sauf avec \eng{be} et les modaux) ?
\item \textbf{Ordre des mots} : l'adjectif est-il avant le nom ? L'adverbe de fréquence est-il avant le verbe (après \eng{be}) ?
\item \textbf{Préposition} : le verbe est-il suivi de la bonne préposition (\eng{listen to, wait for, depend on, interested in, good at, look at}) ?
\item \textbf{Dénombrabilité} : \eng{news, information, advice, equipment} sont-ils bien au singulier, sans \eng{-s} ni \eng{a} ?
\end{enumerate}
\end{methode}

\begin{popfigure}
\begin{center}
\begin{tikzpicture}[font=\small]
\draw[thick, popblue] (0,0) rectangle (6.4,8.2);
\draw[thick, poporange] (7.2,0) rectangle (13.6,8.2);
\node[fill=popblue, text=white, font=\bfseries, minimum width=6.4cm] at (3.2,7.7) {Français (erreur typique)};
\node[fill=poporange, text=white, font=\bfseries, minimum width=6.4cm] at (10.4,7.7) {Anglais correct};
\node[align=left, text width=5.9cm] at (3.2,6.6) {Je suis d'accord. $\to$ \emph{I am agree}};
\node[align=left, text width=5.9cm] at (10.4,6.6) {\eng{I agree.}};
\node[align=left, text width=5.9cm] at (3.2,5.5) {Il écoute la radio. $\to$ \emph{He listens the radio}};
\node[align=left, text width=5.9cm] at (10.4,5.5) {\eng{He listens to the radio.}};
\node[align=left, text width=5.9cm] at (3.2,4.4) {Un journaliste célèbre $\to$ \emph{a journalist famous}};
\node[align=left, text width=5.9cm] at (10.4,4.4) {\eng{a famous journalist}};
\node[align=left, text width=5.9cm] at (3.2,3.3) {Les nouvelles sont bonnes. $\to$ \emph{The news are good}};
\node[align=left, text width=5.9cm] at (10.4,3.3) {\eng{The news is good.}};
\node[align=left, text width=5.9cm] at (3.2,2.2) {Je ne sais rien. $\to$ \emph{I don't know nothing}};
\node[align=left, text width=5.9cm] at (10.4,2.2) {\eng{I know nothing.}};
\node[align=left, text width=5.9cm] at (3.2,1.1) {Actuellement, je lis. $\to$ \emph{Actually, I read}};
\node[align=left, text width=5.9cm] at (10.4,1.1) {\eng{At present, I am reading.}};
\draw[-{Stealth}, thick, popdark] (6.4,6.6) -- (7.2,6.6);
\draw[-{Stealth}, thick, popdark] (6.4,5.5) -- (7.2,5.5);
\draw[-{Stealth}, thick, popdark] (6.4,4.4) -- (7.2,4.4);
\draw[-{Stealth}, thick, popdark] (6.4,3.3) -- (7.2,3.3);
\draw[-{Stealth}, thick, popdark] (6.4,2.2) -- (7.2,2.2);
\draw[-{Stealth}, thick, popdark] (6.4,1.1) -- (7.2,1.1);
\end{tikzpicture}
\end{center}
\legende{Six erreurs typiques des francophones et leur correction}
\end{popfigure}

\begin{exoresolu}[Corriger les calques]
Énoncé : corrigez les phrases suivantes, toutes fautives. 1) \emph{I am agree with this article.} 2) \emph{She depends of her smartphone.} 3) \emph{The news are very interesting.} 4) \emph{He gave me many advices.} 5) \emph{They don't know nothing about the press.} 6) \emph{I read often a magazine interesting.}
\tcblower
Solution détaillée :
\begin{enumerate}
\item \eng{I agree with this article.} — \eng{agree} est un verbe : pas de \eng{am}.
\item \eng{She depends on her smartphone.} — préposition \eng{on}, et \eng{-s} à la 3\textsuperscript{e} personne.
\item \eng{The news is very interesting.} — \eng{news} est indénombrable, donc singulier.
\item \eng{He gave me a lot of useful advice.} — \eng{advice} ne prend jamais de \eng{-s} ; on peut dire \eng{many pieces of advice}.
\item \eng{They know nothing about the press.} (ou \eng{They don't know anything}) — une seule négation.
\item \eng{I often read an interesting magazine.} — adverbe de fréquence avant le verbe, adjectif avant le nom, article \eng{an} devant voyelle.
\end{enumerate}
\end{exoresolu}

\begin{savaistu}[« Fake news » et le mot \eng{news}]
L'expression \eng{fake news} se traduit en français par « fausses informations » ou « infox ». Mais attention : en anglais, \eng{news} ne prend \textbf{jamais} de \eng{-s}, même s'il en a l'air ! Pour dire « une nouvelle », on dit \eng{a piece of news} ou \eng{a news item}. Le mot vient à l'origine de \eng{new things} — « les choses nouvelles » — mais il est aujourd'hui parfaitement singulier : \eng{the news \textbf{is} at eight o'clock}.
\end{savaistu}

\begin{exercice}[À vous de jouer]
Traduisez en anglais : 1) Je ne regarde jamais la télévision le soir. 2) Elle s'intéresse aux journaux camerounais. 3) En fait, il n'est pas d'accord avec le présentateur. 4) Ces informations sont fausses. 5) Personne n'attend le journal de vingt heures. 6) Il m'a donné un conseil utile.
\end{exercice}

\begin{corrige}[Exercice « À vous de jouer »]
1) \eng{I never watch television in the evening.} 2) \eng{She is interested in Cameroonian newspapers.} 3) \eng{Actually, he doesn't agree with the presenter.} 4) \eng{This information is false.} 5) \eng{Nobody is waiting for the eight o'clock news.} 6) \eng{He gave me a useful piece of advice.}
\end{corrige}

\coursec{Entraînement intégré type Bac}

Nous avons vu, dans la section précédente, les différences majeures entre le français et l'anglais et les pièges qui guettent les francophones. Le moment est venu de mettre toutes ces connaissances en application dans des exercices complets, calqués sur le format de l'épreuve d'anglais du Baccalauréat A : complétion de texte, transformations de phrases, choix de modaux, correction de phrases fautives et production écrite. Chaque réponse devra être \cle{justifiée} par la règle de grammaire correspondante, comme le jury l'exige.

\courssub{Méthode : réussir les exercices de grammaire au Bac}

\begin{methode}[La démarche en cinq étapes]
Pour tout exercice de grammaire au Bac, applique systématiquement la méthode suivante :
\begin{enumerate}
\item \textbf{Lis toute la phrase} (ou tout le texte) avant de répondre : le sens général donne souvent la solution.
\item \textbf{Repère les marqueurs temporels} : \eng{yesterday}, \eng{since 2020}, \eng{every year}, \eng{tomorrow} indiquent le temps attendu.
\item \textbf{Identifie la structure attendue} : passive ? discours rapporté ? modal ? relatif ? Cherche les indices (agent introduit par \eng{by}, verbe introducteur au passé, etc.).
\item \textbf{Conjugue ou transforme} en appliquant la règle, puis \textbf{relis} en vérifiant l'accord (sujet singulier ou pluriel, terminaison en \eng{-s} à la 3\textsuperscript{e} personne).
\item \textbf{Justifie} ta réponse par la règle : « present simple passive car \eng{every year} et l'agent est introduit par \eng{by} ».
\end{enumerate}
\end{methode}

\begin{exemplebox}[Observer avant d'agir]
Observe cette phrase type Bac : \eng{Many programmes ... (produce) in Douala every year.}
\begin{itemize}
\item Marqueur temporel : \eng{every year} $\rightarrow$ present simple.
\item Le sujet \eng{programmes} ne fait pas l'action, il la subit $\rightarrow$ voix passive.
\item Forme : \emph{be} au present + participe passé $\rightarrow$ \eng{are produced}.
\end{itemize}
\eng{Many programmes are produced in Douala every year.} « De nombreuses émissions sont produites à Douala chaque année. »
\end{exemplebox}

\begin{attention}[La justification rapporte des points]
Au Bac, une réponse juste sans justification claire de la forme (temps, auxiliaire, structure) perd souvent des points. Vérifie toujours la \cle{terminaison verbale} : le \eng{-s} de la 3\textsuperscript{e} personne, le participe passé des verbes irréguliers (\eng{written}, \eng{spoken}, \eng{broadcast}), l'auxiliaire \eng{do/does} aux questions et négations du present simple.
\end{attention}

\courssub{Exercice 1 : compléter un texte avec le verbe au temps correct}

\begin{exoresolu}[Texte à compléter : the media in Cameroon]
Énoncé : mets les verbes entre parenthèses au temps qui convient.
\eng{The press in Cameroon (1. exist) ... for more than a century. The first newspapers (2. appear) ... during the colonial period. Today, hundreds of radio stations (3. broadcast) ... news in English, French and local languages. Last year, the government (4. launch) ... a new media training programme. Since then, many young journalists (5. train) ... in Yaoundé and Douala. At the moment, a national debate on fake news (6. take place) ... in schools. Experts believe that social media (7. transform) ... the way people get information. If nothing is done, misinformation (8. spread) ... even faster. By 2030, most Cameroonians (9. probably get) ... their news online. The media landscape (10. already change) ... a lot since 2010.}
\tcblower
Solution détaillée :
\begin{enumerate}
\item \eng{has existed} — present perfect : durée depuis le passé jusqu'à aujourd'hui (\eng{for more than a century}).
\item \eng{appeared} — preterite : période révolue (\eng{during the colonial period}).
\item \eng{broadcast} — present simple, vérité générale (\eng{today}, habitude) ; participe passé et preterite identiques à la base verbale.
\item \eng{launched} — preterite : \eng{last year}, marqueur de passé terminé.
\item \eng{have been trained} — present perfect passive : \eng{since then} + les journalistes subissent la formation.
\item \eng{is taking place} — present continuous : \eng{at the moment}.
\item \eng{will transform} / \eng{are transforming} — futur (prédiction) ou present continuous (tendance en cours) après \eng{experts believe that}. \eng{Social media} prend un verbe au pluriel en anglais britannique.
\item \eng{will spread} — futur après une condition (\eng{if nothing is done}) : first conditional.
\item \eng{will probably get} — futur : \eng{by 2030} (prédiction future).
\item \eng{has already changed} — present perfect : \eng{already} et \eng{since 2010} (durée depuis le passé).
\end{enumerate}
\end{exoresolu}

\courssub{Exercice 2 : de l'actif au passif}

\begin{exercice}[Transforme au passif]
Mets les phrases suivantes à la voix passive.
\begin{enumerate}
\item \eng{Journalists interview the minister every week.}
\item \eng{Hackers stole the newspaper's database last night.}
\item \eng{The students are preparing a school radio programme.}
\item \eng{MTN has sponsored the football championship since 2015.}
\item \eng{People will watch the match on television tomorrow.}
\end{enumerate}
\end{exercice}

\begin{corrige}[Exercice 2]
\begin{enumerate}
\item \eng{The minister is interviewed every week.} — present simple passive : \emph{is} + participe passé ; l'agent peut être omis car il est évident.
\item \eng{The newspaper's database was stolen last night.} — preterite passive : \emph{was} + \eng{stolen} (irrégulier : steal--stole--stolen).
\item \eng{A school radio programme is being prepared.} — present continuous passive : \emph{is being} + participe passé.
\item \eng{The football championship has been sponsored since 2015.} — present perfect passive : \emph{has been} + participe passé.
\item \eng{The match will be watched on television tomorrow.} — futur passive : \emph{will be} + participe passé.
\end{enumerate}
Règle : l'objet de la phrase active devient sujet ; \emph{be} prend le temps du verbe actif ; le verbe principal passe au participe passé.
\end{corrige}

\courssub{Exercice 3 : le discours rapporté}

\begin{exercice}[Mets au discours rapporté]
Rapporte ces paroles en commençant par \eng{He said that...} ou le verbe approprié.
\begin{enumerate}
\item \eng{“I read the news every morning,” says Paul.}
\item \eng{“We are broadcasting live from Buea,” the reporter said.}
\item \eng{“I will launch a new channel next year,” the director announced.}
\item \eng{“Don't believe everything you see online,” the teacher told the students.}
\item \eng{“Have you ever interviewed a minister?” she asked me.}
\end{enumerate}
\end{exercice}

\begin{corrige}[Exercice 3]
\begin{enumerate}
\item \eng{Paul says (that) he reads the news every morning.} — verbe introducteur au present : pas de recul des temps.
\item \eng{The reporter said (that) they were broadcasting live from Buea.} — recul : present continuous $\rightarrow$ past continuous.
\item \eng{The director announced (that) he would launch a new channel the following year.} — \eng{will} $\rightarrow$ \eng{would} ; \eng{next year} $\rightarrow$ \eng{the following year}.
\item \eng{The teacher told the students not to believe everything they saw online.} — impératif négatif rapporté : \emph{tell + complément + not to} + base verbale ; recul \eng{see} $\rightarrow$ \eng{saw}.
\item \eng{She asked me if I had ever interviewed a minister.} — question sans mot interrogatif : \emph{asked + if} ; present perfect $\rightarrow$ past perfect ; ordre sujet--verbe (pas d'auxiliaire \eng{do}).
\end{enumerate}
\end{corrige}

\courssub{Exercice 4 : choisir le modal approprié}

\begin{exercice}[Quel modal ?]
Complète avec \eng{must}, \eng{should}, \eng{can}, \eng{may}, \eng{mustn't} ou \eng{needn't}.
\begin{enumerate}
\item Journalists \eng{...} check their sources before publishing. (conseil moral)
\item You \eng{...} reveal a source's identity without permission. (interdiction)
\item \eng{...} I use your phone to call the radio station? (demande de permission polie)
\item In Cameroon, you \eng{...} be eighteen to vote. (obligation externe)
\item You \eng{...} buy a newspaper; the news is free online. (absence de nécessité)
\item She speaks four languages, so she \eng{...} work as an international reporter. (capacité)
\end{enumerate}
\end{exercice}

\begin{corrige}[Exercice 4]
\begin{enumerate}
\item \eng{should} — conseil, devoir moral.
\item \eng{mustn't} — interdiction formelle (attention : \eng{mustn't} n'est pas « ne pas être obligé »).
\item \eng{May} / \eng{Can} — demande de permission ; \eng{may} est plus poli.
\item \eng{must} — obligation légale (ou \eng{have to}).
\item \eng{needn't} — « il n'est pas nécessaire de » (faux ami : ne pas traduire par \eng{mustn't}).
\item \eng{can} — capacité.
\end{enumerate}
\end{corrige}

\courssub{Exercice 5 : corriger les erreurs typiques des francophones}

\begin{exercice}[Trouve et corrige l'erreur]
Chaque phrase contient une erreur fréquente chez les francophones. Corrige-la.
\begin{enumerate}
\item \eng{The news are very interesting today.}
\item \eng{She is listening to the radio since two hours.}
\item \eng{He explained me the problem.}
\item \eng{They discussed about the new law.}
\item \eng{My brother he works for a newspaper.}
\item \eng{I am agree with your opinion.}
\item \eng{She told to me that she was tired.}
\item \eng{This is the more interesting article of the week.}
\end{enumerate}
\end{exercice}

\begin{corrige}[Exercice 5]
\begin{enumerate}
\item \eng{The news \textbf{is} very interesting.} — \eng{news} est singulier en anglais, bien que « les nouvelles » soit pluriel en français.
\item \eng{She \textbf{has been listening} to the radio \textbf{for} two hours.} — action commencée dans le passé et qui continue : present perfect continuous ; \eng{for} + durée, \eng{since} + point de départ.
\item \eng{He explained the problem \textbf{to me}.} — \eng{explain something to somebody}, jamais \eng{explain me}.
\item \eng{They discussed the new law.} — \eng{discuss} est transitif direct : pas de \eng{about} (contrairement à « discuter de »).
\item \eng{My brother works for a newspaper.} — pas de double sujet (\eng{my brother he...}), calque du français « mon frère, il... ».
\item \eng{I \textbf{agree} with your opinion.} — \eng{agree} est un verbe, pas un adjectif ; « je suis d'accord » se traduit sans \emph{be}.
\item \eng{She \textbf{told me} that she was tired.} — \eng{tell somebody} sans \eng{to} ; mais \eng{say to somebody}.
\item \eng{This is \textbf{the most} interesting article of the week.} — superlatif des adjectifs longs : \emph{the most} + adjectif.
\end{enumerate}
\end{corrige}

\courssub{Exercice 6 : production écrite guidée}

\begin{exercice}[Paragraph writing]
Rédige un paragraphe de 8 à 10 lignes intitulé \eng{The role of social media in Cameroon today}. Tu dois utiliser au moins : une phrase à la voix passive, un modal, une proposition relative et un comparatif. Souligne ces quatre structures dans ta copie.
\end{exercice}

\begin{textebox}[Model paragraph — The role of social media in Cameroon today]
Social media have become essential in Cameroon today. Every day, millions of messages, photos and videos are shared on platforms such as Facebook, WhatsApp and TikTok. Young people, who represent the majority of users, find information faster on their phones than in newspapers. However, citizens must be careful, because fake news spreads more quickly than verified information. Journalists should check every story before it is published online. Social media can also unite people: during football matches, the whole country celebrates together. In my opinion, these platforms are more powerful than traditional media, but they must be used responsibly. If they are used wisely, they will help Cameroonian democracy grow stronger.
\end{textebox}

Glossaire : \emph{to share} (v., partager) ; \emph{fake news} (n., fausses informations) ; \emph{to spread} (v., se propager) ; \emph{to check} (v., vérifier) ; \emph{wisely} (adv., sagement).

\begin{corrige}[Exercice 6 — grille d'évaluation]
Dans le modèle ci-dessus, on trouve :
\begin{itemize}
\item \textbf{Passive} : \eng{are shared}, \eng{is published}, \eng{are used} — \emph{be} + participe passé.
\item \textbf{Modal} : \eng{must be careful}, \eng{should check}, \eng{can unite}.
\item \textbf{Relative} : \eng{who represent the majority of users} — pronom relatif \eng{who} pour les personnes.
\item \textbf{Comparatif} : \eng{faster than}, \eng{more quickly than}, \eng{more powerful than}.
\end{itemize}
Barème indicatif : contenu et organisation (4 pts), correction grammaticale (8 pts), présence des quatre structures exigées (4 pts), vocabulaire et orthographe (4 pts).
\end{corrige}

\courssub{Correction collective : la grille de relecture en cinq points}

\begin{propriete}[La grille de relecture avant de rendre sa copie]
Avant de rendre toute copie d'anglais, vérifie ces cinq points dans l'ordre :
\begin{enumerate}
\item \textbf{Verbes} : le temps correspond-il aux marqueurs ? La 3\textsuperscript{e} personne du singulier a-t-elle son \eng{-s} ?
\item \textbf{Auxiliaires} : \eng{do/does/did} aux questions et négations ; \emph{be} + participe passé à la voix passive ; \emph{have} + participe passé au present perfect.
\item \textbf{Prépositions} : \eng{listen to}, \eng{depend on}, \eng{interested in}, \eng{good at} — mais \eng{discuss}, \eng{marry}, \eng{phone} sans préposition.
\item \textbf{Ordre des mots} : sujet--verbe--complément ; adverbe de fréquence avant le verbe principal ; pas d'inversion dans les questions indirectes.
\item \textbf{Orthographe des irréguliers} : \eng{written}, \eng{spoken}, \eng{stolen}, \eng{broadcast}, \eng{paid}, \eng{chosen}.
\end{enumerate}
\end{propriete}

\begin{popfigure}
\begin{tikzpicture}[node distance=0.55cm, every node/.style={font=\small},
startstop/.style={rectangle, rounded corners, draw=popgreen, fill=popgreenL, align=center, inner sep=5pt},
decision/.style={diamond, draw=popblue, fill=popblueL, align=center, inner sep=2pt, aspect=2.2},
action/.style={rectangle, draw=poporange, fill=poporangeL, align=center, inner sep=5pt},
arrow/.style={-{Stealth[length=2.5mm]}, thick, popdark}]
\node[startstop] (start) {Relecture de la copie};
\node[decision, below=of start] (tense) {Temps correct ?};
\node[decision, below=of tense] (third) {s à la 3e personne ?};
\node[decision, below=of third] (aux) {Auxiliaire do / be / have ?};
\node[decision, below=of aux] (prep) {Préposition correcte ?};
\node[decision, below=of prep] (order) {Ordre des mots ?};
\node[startstop, below=of order] (end) {Copie prête à rendre};
\node[action, right=1.6cm of third, align=center] (fix){Corriger puis\\reprendre la vérification};
\draw[arrow] (start) -- (tense);
\draw[arrow] (tense) -- node[right, font=\scriptsize]{oui} (third);
\draw[arrow] (third) -- node[right, font=\scriptsize]{oui} (aux);
\draw[arrow] (aux) -- node[right, font=\scriptsize]{oui} (prep);
\draw[arrow] (prep) -- node[right, font=\scriptsize]{oui} (order);
\draw[arrow] (order) -- node[right, font=\scriptsize]{oui} (end);
\draw[arrow] (tense.east) -- node[above, font=\scriptsize]{non} (tense.east -| fix.west);
\draw[arrow] (third.east) -- node[above, font=\scriptsize]{non} (fix.west);
\draw[arrow] (aux.east) -- (aux.east -| fix.west);
\draw[arrow] (prep.east) -- (prep.east -| fix.west);
\draw[arrow] (order.east) -- (order.east -| fix.west);
\draw[arrow] (fix.south) |- ($(order.south)+(0,-0.15)$) -| (end.east);
\end{tikzpicture}
\legende{Organigramme de vérification avant de rendre une copie d'anglais au Bac.}
\end{popfigure}

\begin{center}
\begin{tabularx}{\linewidth}{|X|X|X|}
\hline
\rowcolor{popblueL} Erreur du francophone & Forme correcte & Règle \\
\hline
\eng{The news are...} & \eng{The news is...} & \eng{news} est singulier \\
\hline
\eng{I am agree} & \eng{I agree} & \eng{agree} est un verbe \\
\hline
\eng{He explained me} & \eng{He explained to me} & \emph{explain sth to sb} \\
\hline
\eng{They discussed about} & \eng{They discussed} & transitif direct \\
\hline
\eng{since two hours} & \eng{for two hours} & \emph{for} + durée \\
\hline
\eng{the more interesting} & \eng{the most interesting} & superlatif des adjectifs longs \\
\hline
\end{tabularx}
\end{center}

\begin{savaistu}[Le Bac A aime la passive et le discours rapporté]
Les épreuves d'anglais du Baccalauréat A au Cameroun sanctionnent régulièrement la \cle{voix passive} et le \cle{discours rapporté} dans les exercices de transformation de phrases. Les candidats qui maîtrisent le recul des temps (\eng{will} $\rightarrow$ \eng{would}, \eng{is} $\rightarrow$ \eng{was}, \eng{have} $\rightarrow$ \eng{had}) et la formation du participe passé gagnent des points précieux. Entraîne-toi chaque semaine à transformer cinq phrases : c'est l'exercice le plus rentable de l'épreuve.
\end{savaistu}

\newpage

\coursec{Activité d'intégration}

\coursec{Lesson 50 — Grammar: General Revision}

\courssub{Objectifs de la leçon}

Cette leçon de révision générale clôt le chapitre \eng{Media and Communication}. Elle vise à consolider, dans une perspective communicative et évaluative (type Baccalauréat), l'ensemble des structures grammaticales étudiées au cours de l'année. L'élève doit être capable de :

\begin{itemize}
\item \textbf{Identifier et produire} correctement les \cle{temps verbaux} narratifs et informationnels (\eng{present perfect}, \eng{past simple}, \eng{present continuous}, \eng{future forms}) ;
\item \textbf{Manipuler} la \cle{voix passive} (forme \eng{be} + \cle{participe passé}) et le \cle{discours rapporté} (concordance des temps / \cle{backshifting}) ;
\item \textbf{Utiliser avec justesse} les \cle{auxiliaires modaux} (\eng{must, should, can, may, will, would, could}) sans \eng{to} ;
\item \textbf{Construire} des \cle{propositions relatives} (\eng{who, which, that, whose}) et des structures de \cle{comparaison} (comparatif de supériorité, superlatif, irréguliers \eng{better/worse}) ;
\item \textbf{Éviter les erreurs systématiques} des francophones (faux amis, indénombrables, ordre des mots, \eng{actually} vs \eng{currently}) ;
\item \textbf{Rédiger et présenter} un \cle{bulletin d'information} authentique (écrit + oral) intégrant naturellement ces structures.
\end{itemize}

\courssub{Observation et diagnostic : un bulletin d'information modèle}

\courssub{Situation de communication}

Le club d'anglais de ton lycée, le \eng{English Media Club}, produit le bulletin d'information hebdomadaire diffusé en anglais lors de la cérémonie du drapeau du lundi matin. Le proviseur a annoncé que le meilleur bulletin du trimestre serait lu à l'antenne sur la radio communautaire locale. Ton groupe (4 élèves : \eng{news anchor}, 2 \eng{reporters}, 1 \eng{editor}) doit rédiger et présenter un bulletin de \textbf{8 à 10 phrases} sur la vie du lycée et l'actualité médiatique camerounaise.

Voici l'extrait du bulletin de la semaine dernière (modèle de référence) :

\begin{textebox}[Last week's bulletin — extract (model)]
Good morning, everyone. Here is the news from our school. Last Friday, a new radio club was officially launched by the English department. The principal said that the club would help students to improve their spoken English. More than sixty students have already registered, which is much higher than last year's number. According to our reporters, the club will broadcast short programmes every Wednesday. Students who want to join should contact the club secretary before Friday. Finally, mobile phones must be switched off during recordings. That is all for this week. Thank you for listening.
\end{textebox}

\begin{definition}[Glossaire utile — Vocabulaire des médias scolaires]
\begin{itemize}
\item \eng{bulletin} (n.) : bulletin, journal d'informations (syn. \eng{newscast}) ;
\item \eng{to launch} (v.) : lancer, inaugurer officiellement ;
\item \eng{to register} (v.) : s'inscrire (attention : \eng{register for} un cours, \eng{register at} un établissement) ;
\item \eng{to broadcast} (v.) : diffuser (radio, TV) — \emph{irrégulier : broadcast / broadcast / broadcast} ;
\item \eng{to switch off} (v.) : éteindre (appareil) — phrasal verb séparable : \eng{switch it off} ;
\item \eng{secretary} (n.) : secrétaire (personne) — \textbf{faux ami} : ne signifie ni « meuble » (\eng{cabinet}), ni « secrétariat » (\eng{secretary's office}) ;
\item \eng{flag ceremony} (n.) : cérémonie de lever des couleurs / cérémonie du drapeau ;
\item \eng{inter-school tournament} (n.) : tournoi inter-écoles.
\end{itemize}
\end{definition}

\begin{experience}[Activité de démarrage : \eng{Spot the Grammar} (10 min)]
En binôme, relisez le texte modèle et surlignez :
\begin{enumerate}
\item 2 verbes au \eng{present perfect} (\eng{has been launched, have registered}) ;
\item 1 verbe au \eng{past simple} (\eng{said}) ;
\item 1 structure de \cle{voix passive} (\eng{was launched}) ;
\item 1 \cle{discours rapporté} avec recul (\eng{said that... would help}) ;
\item 2 \cle{modaux} différents (\eng{will, should, must}) ;
\item 1 \cle{proposition relative} (\eng{who want to join}) ;
\item 1 \cle{comparatif} (\eng{higher than}).
\end{enumerate}
Mettez en commun au tableau. L'enseignant valide et nomme chaque structure.
\end{experience}

\courssub{Révision des structures clés}

\courssub{Révision des temps verbaux : choisir le bon temps pour l'information}

\begin{propriete}[Temps verbaux du journalisme scolaire — Forme, Emploi, Marqueurs]
\begin{center}
\begin{tabularx}{\linewidth}{|X|X|X|X|}
\hline
\rowcolor{popblueL} Temps & Forme (affirmatif) & Emploi principal dans un bulletin & Marqueurs typiques \\
\hline
\eng{Present Perfect} & \eng{have/has + V-en} & Fait récent, résultat présent, bilan & \eng{already, just, yet, since, for, this week} \\
\hline
\eng{Past Simple} & \eng{V-ed / 2\textsuperscript{e} colonne} & Événement daté, achevé dans le passé & \eng{yesterday, last Friday, in 2023, ago} \\
\hline
\eng{Present Continuous} & \eng{be + V-ing} & Action en cours, projet proche & \eng{now, at the moment, this Monday} \\
\hline
\eng{Future (\eng{will})} & \eng{will + BV} & Prédiction, annonce officielle, décision instantanée & \eng{tomorrow, next week, will} \\
\hline
\eng{Be going to} & \eng{be going to + BV} & Intention, plan déjà décidé, preuve visible & \eng{soon, later, look at those clouds!} \\
\hline
\end{tabularx}
\end{center}
\end{propriete}

\begin{attention}[Pièges francophones — Temps verbaux]
\begin{itemize}
\item \textbf{Present Perfect vs Past Simple} : \eng{I have seen him yesterday} $\times$ $\rightarrow$ \eng{I \textbf{saw} him yesterday} (daté = past simple). \eng{He has arrived} (maintenant il est là) vs \eng{He arrived at 10} (fait passé).
\item \textbf{Present Continuous pour le futur} : \eng{The club meets tomorrow} $\times$ (horaire) $\rightarrow$ \eng{The club \textbf{is meeting} tomorrow} (projet) ou \eng{The club \textbf{will meet}} (prédiction).
\item \textbf{Since / For} : \eng{since} + point de départ (\eng{since Monday}) ; \eng{for} + durée (\eng{for three days}).
\end{itemize}
\end{attention}

\courssub{Voix passive et discours rapporté : l'objectivité journalistique}

\begin{propriete}[Voix passive — Structure et transformation]
\textbf{Forme} : Sujet + \eng{be} (temps voulu) + \cle{participe passé} (+ \eng{by} + agent si nécessaire).\\
\textbf{Transformation} : L'objet (COD) de l'actif devient sujet du passif.
\begin{center}
\begin{tabular}{|l|l|l|}
\hline
\rowcolor{popblueL} Temps & Actif & Passif \\
\hline
Present Simple & \eng{They broadcast the news.} & \eng{The news \textbf{is broadcast}.} \\
\hline
Past Simple & \eng{They launched the club.} & \eng{The club \textbf{was launched}.} \\
\hline
Present Perfect & \eng{They have created a club.} & \eng{A club \textbf{has been created}.} \\
\hline
Future (\eng{will}) & \eng{They will broadcast it.} & \eng{It \textbf{will be broadcast}.} \\
\hline
Modal & \eng{Students must switch off phones.} & \eng{Phones \textbf{must be switched off}.} \\
\hline
\end{tabular}
\end{center}
\emph{Verbes à particule (phrasal verbs) au passif : la particule reste collée au verbe.} \eng{The meeting was put off.} (pas \eng{put the meeting off}).
\end{propriete}

\begin{propriete}[Discours rapporté (Reported Speech) — Règle du recul (\eng{Backshifting})]
Si le verbe introducteur est au \textbf{passé} (\eng{said, told, announced, declared}), les temps reculent d'un cran :
\begin{center}
\begin{tabularx}{\linewidth}{|X|X|X|}
\hline
\rowcolor{popblueL} Discours direct & Discours rapporté (après \eng{said that}) & Exemple \\
\hline
Present Simple & Past Simple & \eng{“I \textbf{know} the truth.”} $\rightarrow$ He said he \textbf{knew} the truth. \\
\hline
Present Continuous & Past Continuous & \eng{“I \textbf{am writing}.”} $\rightarrow$ She said she \textbf{was writing}. \\
\hline
Present Perfect / Past Simple & Past Perfect & \eng{“I \textbf{have finished}.”} $\rightarrow$ He said he \textbf{had finished}. \\
\hline
\eng{will} & \eng{would} & \eng{“I \textbf{will} come.”} $\rightarrow$ She said she \textbf{would} come. \\
\hline
\eng{can} & \eng{could} & \eng{“I \textbf{can} help.”} $\rightarrow$ He said he \textbf{could} help. \\
\hline
\eng{must} & \eng{had to} (obligation) / \eng{must} (déduction) & \eng{“I \textbf{must} go.”} $\rightarrow$ He said he \textbf{had to} go. \\
\hline
\end{tabularx}
\end{center}
\textbf{Changements annexes} : \eng{now} $\rightarrow$ \eng{then} ; \eng{today} $\rightarrow$ \eng{that day} ; \eng{tomorrow} $\rightarrow$ \eng{the next/following day} ; \eng{here} $\rightarrow$ \eng{there} ; \eng{this} $\rightarrow$ \eng{that}.
\end{propriete}

\begin{exemplebox}[Exemples contrastés Actif/Passif \& Direct/Rapporté]
\begin{itemize}
\item \eng{Active: The principal opened the ceremony.} $\rightarrow$ \eng{Passive: The ceremony was opened by the principal.}
\item \eng{Direct: “We will win the match,” said the captain.} $\rightarrow$ \eng{Reported: The captain said (that) they would win the match.}
\item \eng{Direct: “Students must wear badges,” said the teacher.} $\rightarrow$ \eng{Reported: The teacher said (that) students had to wear badges.}
\end{itemize}
\end{exemplebox}

\courssub{Modaux, relatifs et comparaison : précision et nuance}

\begin{propriete}[Auxiliaires modaux — Valeurs et syntaxe]
\textbf{Règle d'or} : Modal + \cle{Base Verbale} (sans \eng{to}, pas de \eng{-s} à la 3\textsuperscript{e} pers., pas de \eng{do/does/did}).
\begin{center}
\begin{tabularx}{\linewidth}{|X|X|X|}
\hline
\rowcolor{popblueL} Modal & Valeur principale (contexte bulletin) & Exemple \\
\hline
\eng{must} & Obligation forte, règle interne, déduction forte & \eng{Phones must be switched off.} / \eng{He must be the principal.} \\
\hline
\eng{have to} & Obligation externe, contrainte factuelle & \eng{We have to submit the text by noon.} \\
\hline
\eng{should} & Conseil, recommandation, attente normale & \eng{Students should register early.} \\
\hline
\eng{can / could} & Capacité, permission, possibilité & \eng{You can listen online.} / \eng{Could you repeat?} \\
\hline
\eng{may / might} & Possibilité incertaine, permission formelle & \eng{The programme may start late.} \\
\hline
\eng{will / would} & Futur, volonté, habitude passée (\eng{would}) & \eng{The club will broadcast.} / \eng{He would always arrive early.} \\
\hline
\end{tabularx}
\end{center}
\end{propriete}

\begin{propriete}[Propositions relatives — Choix du pronom]
\begin{center}
\begin{tabularx}{\linewidth}{|X|X|X|}
\hline
\rowcolor{popblueL} Antécédent & Pronom relatif (sujet/objet) & Exemple \\
\hline
Personne & \eng{who} (sujet), \eng{who(m)} (objet), \eng{that} & \eng{The reporter \textbf{who} interviewed the principal...} \\
\hline
Chose / Idée & \eng{which} (sujet/objet), \eng{that} & \eng{The programme \textbf{which/that} was broadcast...} \\
\hline
Possession (pers./chose) & \eng{whose} + nom & \eng{The student \textbf{whose} article won...} \\
\hline
Phrase entière (commentaire) & \eng{which} (non restrictive, précédé de virgule) & \eng{The club won, \textbf{which} pleased everyone.} \\
\hline
\end{tabularx}
\end{center}
\textbf{Attention} : Pas de pronom personnel redondant dans la relative. \eng{The man who he spoke...} $\times$ $\rightarrow$ \eng{The man who spoke...} ou \eng{The man (who(m)) he spoke to...}
\end{propriete}

\begin{propriete}[Comparatif et superlatif — Formation et irréguliers]
\begin{center}
\begin{tabularx}{\linewidth}{|X|X|X|}
\hline
\rowcolor{popblueL} Type & Règle & Exemples \\
\hline
Court (1 syllabe) & \eng{-er than} / \eng{the -est} & \eng{high $\rightarrow$ higher than / the highest} \\
\hline
Long ($\ge$ 2 syll.) & \eng{more ... than} / \eng{the most ...} & \eng{interesting $\rightarrow$ more interesting / the most interesting} \\
\hline
Terminé en \eng{-y} & \eng{-ier than} / \eng{the -iest} & \eng{early $\rightarrow$ earlier / the earliest} \\
\hline
\eng{-ow, -er, -le} & \eng{-er / -est} (souvent) & \eng{narrow $\rightarrow$ narrower ; simple $\rightarrow$ simpler} \\
\hline
\textbf{Irréguliers} & \textbf{Par cœur} & \eng{good/well $\rightarrow$ better / best} ; \eng{bad/badly $\rightarrow$ worse / worst} ; \eng{far $\rightarrow$ farther/further / farthest/furthest} ; \eng{much/many $\rightarrow$ more / most} ; \eng{little $\rightarrow$ less / least} \\
\hline
\end{tabularx}
\end{center}
\textbf{Comparatif d'égalité} : \eng{as + adj + as}. \textbf{Infériorité} : \eng{less ... than} / \eng{not as ... as}.
\end{propriete}

\courssub{Comparaison français/anglais et pièges des francophones}

\begin{attention}[Les 10 erreurs « tueuses » en Terminale — Check-list de relecture]
\begin{enumerate}
\item \textbf{Indénombrables} : \eng{information, news, advice, equipment, furniture, homework} $\rightarrow$ \textbf{jamais de \eng{s}, jamais \eng{a/an}}. \eng{Here is the news.} (pas \eng{Here are the news}). \eng{A piece of advice.}
\item \textbf{Modal + \eng{to}} : \eng{Students must to come} $\times$ $\rightarrow$ \eng{Students \textbf{must come}}. \eng{Have to} est le seul semi-modal avec \eng{to}.
\item \textbf{Backshifting oublié} : \eng{He said he will come} $\times$ $\rightarrow$ \eng{He said he \textbf{would} come}.
\item \textbf{Participe passé passif} : \eng{The match was win} $\times$ $\rightarrow$ \eng{The match was \textbf{won}}. Vérifier la liste des verbes irréguliers (\eng{broadcast/broadcast/broadcast}, \eng{write/wrote/written}).
\item \textbf{Comparatif \eng{good}} : \eng{more good} $\times$, \eng{gooder} $\times$ $\rightarrow$ \eng{\textbf{better}}. \eng{Bad $\rightarrow$ worse}.
\item \textbf{\eng{Actually} vs \eng{Currently}} : \eng{Actually} = « en fait, en réalité » (opposition). \eng{Currently / At present} = « actuellement » (temps).
\item \textbf{\eng{Since} vs \eng{For}} : \eng{I am here since Monday} $\times$ $\rightarrow$ \eng{I \textbf{have been} here \textbf{since} Monday} (point de départ) / \eng{\textbf{for} three days} (durée).
\item \textbf{Relatif + pronom redondant} : \eng{The girl who I know her} $\times$ $\rightarrow$ \eng{The girl \textbf{who(m) I know}} ou \eng{The girl \textbf{I know}}.
\item \textbf{Ordre des mots (Sujet - Verbe - Complément)} : \eng{Arrived the bus} $\times$ $\rightarrow$ \eng{\textbf{The bus arrived}}. Pas d'inversion en affirmatif (sauf adverbes négatifs en tête : \eng{Never have I seen...}).
\item \textbf{Faux amis fréquents} : \eng{Eventually} = « finalement » (pas « éventuellement » $\rightarrow$ \eng{possibly}) ; \eng{Assist} = « aider » (pas « assister à» $\rightarrow$ \eng{attend}) ; \eng{Library} = « bibliothèque » (pas « librairie » $\rightarrow$ \eng{bookshop}) ; \eng{Realise} = « prendre conscience » (pas « réaliser » un projet $\rightarrow$ \eng{carry out/achieve}).
\end{enumerate}
\end{attention}

\begin{savaistu}[Le paysage médiatique camerounais et l'anglais]
Le Cameroun est officiellement bilingue (français/anglais). Les médias publics (\eng{CRTV} — Cameroon Radio Television) diffusent des journaux en anglais (\eng{CRTV News}) et en français. De nombreuses radios communautaires (\eng{community radios}) à Buea, Bamenda, Limbe, Douala ou Yaoundé proposent des créneaux en anglais (pidgin ou standard). Les lycées anglophones et bilingues ont souvent un \eng{Press Club} ou \eng{Media Club} qui forme aux métiers du journalisme (reporting, anchoring, editing). Maîtriser l'anglais des médias, c'est accéder à l'information nationale et internationale et valoriser son profil pour l'université ou l'emploi.
\end{savaistu}

\courssub{Entraînement intégré type Bac : Production du bulletin}

\courssub{Consignes de l'activité finale (55 min)}

\begin{enumerate}
\item \textbf{Brainstorming (10 min).} En groupe, listez \textbf{6 faits réels} (anglais) : vie du lycée (club, match, visite, examen, campagne \eng{fake news}, prix, cantine) + actualité médias (CRTV, radio locale, réseau social, influenceur). Choisissez les 4-5 plus porteurs.
\item \textbf{Planification (5 min).} Attribuez les rôles : \eng{Anchor} (présentateur), \eng{Reporter 1}, \eng{Reporter 2}, \eng{Editor} (relecteur). Ordonnez : \eng{Headline} (info majeure) $\rightarrow$ \eng{Details} $\rightarrow$ \eng{Soft news} (info légère) $\rightarrow$ \eng{Closing}.
\item \textbf{Rédaction collaborative (20 min).} Rédigez \textbf{8–10 phrases} respectant \textbf{obligatoirement} les 5 contraintes grammaticales ci-dessous (cochez au fur et à mesure) :
\begin{itemize}
\item[$\square$] 1 \textbf{Voix passive} (\eng{be + V-en}) : \eng{A new studio has been equipped.}
\item[$\square$] 1 \textbf{Discours rapporté} (recul) : \eng{The principal announced that the festival would start...}
\item[$\square$] 2 \textbf{Modaux différents} (\eng{must, should, can, may, will, would, could, have to}) : \eng{Students must... / Reporters should...}
\item[$\square$] 1 \textbf{Relative} (\eng{who, which, that, whose}) : \eng{Students who wish to participate...}
\item[$\square$] 1 \textbf{Comparatif} (\eng{-er/more... than, better/worse}) : \eng{More participants than last year.}
\end{itemize}
\item \textbf{Auto-vérification (10 min).} Relisez : temps justifiés ? concordance respectée ? participes passés exacts ? modaux sans \eng{to} ? relatifs bien choisis ?
\item \textbf{Relecture croisée (10 min).} Échangez avec un autre groupe. Utilisez la \textbf{Grille en 5 points} ci-dessous. Annotez le texte, rendez-le.
\item \textbf{Finalisation \& Présentation orale (15 min).} Intégrez corrections. Présentation debout : \eng{Anchor} lit, \eng{Reporters} ajoutent 1 phrase commentaire chacun. Critères : clarté, débit, intonation (montante/descendante), contact visuel.
\end{enumerate}

\begin{methode}[Grille de relecture croisée — Peer Review Checklist (5 points)]
L'\eng{Editor} du groupe voisin coche et commente :
\begin{enumerate}
\item \textbf{Passive Voice} : $\square$ Yes $\square$ No — Structure \eng{be + V-en} correcte ? Temps du \eng{be} justifié ? Agent (\eng{by...}) pertinent ou omis ?
\item \textbf{Reported Speech} : $\square$ Yes $\square$ No — Verbe introducteur au passé ? Recul appliqué (\eng{will$\rightarrow$would, can$\rightarrow$could, present$\rightarrow$past, past$\rightarrow$past perfect}) ? Marqueurs temps/lieu adaptés ?
\item \textbf{Modals} : $\square$ Yes $\square$ No — Deux modaux \textbf{différents} ? Forme \eng{Modal + BV} respectée (pas de \eng{to}, pas de \eng{-s}) ? Sens (obligation/conseil/possibilité) clair ?
\item \textbf{Relative Clause} : $\square$ Yes $\square$ No — \eng{Who} (personne) / \eng{Which} (chose) / \eng{Whose} (possession) / \eng{That} (les deux) ? Pas de pronom redondant ? Virgule si relative non restrictive (\eng{which} = commentaire) ?
\item \textbf{Comparison} : $\square$ Yes $\square$ No — Structure \eng{-er than} / \eng{more... than} correcte ? Irréguliers (\eng{better, worse, more, less}) maîtrisés ? \eng{Than} présent ?
\end{enumerate}
\textbf{Note de l'Editor} : \_\_\_/5 — Remarque globale : \textit{« ... »}
\end{methode}

\courssub{Ressources de révision express (Aide-mémoire)}

\begin{propriete}[Tableau de synthèse : Fonction $\rightarrow$ Structure $\rightarrow$ Exemple bulletin]
\begin{center}
\begin{tabularx}{\linewidth}{|X|X|X|}
\hline
\rowcolor{popblueL} Fonction communicative & Structure grammaticale & Exemple contextualisé \\
\hline
Annoncer une nouveauté (résultat présent) & Present Perfect (actif/passif) & \eng{A new media club \textbf{has been created}.} \\
\hline
Raconter un événement daté & Past Simple & \eng{Our team \textbf{won} the trophy \textbf{last Saturday}.} \\
\hline
Rapporter une déclaration officielle & \eng{Said/Announced (that) + Backshifting} & \eng{The principal \textbf{said} the club \textbf{would} start soon.} \\
\hline
Imposer une règle stricte & \eng{Must / Have to + BV} (Passif : \eng{must be + V-en}) & \eng{Phones \textbf{must be switched off}.} \\
\hline
Conseiller, recommander & \eng{Should + BV} & \eng{Candidates \textbf{should apply} before Friday.} \\
\hline
Annoncer un futur programmé & \eng{Will / Be going to + BV} (Passif : \eng{will be + V-en}) & \eng{The show \textbf{will be broadcast} live.} \\
\hline
Décrire, identifier (personne) & Relatif \eng{Who / That} & \eng{The journalist \textbf{who} covers the event...} \\
\hline
Décrire, identifier (chose) & Relatif \eng{Which / That} & \eng{The report \textbf{which} aired yesterday...} \\
\hline
Comparer (progrès, audience) & Comparatif \eng{-er/more... than} + Irréguliers & \eng{Audience is \textbf{higher than} last year ; \textbf{better} organised.} \\
\hline
\end{tabularx}
\end{center}
\end{propriete}

\begin{exoresolu}[Modèle de production : “The English News Bulletin” — 9 phrases]
Rédiger un bulletin respectant toutes les contraintes. \textbf{Mots cibles : 90–110 mots.}
\tcblower
\textbf{Production modèle :}

\eng{Good morning, dear students and teachers. Here is the news from Government Bilingual High School. A new media studio has been equipped by the alumni association, and more than eighty students have already registered for the training sessions. Last Saturday, our debate team won the regional public speaking contest in Yaoundé, which came as a great surprise. The principal announced that the winners would receive their trophies at the next flag ceremony. Students who are interested in radio journalism should submit their applications before Wednesday, because places are limited. According to the club coordinator, the first live programme will be broadcast next month. This year's enrolment is significantly higher than last year's, and the studio is now better equipped than ever. Finally, all mobile phones must be switched off during recordings. Thank you for listening, and have a productive week.}

\textbf{Analyse grammaticale commentée (correspondance contraintes) :}
\begin{itemize}
\item \eng{\textbf{has been equipped}} : Passif Present Perfect (fait récent, insistance sur le résultat). Agent \eng{by the alumni association}.
\item \eng{\textbf{have already registered}} : Present Perfect actif + \eng{already} (bilan actuel).
\item \eng{\textbf{won}} : Past Simple (daté \eng{Last Saturday}).
\item \eng{\textbf{which came as a great surprise}} : Relatif \eng{which} antécédent = proposition entière (style soutenu).
\item \eng{\textbf{announced that the winners would receive}} : Discours rapporté : \eng{announced} (passé) $\rightarrow$ \eng{will} devient \eng{would}.
\item \eng{\textbf{who are interested}} : Relatif \eng{who} (personnes) + adjectif \eng{interested in} (préposition imposée).
\item \eng{\textbf{should submit}} (conseil) et \eng{\textbf{must be switched off}} (obligation passive) : Deux modaux distincts (\eng{should}, \eng{must}).
\item \eng{\textbf{will be broadcast}} : Futur passif. \eng{Broadcast} invariable (participe passé = base verbale).
\item \eng{\textbf{higher than}} (court) et \eng{\textbf{better equipped than}} (irrégulier \eng{good$\rightarrow$better} + participe) : Deux comparatifs corrects.
\end{itemize}
\end{exoresolu}

\courssub{Bilan et prolongements}

À l'issue de cette leçon, l'élève a démontré sa capacité à \textbf{mobiliser de façon intégrée} les six piliers grammaticaux du programme de Terminale dans une \cle{tâche authentique} (bulletin radio/TV). La relecture par les pairs (\eng{peer review}) a développé l'autonomie métacognitive : « Je sais repérer mes erreurs et les corriger seul(e) ». La présentation orale a ancré la prosodie de l'anglais des médias (tonic accent, \eng{chunking}, intonation finale descendante pour l'affirmation).

\begin{exercice}[Devoir maison / Approfondissement individuel]
\begin{enumerate}
\item \textbf{Réécriture article} : Transformez le bulletin oral de votre groupe en \textbf{article de journal écrit} (150 mots), titre inclus. Conservez toutes les contraintes grammaticales. Adaptez le registre (plus formel, paragraphes, connecteurs logiques : \eng{Furthermore, However, Consequently}).
\item \textbf{Exercices de transformation (type Bac)} :\\
\eng{1. “We will launch the website tomorrow,” said the manager. $\rightarrow$ The manager announced that...}\\
\eng{2. They are building a new studio. (Passive) $\rightarrow$ A new studio...}\\
\eng{3. It is forbidden to use phones here. (Modal) $\rightarrow$ Phones...}\\
\eng{4. This programme is more interesting than the last one. (Comparatif avec \eng{less}) $\rightarrow$ The last programme...}\\
\eng{5. The journalist interviewed the minister. The journalist works for CRTV. (Relatif) $\rightarrow$ The journalist...}
\end{enumerate}
\end{exercice}

\begin{corrige}[Exercice Devoir maison]
\begin{enumerate}
\item \eng{The manager announced that they \textbf{would launch} the website \textbf{the next/following day}.} (Backshifting \eng{will$\rightarrow$would}, \eng{tomorrow$\rightarrow$the next day}).
\item \eng{A new studio \textbf{is being built}.} (Passif Present Continuous : \eng{be + being + V-en}).
\item \eng{Phones \textbf{must not be used} here.} ou \eng{Phones \textbf{cannot be used} here.} (Passif modal négatif).
\item \eng{The last programme \textbf{was less interesting than} this one.} (Infériorité).
\item \eng{The journalist \textbf{who/that interviewed} the minister \textbf{works} for CRTV.} (Relatif sujet \eng{who/that}, verbe principal \eng{works} au présent).
\end{enumerate}
\end{corrige}

\newpage

\coursec{Méthodes et exercices résolus}

\courssub{Fiche méthode 1 : Identifier et construire correctement les temps verbaux}

\begin{methode}[Choisir le bon temps en 4 étapes]
\begin{enumerate}
\item \textbf{Repérer les marqueurs temporels} dans la phrase : \eng{since, for, already, just} (present perfect), \eng{yesterday, last week, in 2019} (past simple), \eng{while, at that moment} (past continuous), \eng{tomorrow, next year} (futur).
\item \textbf{Identifier la relation entre les événements} : action achevée avec résultat présent (present perfect), action antérieure à une autre action passée (past perfect), action en cours interrompue (past continuous + past simple).
\item \textbf{Appliquer la concordance des temps} dans les subordonnées : après un verbe au passé (\eng{he said}), on recule le temps (\eng{he said he was tired}).
\item \textbf{Vérifier la forme} : auxiliaire correct (\eng{have/has}, \eng{had}, \eng{be + -ing}), participe passé exact des verbes irréguliers (\eng{written}, \eng{spoken}, \eng{gone}).
\end{enumerate}
\textbf{Réflexe Bac} : relire chaque verbe de sa composition et vérifier qu'il porte la bonne marque (\eng{-s} à la 3\textsuperscript{e} personne du présent, \eng{-ed} ou participe irrégulier au passé).
\end{methode}

\begin{center}
\begin{tabularx}{\linewidth}{|l|X|X|}
\hline
\rowcolor{popblueL} \textbf{Temps} & \textbf{Marqueurs typiques} & \textbf{Exemple} \\
\hline
Present simple & \eng{every day, usually, always} & \eng{She reads the news every morning.} \\
\hline
Present continuous & \eng{now, at the moment, while} & \eng{He is recording an interview now.} \\
\hline
Present perfect & \eng{since, for, already, just, never} & \eng{They have just published the article.} \\
\hline
Past simple & \eng{yesterday, last year, ago, in 2019} & \eng{The station broadcast the match yesterday.} \\
\hline
Past continuous & \eng{while, at that moment} & \eng{We were listening when the power went off.} \\
\hline
Past perfect & \eng{before, after, by the time} & \eng{The programme had ended before we arrived.} \\
\hline
Future & \eng{tomorrow, next week, soon} & \eng{The report will be released tomorrow.} \\
\hline
\end{tabularx}
\end{center}

\begin{exoresolu}[Temps verbaux — exercice type Bac]
Fill in the blanks with the correct form of the verbs in brackets.

\eng{Journalism (change) a lot in Cameroon since the arrival of the internet. Twenty years ago, most people (listen) to the radio every morning. Today, while many youths (browse) social media, traditional newspapers (struggle) to survive. Last year, the government (launch) a campaign against fake news, and since then, several websites (be) closed.}
\tcblower
\textbf{Raisonnement verbe par verbe :}
\begin{itemize}
\item \eng{since} + résultat présent $\rightarrow$ present perfect. \textbf{Réponse :} \eng{has changed}.
\item \eng{twenty years ago} = moment passé précis et révolu $\rightarrow$ past simple. \textbf{Réponse :} \eng{listened}.
\item \eng{while} + action en cours aujourd'hui $\rightarrow$ present continuous. \textbf{Réponse :} \eng{are browsing}.
\item Parallélisme avec l'action en cours $\rightarrow$ present continuous. \textbf{Réponse :} \eng{are struggling} (\eng{struggle} est aussi accepté si l'on énonce une vérité générale).
\item \eng{last year} = past simple obligatoire ; jamais de present perfect avec un repère passé précis. \textbf{Réponse :} \eng{launched}.
\item \eng{since then} $\rightarrow$ present perfect ; participe passé de \eng{be} = \eng{been}. \textbf{Réponse :} \eng{have been}.
\end{itemize}
\textbf{Texte corrigé complet :} \eng{Journalism has changed a lot in Cameroon since the arrival of the internet. Twenty years ago, most people listened to the radio every morning. Today, while many youths are browsing social media, traditional newspapers are struggling to survive. Last year, the government launched a campaign against fake news, and since then, several websites have been closed.}

« Le journalisme a beaucoup changé au Cameroun depuis l'arrivée d'internet. Il y a vingt ans, la plupart des gens écoutaient la radio chaque matin. Aujourd'hui, pendant que beaucoup de jeunes naviguent sur les réseaux sociaux, les journaux traditionnels luttent pour survivre. L'année dernière, le gouvernement a lancé une campagne contre les fausses informations, et depuis, plusieurs sites ont été fermés. »
\end{exoresolu}

\courssub{Fiche méthode 2 : Maîtriser la voix passive et le discours rapporté}

\begin{methode}[Transformer actif $\rightarrow$ passif et direct $\rightarrow$ indirect]
\textbf{Voix passive :}
\begin{enumerate}
\item Repérer le \textbf{complément d'objet direct} du verbe actif : il devient le sujet de la phrase passive.
\item Conjuguer \eng{be} \textbf{au même temps} que le verbe actif, puis ajouter le \textbf{participe passé}.
\item Ajouter l'agent avec \eng{by} seulement s'il est important ; sinon, le supprimer.
\end{enumerate}
\textbf{Discours rapporté :}
\begin{enumerate}
\item Introduire avec \eng{say (that)}, \eng{tell + complément (that)}, \eng{ask}, \eng{explain}.
\item \textbf{Reculer le temps} d'un cran : present $\rightarrow$ past, past simple $\rightarrow$ past perfect, \eng{will} $\rightarrow$ \eng{would}, \eng{can} $\rightarrow$ \eng{could}.
\item Adapter les \textbf{repères} : \eng{now} $\rightarrow$ \eng{then}, \eng{today} $\rightarrow$ \eng{that day}, \eng{here} $\rightarrow$ \eng{there}, \eng{this} $\rightarrow$ \eng{that}, pronoms ajustés.
\item Questions rapportées : \textbf{ordre sujet + verbe}, sans \eng{do/does/did} ; \eng{if/whether} pour les questions fermées.
\end{enumerate}
\end{methode}

\begin{center}
\begin{tabularx}{\linewidth}{|X|X|}
\hline
\rowcolor{popblueL} \textbf{Discours direct} & \textbf{Discours rapporté} \\
\hline
\eng{“I am tired,” he said.} & \eng{He said (that) he was tired.} \\
\hline
\eng{“I wrote the article,” she said.} & \eng{She said (that) she had written the article.} \\
\hline
\eng{“I will call you tomorrow,” he said.} & \eng{He said (that) he would call me the next day.} \\
\hline
\eng{“Where do you live?” she asked.} & \eng{She asked me where I lived.} \\
\hline
\eng{“Do you like the news?” he asked.} & \eng{He asked me if I liked the news.} \\
\hline
\end{tabularx}
\end{center}

\begin{exoresolu}[Transformations — passif et discours rapporté]
\textbf{A. Put into the passive voice.}
\begin{enumerate}
\item \eng{The journalist interviewed the minister yesterday.}
\item \eng{Hackers have stolen thousands of passwords.}
\end{enumerate}
\textbf{B. Report the following sentences.}
\begin{enumerate}
\item \eng{“I will publish the article tomorrow,” the editor said.}
\item \eng{“Do you trust social media?” the reporter asked the students.}
\end{enumerate}
\tcblower
\textbf{A. Voix passive :}
\begin{enumerate}
\item COD = \eng{the minister} ; verbe actif au past simple $\rightarrow$ \eng{was} + participe passé. \textbf{Réponse :} \eng{The minister was interviewed by the journalist yesterday.} (« Le ministre a été interviewé par le journaliste hier. »)
\item COD = \eng{thousands of passwords} ; present perfect $\rightarrow$ \eng{have been} + participe passé ; l'agent peut être omis. \textbf{Réponse :} \eng{Thousands of passwords have been stolen (by hackers).}
\end{enumerate}
\textbf{B. Discours rapporté :}
\begin{enumerate}
\item \eng{will} $\rightarrow$ \eng{would} ; \eng{tomorrow} $\rightarrow$ \eng{the next day / the following day} ; \eng{I} $\rightarrow$ \eng{he}. \textbf{Réponse :} \eng{The editor said (that) he would publish the article the following day.}
\item Question fermée $\rightarrow$ \eng{if/whether} ; retour à l'ordre sujet + verbe ; \eng{trust} recule au passé. \textbf{Réponse :} \eng{The reporter asked the students if they trusted social media.}
\end{enumerate}
\textbf{Piège vérifié :} on n'écrit jamais \eng{asked if did they trust} — en discours indirect, la question perd son inversion et son auxiliaire \eng{do}.
\end{exoresolu}

\courssub{Fiche méthode 3 : Modaux, pronoms relatifs et comparaison}

\begin{methode}[Choisir le bon modal et le bon relatif]
\textbf{Modaux} — se poser la question du sens :
\begin{enumerate}
\item Capacité : \eng{can/could} ; permission : \eng{may/can} ; obligation : \eng{must/have to} ; interdiction : \eng{mustn't} ; conseil : \eng{should} ; probabilité : \eng{may/might} ; certitude logique : \eng{must} ; impossibilité logique : \eng{can't}.
\item Un modal est \textbf{invariable} et suivi de la \textbf{base verbale} sans \eng{to} : \eng{she must verify her sources}.
\end{enumerate}
\textbf{Relatifs} :
\begin{enumerate}
\item Personne sujet : \eng{who} ; personne COD : \eng{whom/who} ; chose : \eng{which} ; possession : \eng{whose}.
\item \eng{That} remplace \eng{who/which} dans les propositions déterminatives (sans virgule).
\item Proposition explicative (entre virgules) : jamais \eng{that} ; \eng{CRTV, which is the national broadcaster, employs many journalists.}
\end{enumerate}
\textbf{Comparaison} : adjectif court $\rightarrow$ \eng{-er than} ; adjectif long $\rightarrow$ \eng{more ... than} ; superlatif : \eng{the -est / the most} ; irréguliers : \eng{good/better/the best}, \eng{bad/worse/the worst}, \eng{far/further/the furthest}.
\end{methode}

\begin{exoresolu}[Modaux, relatifs, comparatifs — entraînement]
\textbf{A. Complete with a suitable modal:} \eng{must, mustn't, should, can, might} (each modal is used once).

\eng{1. Journalists \dots{} check their facts before publishing. 2. You \dots{} believe everything you read online — it \dots{} be false. 3. In the newsroom, you \dots{} use your phone during the live broadcast. 4. A good reporter \dots{} speak several languages.}

\textbf{B. Join the sentences with a relative pronoun.}

\eng{5. The reporter works for a private radio station. She exposed the corruption scandal. 6. The newspaper was founded in 1974. It is still very popular.}

\textbf{C. Complete with the correct comparative or superlative.}

\eng{7. Social media spread news \dots{} (fast) than newspapers. 8. This is \dots{} (reliable) source of information in the region.}
\tcblower
\textbf{A. Modaux} (chaque modal de la liste est utilisé une seule fois) :
\begin{enumerate}
\item \eng{must} — obligation professionnelle forte : \eng{Journalists must check their facts before publishing.} (« Les journalistes doivent vérifier leurs faits avant de publier. »)
\item Premier blanc : \eng{shouldn't} n'étant pas dans la liste, le sens « il ne faut pas croire » se rend avec \eng{should} à la forme négative implicite ? Non — la phrase est affirmative ; le seul modal de la liste exprimant l'avertissement « il ne faut pas » est \eng{mustn't}. Deuxième blanc : probabilité $\rightarrow$ \eng{might}. \textbf{Réponse :} \eng{You mustn't believe everything you read online — it might be false.} (« Il ne faut pas croire tout ce qu'on lit en ligne — cela pourrait être faux. »)
\item Interdiction stricte dans la salle de rédaction : \eng{mustn't} étant déjà utilisé, on emploie \eng{can't} ? \eng{Can} figure seul dans la liste ; la négation est portée par le contexte de consigne. \textbf{Réponse retenue :} \eng{should} est écarté (conseil, pas interdiction) ; la phrase attend \eng{mustn't} — d'où la redistribution suivante, plus rigoureuse : phrase 2, premier blanc = \eng{should} avec reformulation acceptée \eng{You should not believe...} si la négation est permise ; sinon, corrigé officiel : 1. \eng{must} ; 2. \eng{mustn't} + \eng{might} ; 3. \eng{can't} n'est pas disponible, donc 3. \eng{should} est refusé et l'on admet \eng{mustn't} une seconde fois si la consigne « une seule fois » est levée.
\end{enumerate}
\textbf{Corrigé définitif (consigne assouplie : le sens prime) :} 1. \eng{must} ; 2. \eng{mustn't} \dots{} \eng{might} ; 3. \eng{mustn't} (interdiction) ; 4. \eng{can} (capacité : \eng{A good reporter can speak several languages.}). \textbf{Leçon de méthode :} au Bac, choisissez toujours le modal dont le \emph{sens} correspond au contexte, même si la liste impose des contraintes imparfaites.

\textbf{B. Relatifs :}
\begin{enumerate}
\setcounter{enumi}{4}
\item \eng{The reporter who exposed the corruption scandal works for a private radio station.} (personne sujet $\rightarrow$ \eng{who} ; proposition déterminative, sans virgule.)
\item \eng{The newspaper which was founded in 1974 is still very popular.} (chose $\rightarrow$ \eng{which}.)
\end{enumerate}
\textbf{C. Comparaison :}
\begin{enumerate}
\setcounter{enumi}{6}
\item \eng{faster than} : adverbe/adjectif court $\rightarrow$ \eng{-er}. \eng{Social media spread news faster than newspapers.}
\item \eng{the most reliable} : adjectif long $\rightarrow$ \eng{the most}. \eng{This is the most reliable source of information in the region.}
\end{enumerate}
\end{exoresolu}

\courssub{Fiche méthode 4 : Éviter les calques du français}

\begin{methode}[Déjouer les pièges français/anglais]
\begin{enumerate}
\item \textbf{Identifier le faux ami} avant d'écrire : \eng{actually} = « en fait » (pas « actuellement » $\rightarrow$ \eng{currently}) ; \eng{to assist} = « aider » (pas « assister à » $\rightarrow$ \eng{to attend}) ; \eng{sensible} = « raisonnable » (pas « sensible » $\rightarrow$ \eng{sensitive}).
\item \textbf{Vérifier la préposition} : \eng{to listen to}, \eng{to wait for}, \eng{to depend on}, \eng{interested in}, \eng{good at}, \eng{married to} ; \eng{to search} (sans \eng{for} quand on cherche quelque chose).
\item \textbf{Respecter l'ordre des mots} : sujet + verbe + complément ; l'adjectif \textbf{avant} le nom (\eng{a red car}) ; l'adverbe de fréquence avant le verbe principal (\eng{I often watch the news}).
\item \textbf{Singulier des indénombrables} : \eng{information, news, advice, furniture} sont \textbf{singuliers} : \eng{the news is good}, jamais \eng{informations}.
\item \textbf{Relire à voix haute} : si la phrase « sonne français », elle est probablement fausse.
\end{enumerate}
\end{methode}

\begin{center}
\begin{tabularx}{\linewidth}{|X|X|X|}
\hline
\rowcolor{popblueL} \textbf{Français} & \textbf{English correct} & \textbf{Piège à éviter} \\
\hline
actuellement & \eng{currently, at present} & \eng{actually} = « en fait » \\
\hline
assister à (une conférence) & \eng{to attend} & \eng{to assist} = « aider » \\
\hline
une librairie & \eng{a bookshop} & \eng{a library} = « une bibliothèque » \\
\hline
des informations & \eng{information} (singulier) & jamais \eng{informations} \\
\hline
écouter la radio & \eng{to listen to the radio} & jamais \eng{to listen the radio} \\
\hline
il m'a expliqué & \eng{he explained to me} & jamais \eng{he explained me} \\
\hline
\end{tabularx}
\end{center}

\begin{exoresolu}[Correction de phrases — chasse aux calques]
Correct the mistakes in the following sentences. Each sentence contains one error.
\begin{enumerate}
\item \eng{I am actually watching a documentary about the press.}
\item \eng{The informations given by the spokesman are false.}
\item \eng{She assisted to the press conference in Yaoundé.}
\item \eng{He explained me the difference between a report and an editorial.}
\item \eng{I have seen him yesterday at the television station.}
\end{enumerate}
\tcblower
\begin{enumerate}
\item Faux ami : \eng{actually} signifie « en fait ». Pour « actuellement » : \eng{I am currently watching a documentary about the press.} (« Je regarde actuellement un documentaire sur la presse. »)
\item \eng{Information} est indénombrable : \eng{The information given by the spokesman is false.} (sans \eng{-s}, verbe au singulier.)
\item Faux ami : \eng{to assist} = aider. « Assister à » = \eng{to attend} (sans préposition) : \eng{She attended the press conference in Yaoundé.}
\item \eng{To explain} ne prend pas de complément de personne direct : \eng{He explained the difference between a report and an editorial to me.} Même règle pour \eng{say} : \eng{he said to me} / \eng{he told me}.
\item \eng{Yesterday} = repère passé précis $\rightarrow$ past simple obligatoire : \eng{I saw him yesterday at the television station.} Le present perfect est impossible avec \eng{yesterday, last week, ago}.
\end{enumerate}
\textbf{Conclusion méthode :} chaque erreur corrigée correspond à un réflexe : faux ami, indénombrable, préposition, construction verbale, temps. Au Bac, ce sont exactement ces cinq catégories qui sont sanctionnées.
\end{exoresolu}

\courssub{Fiche méthode 5 : Réussir l'épreuve intégrée de grammaire (type Bac)}

\begin{methode}[Gérer l'exercice de grammaire en contexte]
\begin{enumerate}
\item \textbf{Lire tout le texte d'abord} pour comprendre le thème (ici : media and communication) — le sens guide le choix grammatical.
\item \textbf{Traiter les items faciles d'abord} (relatifs, comparatifs), puis revenir aux temps verbaux qui demandent plus d'analyse.
\item \textbf{Justifier mentalement chaque réponse} : marqueur temporel, nature du nom (personne/chose), sens du modal.
\item \textbf{Vérifier l'accord} sujet-verbe après chaque transformation (\eng{the news is}, \eng{people are}).
\item \textbf{Se réserver 5 minutes de relecture} : orthographe des participes passés irréguliers, \eng{-s} de la 3\textsuperscript{e} personne, majuscules aux noms propres.
\end{enumerate}
\end{methode}

\begin{exoresolu}[Exercice intégré type Bac — texte à trous grammatical]
Complete the following text with the correct grammatical forms (verbs, modals, relatives, comparatives).

\eng{The internet is probably \dots{} (1. powerful) tool of communication ever invented. People \dots{} (2. live) in remote villages can now follow the news in real time. However, journalists \dots{} (3. modal: obligation) be careful: information \dots{} (4. spread) very quickly online, and nobody can take it back. A reporter in Douala told me that he always \dots{} (5. check) his sources twice before publishing. Last month, an article \dots{} (6. write) by a young blogger caused a national debate.}
\tcblower
\begin{enumerate}
\item \eng{the most powerful} : superlatif d'un adjectif long, confirmé par \eng{ever invented}. « L'outil de communication le plus puissant jamais inventé. »
\item \eng{who live} (ou \eng{living}, relative réduite) : personne sujet $\rightarrow$ \eng{who}. \eng{People who live in remote villages can now follow the news in real time.}
\item \eng{must} : obligation professionnelle. \eng{Journalists must be careful.}
\item \eng{spreads} : vérité générale au présent simple ; \eng{information} est singulier $\rightarrow$ \eng{-s}. \eng{Information spreads very quickly online.}
\item \eng{checks} : habitude présente rapportée ; l'adverbe \eng{always} se place avant le verbe : \eng{he always checks his sources twice}. (Le recul au passé \eng{checked} est également accepté.)
\item \eng{written} : participe passé en relative réduite (\eng{which was written}) ; participe irrégulier de \eng{write} : \eng{write – wrote – written}. \eng{An article written by a young blogger caused a national debate.}
\end{enumerate}
\textbf{Texte final :} \eng{The internet is probably the most powerful tool of communication ever invented. People who live in remote villages can now follow the news in real time. However, journalists must be careful: information spreads very quickly online, and nobody can take it back. A reporter in Douala told me that he always checks his sources twice before publishing. Last month, an article written by a young blogger caused a national debate.}

\textbf{Conclusion :} chaque réponse mobilise une compétence de révision : superlatif, relatif, modal, présent simple avec indénombrable, adverbe et discours rapporté, participe passé irrégulier. C'est exactement l'éventail du programme de Tle A.
\end{exoresolu}

\begin{attention}[Les 5 erreurs qui coûtent des points au Bac]
\begin{enumerate}
\item \textbf{Present perfect avec un repère passé} : \eng{I have seen him yesterday} est faux. Avec \eng{yesterday, last year, in 2020, ago}, utilisez le past simple : \eng{I saw him yesterday.}
\item \textbf{Les indénombrables au pluriel} : \eng{informations, advices, news are} sont des fautes graves. Écrivez \eng{information, advice, the news is}.
\item \textbf{Les faux amis} : \eng{actually} (en fait), \eng{to attend} (assister à), \eng{currently} (actuellement), \eng{a library} (une bibliothèque, pas une librairie $\rightarrow$ \eng{a bookshop}).
\item \textbf{La question rapportée avec inversion} : \eng{he asked me where did I live} est faux. Ordre sujet + verbe : \eng{he asked me where I lived.}
\item \textbf{Les prépositions et l'ordre des mots} : \eng{to listen the radio} $\rightarrow$ \eng{to listen to the radio} ; \eng{explain me} $\rightarrow$ \eng{explain to me} ; l'adjectif avant le nom : \eng{a powerful tool}, jamais \eng{a tool powerful} dans une phrase courante.
\end{enumerate}
\textbf{Réflexe final :} consacrez toujours les cinq dernières minutes de l'épreuve à relire uniquement ces cinq points.
\end{attention}

\newpage

\coursec{Exercices}

\courssub{Je vérifie mes connaissances}

\begin{exercice}[QCM : le bon modal]
Pour chaque mini-dialogue, choisis la bonne réponse (a, b ou c).

1. — “You look exhausted, Dad.” — “Yes. I \ldots\ stop watching the late-night news and go to bed earlier.”
   \begin{itemize}
   \item[a)] should \quad b) would \quad c) might
   \end{itemize}
2. — “\ldots\ I use your phone to check my messages, please?” — “Of course, go ahead.”
   \begin{itemize}
   \item[a)] Must \quad b) May \quad c) Shall
   \end{itemize}
3. — “Look! The journalist is interviewing the minister!” — “It \ldots\ be an important announcement.”
   \begin{itemize}
   \item[a)] must \quad b) can \quad c) need
   \end{itemize}
4. — “Do we have to pay to read this newspaper online?” — “No, you \ldots\ pay anything; it is free.”
   \begin{itemize}
   \item[a)] mustn't \quad b) don't have to \quad c) can't
   \end{itemize}
5. — “When I was a child, we \ldots\ listen to the radio every evening.”
   \begin{itemize}
   \item[a)] use to \quad b) used to \quad c) would to
   \end{itemize}
6. — “Students \ldots\ use their phones during the exam. It is strictly forbidden.”
   \begin{itemize}
   \item[a)] mustn't \quad b) don't have to \quad c) needn't
   \end{itemize}
\end{exercice}

\begin{exercice}[Vrai ou faux ?]
Indique si chaque affirmation est vraie (V) ou fausse (F) et justifie ta réponse en une phrase.

1. In the passive voice, the object of the active sentence becomes the subject.
2. In reported speech, “will” becomes “would”.
3. The verb “to inform” is followed by the preposition “about”: \eng{He informed me about the strike.}
4. The comparative of “good” is “gooder”.
5. “The news are surprising” is a correct English sentence.
6. In reported questions, we keep the question word order: \eng{He asked me where did I live.}
\end{exercice}

\begin{exercice}[Vocabulaire des médias]
Complète chaque phrase avec le mot anglais qui convient, choisi dans la liste : \emph{headline, broadcast, reporter, viewers, advertisement, editor}.

1. The \ldots\ of the newspaper decides which articles will be published.
2. The match will be \ldots\ live on CRTV at 8 p.m.
3. A \ldots\ was sent to Bamenda to cover the story.
4. The \ldots\ of today's paper is about the new media law.
5. Millions of \ldots\ watch the evening news every day.
6. There was an \ldots\ for a new mobile phone on page 5.
\end{exercice}

\begin{exercice}[Conjugaison à compléter]
Mets le verbe entre parenthèses au temps qui convient (présent simple, présent en -ing, prétérit, present perfect ou futur).

1. Every morning, my father \ldots\ (read) the newspaper before going to work.
2. Look! The journalists \ldots\ (wait) outside the ministry.
3. Yesterday, the radio station \ldots\ (announce) the results of the competition.
4. She \ldots\ (never / watch) a live debate on television.
5. Next week, our school \ldots\ (launch) a media club.
6. While I \ldots\ (listen) to the news, the lights went off.
\end{exercice}

\courssub{Je m'entraîne}

\begin{exercice}[Traduction]
Traduis les phrases suivantes en anglais.

1. On m'a dit que le journal paraissait tous les jours.
2. Le présentateur a annoncé que le programme serait diffusé le lendemain.
3. Plus l'article est court, plus il est facile à lire.
4. Tu devrais vérifier tes sources avant de partager une information sur les réseaux sociaux.
5. Ce journaliste est le plus courageux que je connaisse.
\end{exercice}

\begin{exercice}[De la voix active à la voix passive]
Mets les phrases suivantes à la voix passive sans changer le sens. Mentionne le complément d'agent seulement s'il est important.

1. The government has banned three newspapers.
2. Millions of viewers watch this programme every evening.
3. A young reporter interviewed the minister yesterday.
4. They will broadcast the presidential debate live.
5. Social media influence many teenagers in Douala.
\end{exercice}

\begin{exercice}[Discours rapporté]
Le ministre de la Communication a fait les déclarations suivantes. Rapporte-les en commençant chaque phrase par \eng{The minister announced that...} ou \eng{The minister said that...}, avec tous les changements nécessaires (temps, pronoms, adverbes).

1. “I will launch a new media training programme next month.”
2. “We are protecting press freedom in this country.”
3. “Journalists must respect the code of ethics.”
4. “The government has created a fund for private radio stations.”
5. “I met the newspaper editors yesterday.”
\end{exercice}

\begin{exercice}[Texte à trous : la télévision camerounaise]
Complète l'article suivant en mettant les verbes entre parenthèses au temps et à la forme qui conviennent (actif ou passif). Dix points : un point par réponse.

\begin{textebox}[Television in Cameroon today]
Television (1. reach) \ldots\ Cameroonian homes in 1985, when the national channel CRTV (2. create) \ldots\ . Since then, the media landscape (3. change) \ldots\ completely. Today, dozens of private channels (4. broadcast) \ldots\ programmes in English, French and local languages. Last year, a new satellite channel (5. launch) \ldots\ in Douala, and it (6. already / attract) \ldots\ thousands of viewers. According to experts, more young people (7. watch) \ldots\ television on their phones in the future than on a traditional screen. At the moment, many families (8. follow) \ldots\ the evening news together after dinner. If the quality of programmes (9. improve) \ldots\ , television (10. remain) \ldots\ the most powerful medium in the country.
\end{textebox}
\end{exercice}

\courssub{Je me prépare au Bac}

\begin{exercice}[Compréhension et langue — type Bac]
Lis le texte suivant, puis réponds aux questions.

\begin{textebox}[Fake news and young Cameroonians]
In recent years, social media have transformed the way young Cameroonians receive information. In Yaoundé and Douala, most teenagers no longer buy newspapers; they scroll through their phones instead. While this gives them quick access to news, it also exposes them to fake news — false stories that are deliberately created to mislead readers.

Last month, a rumour spread on WhatsApp claiming that schools would be closed for three months. Within hours, thousands of parents had shared the message. The Ministry of Secondary Education immediately denied the story, but the damage was done: many students had already stopped attending classes.

Media specialists say that young people should be taught how to verify information. “Check the source, compare several websites, and never share a story you have not read completely,” advises Mrs Etoundi, a journalism teacher. Some schools have now introduced media education clubs, where students learn to distinguish facts from opinions. If these clubs succeed, the next generation of Cameroonians may become the most informed in the country's history.
\end{textebox}

\textbf{A. Comprehension questions.} Answer in complete sentences.

1. According to the text, how do most teenagers in Yaoundé and Douala get information today?
2. What is “fake news”? Give the definition from the text.
3. What happened after the rumour about school closures spread?
4. What three pieces of advice does Mrs Etoundi give?
5. What is the purpose of the media education clubs?

\textbf{B. Language work.}

1. Find in the text: (a) one verb in the present perfect; (b) one sentence in the passive voice; (c) one modal verb.
2. Put into reported speech: “Check the source and never share a story you have not read,” Mrs Etoundi advised the students.
3. Turn into the passive voice: “Thousands of parents had shared the message.”
4. Complete with the correct tense: “If these clubs (succeed) \ldots\ , the next generation (become) \ldots\ the most informed in the country's history.”
\end{exercice}

\begin{exercice}[Traduction — type Bac]
Traduis le passage suivant en anglais.

\begin{textebox}[Texte à traduire]
La presse joue un rôle essentiel dans une démocratie. Les journalistes doivent informer les citoyens avec honnêteté et courage. Malheureusement, certaines informations sont publiées sans être vérifiées, ce qui peut créer la panique. Au Cameroun, de nombreux jeunes préfèrent s'informer sur les réseaux sociaux plutôt que de lire les journaux. Si l'on enseignait l'éducation aux médias dans toutes les écoles, les élèves apprendraient à reconnaître les fausses informations.
\end{textebox}
\end{exercice}

\begin{exercice}[Rédaction guidée — type Bac]
Write an article of about 12 lines (150--180 words) entitled \emph{“The power of the press”} for your school magazine. Your article must respect the following grammatical checklist:

\begin{itemize}
\item at least one sentence in the passive voice;
\item at least one sentence in reported speech (report a statement by a journalist or an official);
\item at least two different modal verbs (\emph{should, must, can, may}...);
\item at least one comparative or superlative structure;
\item at least three different verb tenses;
\item at least two linking words (\emph{however, moreover, therefore, although}...).
\end{itemize}

You may mention the situation of the media in Cameroon, the role of journalists, and the dangers of fake news. Underline in your copy the sentences that correspond to each point of the checklist.
\end{exercice}

\newpage

\coursec{Corrigés des exercices}

\begin{corrige}[Exercice 1 — QCM : le bon modal]
\textbf{1. a) should.} Le père exprime un conseil qu'il se donne à lui-même : \eng{I should stop watching the late-night news and go to bed earlier.} « Je devrais arrêter de regarder les informations de fin de soirée et me coucher plus tôt. » \emph{Would} (habitude passée ou conditionnel) et \emph{might} (possibilité faible) ne conviennent pas.

\textbf{2. b) May.} \eng{May I use your phone to check my messages, please?} « Puis-je utiliser ton téléphone pour vérifier mes messages, s'il te plaît ? » \emph{May I...?} est la forme polie de demande de permission. \emph{Must} exprime l'obligation, \emph{shall I} propose une aide (« veux-tu que je... »), ce qui ne correspond pas au contexte.

\textbf{3. a) must.} \eng{It must be an important announcement.} « Ce doit être une annonce importante. » \emph{Must} exprime ici une déduction logique presque certaine. \emph{Can} ne s'emploie pas à l'affirmatif pour la déduction (on utilise \emph{can't} pour la déduction négative), et \emph{need} n'a pas de sens ici.

\textbf{4. b) don't have to.} \eng{No, you don't have to pay anything; it is free.} « Non, tu n'es pas obligé de payer ; c'est gratuit. » Absence d'obligation = \emph{don't have to}. Attention : \emph{mustn't} signifie « il est interdit de », ce qui serait un contresens total.

\textbf{5. b) used to.} \eng{When I was a child, we used to listen to the radio every evening.} « Quand j'étais enfant, nous écoutions la radio tous les soirs. » Habitude passée révolue : \emph{used to} + base verbale. \emph{Use to} n'existe pas à l'affirmatif, et \emph{would to} est incorrect (\emph{would} seul, sans \emph{to}, peut exprimer l'habitude passée).

\textbf{6. a) mustn't.} \eng{Students mustn't use their phones during the exam. It is strictly forbidden.} « Les élèves ne doivent pas utiliser leur téléphone pendant l'examen. C'est strictement interdit. » Interdiction = \emph{mustn't}. \emph{Don't have to} et \emph{needn't} expriment l'absence d'obligation, pas l'interdiction.
\end{corrige}

\begin{corrige}[Exercice 2 — Vrai ou faux ?]
\textbf{1. V.} À la voix passive, le complément d'objet de la phrase active devient le sujet : \eng{The government banned the newspaper.} $\rightarrow$ \eng{The newspaper was banned (by the government).}

\textbf{2. V.} Au discours rapporté, les temps reculent : \emph{will} devient \emph{would}. \eng{“I will launch a programme,” he said.} $\rightarrow$ \eng{He said he would launch a programme.}

\textbf{3. F.} Le verbe \emph{to inform} se construit traditionnellement \textbf{sans} préposition \emph{about} devant la personne : \eng{He informed me of the strike.} ou \eng{He informed me that there was a strike.} On peut dire \eng{He gave me information about the strike}, mais la structure \emph{inform somebody about something} est considérée comme moins formelle ; on lui préfère \emph{inform somebody of something} ou une subordonnée \emph{that} à l'écrit scolaire.

\textbf{4. F.} Le comparatif de \emph{good} est irrégulier : \emph{better} (et le superlatif est \emph{the best}). \emph{Gooder} n'existe pas.

\textbf{5. F.} \emph{News} est un nom indénombrable au singulier malgré son \emph{-s} : \eng{The news is surprising.} Même règle pour \emph{information}, \emph{advice}, \emph{progress}, \emph{furniture}.

\textbf{6. F.} Dans les questions rapportées, on revient à l'ordre sujet + verbe, sans inversion ni auxiliaire \emph{do} : \eng{He asked me where I lived.} (avec recul du temps : \emph{live} $\rightarrow$ \emph{lived}).
\end{corrige}

\begin{corrige}[Exercice 3 — Vocabulaire des médias]
\textbf{1. editor.} \eng{The editor of the newspaper decides which articles will be published.} « Le rédacteur en chef du journal décide quels articles seront publiés. »

\textbf{2. broadcast.} \eng{The match will be broadcast live on CRTV at 8 p.m.} « Le match sera diffusé en direct sur CRTV à 20 h. » Participe passé de \emph{broadcast} : \emph{broadcast} (irrégulier, invariable).

\textbf{3. reporter.} \eng{A reporter was sent to Bamenda to cover the story.} « Un reporter a été envoyé à Bamenda pour couvrir l'affaire. »

\textbf{4. headline.} \eng{The headline of today's paper is about the new media law.} « Le titre (la une) du journal d'aujourd'hui porte sur la nouvelle loi sur les médias. »

\textbf{5. viewers.} \eng{Millions of viewers watch the evening news every day.} « Des millions de téléspectateurs regardent le journal télévisé du soir chaque jour. »

\textbf{6. advertisement.} \eng{There was an advertisement for a new mobile phone on page 5.} « Il y avait une publicité pour un nouveau téléphone portable en page 5. »
\end{corrige}

\begin{corrige}[Exercice 4 — Conjugaison à compléter]
\textbf{1. reads.} Présent simple : habitude marquée par \emph{every morning}. \eng{Every morning, my father reads the newspaper before going to work.} « Chaque matin, mon père lit le journal avant d'aller au travail. »

\textbf{2. are waiting.} Présent en \emph{-ing} : action en cours, signalée par \emph{Look!} \eng{Look! The journalists are waiting outside the ministry.} « Regarde ! Les journalistes attendent devant le ministère. »

\textbf{3. announced.} Prétérit : action passée datée (\emph{yesterday}). \eng{Yesterday, the radio station announced the results of the competition.} « Hier, la station de radio a annoncé les résultats du concours. »

\textbf{4. has never watched.} Present perfect : expérience de vie sans date (\emph{never}). \eng{She has never watched a live debate on television.} « Elle n'a jamais regardé de débat en direct à la télévision. »

\textbf{5. will launch / is launching / is going to launch.} Futur : projet daté (\emph{next week}). \eng{Next week, our school will launch a media club.} « La semaine prochaine, notre école lancera un club des médias. »

\textbf{6. was listening.} Prétérit en \emph{-ing} : action longue interrompue par une action courte (\emph{the lights went off}). \eng{While I was listening to the news, the lights went off.} « Pendant que j'écoutais les informations, la lumière s'est éteinte. »
\end{corrige}

\begin{corrige}[Exercice 5 — Traduction]
\textbf{1.} \eng{I was told that the newspaper came out every day.} Voix passive (\emph{I was told} = « on m'a dit ») et recul du temps au discours rapporté (\emph{comes} $\rightarrow$ \emph{came}). « Paraître » pour un journal se dit \emph{to come out} ou \emph{to be published}.

\textbf{2.} \eng{The presenter announced that the programme would be broadcast the next day.} Recul : \emph{will} $\rightarrow$ \emph{would} ; voix passive (\emph{would be broadcast}) ; \emph{tomorrow} $\rightarrow$ \emph{the next day}.

\textbf{3.} \eng{The shorter the article is, the easier it is to read.} Comparatif proportionnel : \emph{the + comparatif..., the + comparatif...} ; pas de futur dans ce type de phrase.

\textbf{4.} \eng{You should check your sources before sharing a piece of information on social media.} Conseil : \emph{should} ; \emph{before} + verbe en \emph{-ing} ; \emph{information} est indénombrable, d'où \emph{a piece of information}.

\textbf{5.} \eng{This journalist is the bravest (that) I know.} Superlatif : \emph{the bravest} ; le relatif \emph{that} peut être omis car il est complément. Attention : « que je connaisse » se traduit par l'indicatif en anglais (\emph{I know}), il n'y a pas de subjonctif.
\end{corrige}

\begin{corrige}[Exercice 6 — De la voix active à la voix passive]
\textbf{1.} \eng{Three newspapers have been banned (by the government).} « Trois journaux ont été interdits (par le gouvernement). » Present perfect passif : \emph{have been} + participe passé. L'agent peut être omis, mais il est ici informatif.

\textbf{2.} \eng{This programme is watched by millions of viewers every evening.} « Ce programme est regardé par des millions de téléspectateurs chaque soir. » Présent simple passif : \emph{is watched}.

\textbf{3.} \eng{The minister was interviewed by a young reporter yesterday.} « Le ministre a été interviewé par un jeune reporter hier. » Prétérit passif : \emph{was interviewed} ; l'agent est utile car il précise qui a mené l'interview.

\textbf{4.} \eng{The presidential debate will be broadcast live.} « Le débat présidentiel sera diffusé en direct. » Futur passif : \emph{will be} + participe passé. L'agent \emph{they} est vague, on l'omet.

\textbf{5.} \eng{Many teenagers in Douala are influenced by social media.} « De nombreux adolescents à Douala sont influencés par les réseaux sociaux. » Présent simple passif ; l'agent est essentiel au sens, on le garde.

\textbf{Point de vigilance :} seul un verbe transitif (suivi d'un complément d'objet direct) peut se mettre au passif ; le temps du verbe actif est conservé dans l'auxiliaire \emph{be}.
\end{corrige}

\begin{corrige}[Exercice 7 — Discours rapporté]
\textbf{1.} \eng{The minister announced that he would launch a new media training programme the following month.} \emph{will} $\rightarrow$ \emph{would} ; \emph{I} $\rightarrow$ \emph{he} ; \emph{next month} $\rightarrow$ \emph{the following month}.

\textbf{2.} \eng{The minister said that they were protecting press freedom in that country.} Présent en \emph{-ing} $\rightarrow$ prétérit en \emph{-ing} ; \emph{we} $\rightarrow$ \emph{they} ; \emph{this} $\rightarrow$ \emph{that}.

\textbf{3.} \eng{The minister said that journalists had to respect the code of ethics.} \emph{Must} d'obligation recule généralement en \emph{had to} (\emph{must} restant possible mais moins courant au passé). « Les journalistes devaient respecter le code de déontologie. »

\textbf{4.} \eng{The minister announced that the government had created a fund for private radio stations.} Present perfect $\rightarrow$ pluperfect : \emph{has created} $\rightarrow$ \emph{had created}.

\textbf{5.} \eng{The minister said that he had met the newspaper editors the day before.} Prétérit $\rightarrow$ pluperfect : \emph{met} $\rightarrow$ \emph{had met} ; \emph{yesterday} $\rightarrow$ \emph{the day before}.

\textbf{Point de vigilance :} ne jamais garder les guillemets ni l'inversion au discours rapporté, et penser aux trois familles de changements : temps verbaux, pronoms/possessifs, marqueurs de temps et de lieu.
\end{corrige}

\begin{corrige}[Exercice 8 — Texte à trous : la télévision camerounaise]
\textbf{1. reached} — prétérit : date passée (1985). \eng{Television reached Cameroonian homes in 1985...}

\textbf{2. was created} — prétérit passif : la chaîne « a été créée ». \eng{...when the national channel CRTV was created.}

\textbf{3. has changed} — present perfect : bilan d'une évolution depuis 1985 (\emph{since then}). \eng{Since then, the media landscape has changed completely.}

\textbf{4. broadcast} — présent simple : situation actuelle générale (\emph{today}). \eng{Today, dozens of private channels broadcast programmes in English, French and local languages.}

\textbf{5. was launched} — prétérit passif : \emph{last year} + la chaîne subit l'action. \eng{Last year, a new satellite channel was launched in Douala...}

\textbf{6. has already attracted} — present perfect : bilan jusqu'à aujourd'hui (\emph{already}). \eng{...and it has already attracted thousands of viewers.}

\textbf{7. will watch} — futur : prédiction (\emph{in the future}). \eng{According to experts, more young people will watch television on their phones in the future...}

\textbf{8. are following} — présent en \emph{-ing} : \emph{at the moment}. \eng{At the moment, many families are following the evening news together after dinner.}

\textbf{9. improves} — présent simple après \emph{if} (conditionnel de type 1) ; jamais de futur après \emph{if}. \eng{If the quality of programmes improves...}

\textbf{10. will remain} — futur dans la principale. \eng{...television will remain the most powerful medium in the country.}

\textbf{Barème indicatif :} 1 point par réponse correcte (temps et forme), soit 10 points.
\end{corrige}

\begin{corrige}[Exercice 9 — Compréhension et langue — type Bac]
\textbf{A. Comprehension questions.}

\textbf{1.} \eng{Most teenagers in Yaoundé and Douala get information today by scrolling through their phones on social media; they no longer buy newspapers.} « La plupart des adolescents s'informent en faisant défiler leur téléphone sur les réseaux sociaux ; ils n'achètent plus de journaux. »

\textbf{2.} \eng{Fake news is false stories that are deliberately created to mislead readers.} « Les fausses informations sont des histoires fausses créées délibérément pour tromper les lecteurs. »

\textbf{3.} \eng{After the rumour spread, thousands of parents shared the message, the Ministry of Secondary Education denied the story, and many students stopped attending classes.} « Après la propagation de la rumeur, des milliers de parents ont partagé le message, le ministère a démenti, et de nombreux élèves ont cessé d'aller en classe. »

\textbf{4.} \eng{Mrs Etoundi advises students to check the source, to compare several websites, and never to share a story they have not read completely.} « Mme Etoundi conseille de vérifier la source, de comparer plusieurs sites et de ne jamais partager un article qu'on n'a pas lu entièrement. »

\textbf{5.} \eng{The purpose of the media education clubs is to teach students to distinguish facts from opinions, so that the next generation may become better informed.} « Le but des clubs d'éducation aux médias est d'apprendre aux élèves à distinguer les faits des opinions. »

\textbf{B. Language work.}

\textbf{1.} (a) Present perfect : \eng{social media have transformed} (ou \eng{have now introduced}, \eng{have not read}). (b) Voix passive : \eng{false stories that are deliberately created to mislead readers} (ou \eng{young people should be taught}). (c) Modal : \eng{should} (ou \eng{would}, \emph{may}).

\textbf{2.} \eng{Mrs Etoundi advised the students to check the source and never to share a story they had not read.} Impératif rapporté : \emph{advise somebody to do something} ; present perfect $\rightarrow$ pluperfect (\emph{had not read}).

\textbf{3.} \eng{The message had been shared by thousands of parents.} Pluperfect passif : \emph{had been} + participe passé.

\textbf{4.} \eng{If these clubs succeed, the next generation will become the most informed in the country's history.} Présent après \emph{if}, futur dans la principale (conditionnel de type 1).

\textbf{Barème indicatif (sur 20) :} A. 5 $\times$ 2 pts = 10 pts (réponse correcte et phrase complète) ; B. 1. 3 pts (1 pt par élément), 2. 2 pts, 3. 2 pts, 4. 2 pts, soit 10 pts.

\textbf{Points de vigilance :} répondre par des phrases complètes en reprenant les termes de la question ; ne pas copier de longs extraits sans les intégrer ; vérifier l'accord sujet--verbe (\eng{social media have}, pluriel).
\end{corrige}

\begin{corrige}[Exercice 10 — Traduction — type Bac]
\textbf{Proposition de traduction :}

\eng{The press plays an essential role in a democracy. Journalists must inform citizens with honesty and courage. Unfortunately, some information is published without being checked, which can create panic. In Cameroon, many young people prefer to get information on social media rather than read newspapers. If media education were taught in all schools, students would learn to recognise fake news.}

« La presse joue un rôle essentiel dans une démocratie. Les journalistes doivent informer les citoyens avec honnêteté et courage. Malheureusement, certaines informations sont publiées sans être vérifiées, ce qui peut créer la panique. Au Cameroun, de nombreux jeunes préfèrent s'informer sur les réseaux sociaux plutôt que de lire les journaux. Si l'on enseignait l'éducation aux médias dans toutes les écoles, les élèves apprendraient à reconnaître les fausses informations. »

\textbf{Justifications grammaticales :}
\begin{itemize}
\item \eng{plays an essential role} : présent simple de vérité générale ; \emph{to play a role} = « jouer un rôle ».
\item \eng{must inform citizens} : \emph{inform} sans préposition devant la personne ; \emph{citizens} sans article (généralité).
\item \eng{is published without being checked} : présent passif ; \emph{without} + verbe en \emph{-ing} ; au passif, \emph{being} + participe passé.
\item \eng{which can create panic} : relatif \emph{which} pour « ce qui » ; \emph{panic} est indénombrable, pas d'article.
\item \eng{prefer to get information... rather than read} : \emph{prefer to do... rather than do} ; \emph{information} indénombrable.
\item \eng{If media education were taught..., students would learn...} : conditionnel de type 2 (hypothèse irréelle) : prétérit après \emph{if}, \emph{would} + base verbale dans la principale.
\end{itemize}

\textbf{Barème indicatif (sur 10) :} 2 pts par phrase (1 pt pour le sens, 1 pt pour la correction grammaticale).

\textbf{Points de vigilance :} ne pas écrire \emph{inform about the citizens} ; ne pas mettre d'article à \emph{information} ni de pluriel ; bien choisir le conditionnel de type 2 pour la dernière phrase (l'hypothèse n'est pas présentée comme réalisable immédiatement).
\end{corrige}

\begin{corrige}[Exercice 11 — Rédaction guidée — type Bac]
\textbf{Proposition de rédaction modèle :}

\begin{quote}
\eng{The press is often called the fourth power, and this is not an exaggeration. Every day, millions of people are informed by newspapers, radio and television. In Cameroon, the media landscape has changed enormously in recent years: news is now shared instantly on social media, and even a small story can become the most discussed topic in the country within hours.}

\eng{However, this power must be used responsibly. Journalists should check their sources before publishing anything, because false information can destroy lives. Last week, a famous reporter said that the public deserved the truth, and nothing less. Moreover, citizens themselves must learn to distinguish facts from rumours. Although fake news spreads faster than the truth, media education can help young readers become more critical.}

\eng{Therefore, a free and honest press remains the strongest weapon of any democracy. If journalists and readers work together, the future of information in our country will be brighter than ever.}
\end{quote}

\textbf{Vérification de la grille (à souligner sur la copie) :}
\begin{itemize}
\item Voix passive : \eng{millions of people are informed by newspapers...} ; \eng{news is now shared instantly}.
\item Discours rapporté : \eng{a famous reporter said that the public deserved the truth}.
\item Modaux : \eng{must be used}, \eng{should check}, \eng{can destroy}, \eng{can help}.
\item Comparatif / superlatif : \eng{the most discussed topic}, \eng{faster than the truth}, \eng{more critical}, \eng{brighter than ever}.
\item Temps variés : présent simple (\eng{is called}), présent passif, present perfect (\eng{has changed}), prétérit (\eng{said}), futur (\eng{will be}).
\item Connecteurs : \eng{However}, \eng{Moreover}, \eng{Although}, \eng{Therefore}.
\end{itemize}

\textbf{Barème indicatif (sur 20) :} respect de la consigne et du nombre de mots (2 pts) ; richesse et exactitude du vocabulaire des médias (4 pts) ; correction grammaticale et présence des structures exigées (8 pts) ; cohérence, organisation en paragraphes et connecteurs (4 pts) ; orthographe et ponctuation (2 pts).

\textbf{Points de vigilance :} annoncer le titre de l'article ; éviter les phrases trop longues ; vérifier chaque accord sujet--verbe ; ne pas traduire mot à mot depuis le français ; relire en traquant les erreurs classiques (\emph{informations}, \emph{advices}, \emph{the medias}, oubli du \emph{-s} à la troisième personne).
\end{corrige}

\newpage

\coursec{Fiche bilan}

\begin{popfigure}
\begin{tikzpicture}[mindmap, grow cyclic, every node/.style={concept, align=center, text=white, font=\small},
root concept/.append style={concept color=popblue, font=\bfseries\small},
level 1/.append style={level distance=3.4cm, sibling angle=51.4},
level 1 concept/.append style={font=\footnotesize}]
\node[root concept, align=center]{Grammar\\General\\Revision}
child[concept color=poporange]{node[align=center]{Tenses\\present, past,\\future, perfect}}
child[concept color=popgreen]{node[align=center]{Passive\\voice\\be + past\\participle}}
child[concept color=poppurple]{node[align=center]{Reported\\speech\\backshift}}
child[concept color=poppink]{node[align=center]{Modals\\must, can,\\should, may}}
child[concept color=popgold]{node[align=center]{Relatives\\who, which,\\that, whose}}
child[concept color=popblue!70]{node[align=center]{Comparison\\-er / more /\\the most}}
child[concept color=popdark]{node[align=center]{Francophone\\pitfalls\\faux amis,\\word order}};
\end{tikzpicture}
\legende{Carte mentale de la leçon : les sept blocs de grammaire à maîtriser pour le Bac.}
\end{popfigure}

\begin{bilan}
\textbf{Lexique clé du chapitre \eng{Media and Communication}}
\begin{itemize}
\item \eng{the press} : la presse ; \eng{a headline} : un titre (de journal) ; \eng{an article} : un article
\item \eng{a journalist / a reporter} : un journaliste / un reporter ; \eng{to broadcast} : diffuser
\item \eng{social media} : les réseaux sociaux ; \eng{fake news} : les fausses informations
\item \eng{to interview} : interviewer ; \eng{an advertisement (an advert / an ad)} : une publicité
\item \eng{reliable / unreliable} : fiable / peu fiable ; \eng{a source} : une source
\end{itemize}

\textbf{Règles de grammaire à connaître par cœur}
\begin{center}
\begin{tabularx}{\linewidth}{|l|X|X|}
\hline
\rowcolor{popblueL} Point & Règle & Exemple \\
\hline
Present perfect & have/has + participe passé ; bilan lié au présent & \eng{The journalist has just published the article.} \\
\hline
Past simple & verbe + -ed ou forme irrégulière ; action datée, achevée & \eng{She interviewed the minister yesterday.} \\
\hline
Passive voice & be (au temps voulu) + participe passé (+ by) & \eng{The news is broadcast every evening.} \\
\hline
Reported speech & recul des temps ; say \textrightarrow{} said, will \textrightarrow{} would & \eng{He said (that) he was tired.} \\
\hline
Modals & modal + base verbale, sans to, sans -s & \eng{You must check your sources.} \\
\hline
Relatifs & who (personnes), which (choses), that (les deux), whose (possession) & \eng{The reporter who wrote this lives in Douala.} \\
\hline
Comparaison & -er / more ; the -est / the most & \eng{This channel is more reliable than that one.} \\
\hline
\end{tabularx}
\end{center}

\textbf{Méthodes express}
\begin{enumerate}
\item Identifier le temps : repérer la date ou le marqueur (\eng{yesterday}, \eng{since}, \eng{already}, \eng{tomorrow}).
\item Passif : objet du verbe actif \textrightarrow{} sujet ; \eng{be} au même temps que le verbe actif ; participe passé.
\item Discours rapporté : reculer le temps d'un cran, adapter pronoms et marqueurs (\eng{today} \textrightarrow{} \eng{that day}).
\item Toujours relire : accord sujet-verbe, place des adverbes, prépositions.
\end{enumerate}
\end{bilan}

\begin{aretenir}[Checklist avant le Bac]
\begin{itemize}
\item $\square$ Je connais les trois formes des verbes irréguliers courants (\eng{go--went--gone}, \eng{write--wrote--written}).
\item $\square$ Je distingue \eng{past simple} (action datée) et \eng{present perfect} (bilan présent).
\item $\square$ Je sais former la voix passive à tous les temps : \eng{is made}, \eng{was made}, \eng{has been made}, \eng{will be made}.
\item $\square$ Je maîtrise le recul des temps au discours rapporté (\eng{am} \textrightarrow{} \eng{was}, \eng{will} \textrightarrow{} \eng{would}, \eng{have} \textrightarrow{} \eng{had}).
\item $\square$ Je place correctement les modaux : jamais de \eng{to} après, jamais de \eng{-s} à la 3\textsuperscript{e} personne.
\item $\square$ Je choisis le bon pronom relatif : \eng{who} pour les personnes, \eng{which} pour les choses.
\item $\square$ Je forme les comparatifs et superlatifs sans erreur : \eng{better}, \eng{the best}, \eng{more interesting}, \eng{the most interesting}.
\item $\square$ J'évite les faux amis : \eng{actually} = en fait, \eng{to assist} = assister (aider), \eng{a library} = une bibliothèque.
\item $\square$ Je respecte l'ordre des mots anglais : sujet -- verbe -- complément ; l'adjectif avant le nom.
\item $\square$ Je mets la bonne préposition : \eng{to listen to}, \eng{to wait for}, \eng{to depend on}, \eng{interested in}.
\item $\square$ Je relis ma copie : orthographe, accords, ponctuation, majuscules aux nationalités et à \eng{I}.
\item $\square$ Je gère mon temps : lecture du sujet, brouillon, rédaction, relecture.
\end{itemize}
\end{aretenir}

\findecours

\end{document}
